\documentclass[fontset=fandol, a4paper, 12pt, openright]{ctexbook} 

\usepackage{geometry}
\usepackage{hyperref}
\usepackage{fancyhdr}
\usepackage{titlesec}
\usepackage{enumitem}
\usepackage{setspace}
\usepackage{xcolor}

% --- 页面设置 ---
\geometry{left=2.5cm, right=2.5cm, top=3cm, bottom=3cm}
\linespread{1.5} 

% --- 链接与元数据设置 ---
\hypersetup{
    colorlinks=true,
    linkcolor=black, 
    filecolor=magenta,      
    urlcolor=blue,
    pdftitle={圆通的人际关系},
    pdfauthor={曾仕强},
    pdfsubject={人际交往, 国学应用},
    pdflang={zh-CN}
}

% --- 页眉页脚设置 ---
\pagestyle{fancy}
\fancyhf{}
\fancyhead[CE]{\leftmark}
\fancyhead[CO]{\rightmark}
\fancyfoot[C]{\thepage}

% --- 标题格式调整 ---
\ctexset{
    % 设置“篇”的格式
    part = {
        format = \centering\Huge\bfseries,
        name = {第, 篇},
        number = \chinese{part},
    },
    % 设置“章”的格式
    chapter = {
        format = \centering\huge\bfseries,
        number = \chinese{chapter},
        name = {第, 章},
        beforeskip = 10pt,
        afterskip = 30pt
    },
    section = {
        format = \Large\bfseries,
        beforeskip = 20pt,
        afterskip = 15pt
    }
}

\title{\Huge \textbf{圆通的人际关系}\\ \Large “管人是地狱。”}
\author{\textbf{【中】曾仕强}}
\date{}

\begin{document}

\frontmatter

\maketitle

% ==========================================
% 空白页 (扉页后)
% ==========================================
% \newpage                
% \thispagestyle{empty}   
% \null                   
% \newpage        
% ==========================================

\tableofcontents

% ==========================================
% 空白页 (目录后)
% ==========================================
% \newpage                
% \thispagestyle{empty}   
% \null                   
% \newpage                
% ==========================================

% ==========================================
% 序言 (Preface)
% ==========================================


\chapter*{引言}
\addcontentsline{toc}{chapter}{引言}
\markboth{引言}{引言} % 设置页眉
中国人的思想源自伏羲氏，伏羲氏所创的八卦一直都对中国人有深刻的影响。中华文化是在《易经》的影响下形成的，《易经》的主要思想就是人本位，这与西方人所倡导的神本位截然不同。西方式的管理始终离不开神本位的思想，而中国式管理重视人本位，以人为本。

不要以为中国人的言行乱七八糟、模棱两可，如果用《易经》所阐述的道理来考察，就会发现，中国人的所作所为非常有道理。可惜近几百年来，我们经常用西方人的观点和标准来考察自己的行为，才会认为中国人的一切都是乱七八糟的。更有甚者，有人喜欢拿中国人的缺点与西方人的优点去对比，更觉得中国人一无是处。

这种做法大错特错。如果真要比较世界各个民族的优劣，首先要有一个公平的标准，不能拿自己的缺点去跟别人的优点比。这样做，不但有失公平，而且容易导致我们丧失自信心。

中国人有一套自己的东西，这些东西别人学不去，我们也不能去学别人，否则只会自找倒霉。所以，我一直强调管理中国人要靠中国式的管理，而中国式管理中最重要的部分，就是人际关系的管理。

就因为中西方的人际关系大不相同，才导致中西方的管理大相径庭。西方的人际关系建立在平等的基础上；中国人则普遍认为，人与人是不平等的，如果人人讲平等，那是没大没小，而中国人非常讨厌别人没大没小。

现在有个怪现象，很多人喜欢看书，但是看的书越多自己越倒霉，因为他们看的都是西方人写的书：西方人的技术是值得学习的，魏源提出的“师夷长技以制夷”的观点并没有错，科学无国界。可是与文化有关的东西，就要非常小心，否则很容易掉入西方文化的“陷阱”，离中国的实际越来越远。如果这些人从此

生活在西方人的环境中，那问题不大；如果继续生活在中国人的环境中，就会发现自己时时吃亏、处处碰壁。就像清朝末期，很多人留着长辫子却穿西装，到哪里都显得不伦不类。

\section{中西方的人际关系大不相同}

人际关系对任何人来说，都是十分重要的课题。宇宙万物之中，人类的关系最为复杂，而且各地的风土人情不同，人际关系的表现也不一样。

一般而言，西方的人际关系以个人为主。西方人认为，社会由个人构成，个人自由独立，但是必须加以适当的规范，也就是实施法治，才能够维持整体的秩序。人人在法律许可的范围内自由、平等、独立，是西方的人际基础。

中国的人际关系以伦理为主。中国人认为社会固然由个人所构成，但是个人却很难离开社会而生存，个人的自由实际上相当有限。人与人的互动，也不能完全由法律来控制。人人在法律许可的范围内，衡情论理，以伦理来弥补法律的不足，才是我们的人际基础。

正因为中西方的人际关系基础不同，所以二者的人际关系存在很大差别，了解二者的差别，有助于认清我们的人际关系。中西方的人际关系差别如下：

第一，西方的人际关系是神本位的，中国的人际关系是人本位的。西方人认为上帝高高在上，所有人都是上帝的子民。中国人则不同，伏羲氏一开始就让我们认识到，宇宙之间最了不起的不是神，而是人。既然如此，我们为什么非得弄一个神来做自己的主宰呢？

中国人没有自己的宗教，所有的宗教都是外来的。中国人只崇拜自己的祖先，而不去拜什么神——中国人不是在拜神、拜佛，而是礼神、礼佛，就是看到神、佛时，走过去打个招呼，如此而已。

第二，西方人以个人为单位，中国人以家庭为单位。西方人看到一个小孩，通常会直接问他叫什么名字；中国人看到一个小孩，通常会问他是谁家的小孩。中国人会根据小孩的爸爸来判断他的品性，而不是根据他自己。同理，要是小孩做错事，人们通常把错误归到他父母的头上，责怪他们不会教育小孩。

第三，西方人重视平等，中国人重视合理的不平等。西方人可以直呼爸爸的名字，因为大家是平等的。在中国，这样做就是不孝，是忤逆。中国人不认同西方人那种“人生而平等”的观点，而认为人一出生就不平等，而且是合理的不平等。合理的不平等，大家都能接受。但是，过分的不平等，我们就会反抗。中国人不相信绝对的平等，资源有限，机会太少，怎么可能绝对平等？

人与人之间合理的不平等所体现的正是中国人的伦理。爸爸与儿子、上司与下属、老师与学生……永远有高低上下之分，不可能站在同一水平线上。没有哪个中国人敢站在上司面前说：“我和你是平等的，所以你也要听我的话。”中国人对上的态度和对下的态度不同，你对下属敢讲的话，不一定敢对上司讲，这是很正常的。对上级是一个说法，对下级是另一个说法，这就是伦理的体现。

第四，受神本位的影响，西方人只讲权利义务，而中国人讲彼此对待。西方人就算父子之间也是权利义务的关系——儿子18岁以前，靠父母养活；18岁以后，就得靠自己。中国人不看重权利义务，如果中国的父母一等孩子长到18岁就让他自生自灭，会被别人视为狠心的父母。

中国人所重视的彼此对待，就是说：你对我好，我没有理由对你不好；你对我不好，我也不会对你好。中国人会将心比心，投桃报李，这与权利义务没有任何关系。在西方的企业里，上下级之间也只是权利义务关系，你是我的上级，我就会按照规定向你报告。中国人却不这样，你虽然是我的上司，但我不认同你的时候，我就不会向你报告，你要是强迫我报告，我就会敷衍了事。

第五，西方人重视法律，中国人重视道德。西方的法律规定很明确，而且执行得很严格，完全没有人情可言，也没有任何弹性。从理论上讲，法律不应该有弹性，但是没有一点弹性的法律在中国很难执行。自古以来，中华民族不是靠法律约束的民族。古时候，法律的效力是有限的，虽然说“王子犯法，与庶民同罪”，但这只是理想状态，实际上很少这样执行，多数是“刑不上大夫”。在民间故事中，即便是包公这种刚正不阿的典范，也只是“打龙袍”而不是打皇帝。

中国人一般不遵从什么戒律，但是提倡典范，重视道德。“君子爱财，取之有道”，“大位有德者居之”……胡锦涛同志提出的“八荣八耻”同样是道德劝说，并没有立法，强制国民执行。道德是看不见的约束力，而法律是看得见的约束力。一个中国人，如果不讲良心道德的话，是很难在中国生存发展的。中国人不重视有形的东西，凡是有形的东西，对中国人来说，迟早都会变成形式化的东西。我们只靠无形的东西彼此约束，约束别人，也约束自己。

第六，西方人之间充满好奇，中国人之间充满关怀。西方人对人和对动物的态度是一样的，因为西方人认为人就是动物，彼此之间只有好奇。在西方社会，青年男女之间产生好奇就可以同居，一旦失去好奇，双方就会分开。中国人做不到这一点，因为人与人之间不应该好奇，而应该相互关怀。西方人很有礼貌，一见面就会亲切地同你打招呼，但他一点也不关心你。中国人不注重礼貌形式，但是很关心你。看到你嘴巴破了，中国人会直截了当地问：“嘴巴怎么破了，是不是上火了？来，吃点药吧。”

西方人不会看到你的嘴巴破了，看到了也会视而不见，还是会问候你：“你好吗？”而你也清楚，就算抱怨几句自己很不幸也无济于事，因为他并不关心你。此时，就算你疼痛难忍，也只能说：“我很好，谢谢。”

西方人不关心别人，也不希望别人关心自己。中国人以尊老爱幼为美德，而西方人并不如此。你看到西方的老人走路比较吃力，去帮他，他反而不高兴。他的看法是，我有能力处理自己的事，你少操心，你帮我就是看不起我。

第七，西方人之间自然而然地会产生距离，而中国人常常是亲密无间的。西方人重视隐私权，彼此之间保持着戒心，所以会很疏远。而中国人有很强烈的认识对方的欲望，一回生，两回熟，三回见面是朋友，这样慢慢地由不认识到相互了解，再到亲密无间。

人与人之间是要先建立信任才能共事的，你不相信别人，就无法与其合作。西方人之间建立信任靠的是法律的保障，双方在合作前会签订合同，只要稍有不轨，就会受到法律的制裁。中国人之间建立信任靠的是心意的传递，靠的是相互的了解。所以，中国人一见面就会“套近乎”，问东问西，这在西方是绝对不允许的。

第八，西方人之间是利害关系，中国人之间是势利关系。很多人都认为，中国人看重利害关系，其实错了，西方人才重视利害关系。在西方，国与国、组织与组织之间完全是利害关系，没有任何道义可言，个人之间也是如此。而中国人是很势利的。利害与势利有很大的不同，举个例子，公司里有一个员工，表现很不好，公司马上把他辞退，这是利害。员工表现虽然不好，但是他有后台，辞退他会惹来麻烦，所以只能留着他，这是势利。可以说，势利就是复杂的利害。

总而言之，西方的人际关系相对单纯，而中国的人际关系相当复杂，一着不慎，满盘皆输。西方人以“二分法”区分事物，对就是对，错就是错；中国人早已摆脱“二分法”的陷阱，我们知道“错，绝对不可以；对，常常没有用”。对错之外，还牵涉是否圆通的问题。我们厌恶是非不分的人，也不欢迎是非分明却不圆通的人，“水至清则无鱼，人至察则无徒”就是这个道理。中国人讲求“在圆通中分是非”，把是非分得大家都有面子，不得罪人，但也不讨好人，人际关系才可能良好。

中国人喜欢拉关系、靠关系。这句话很容易引起人们的误解，朝坏的、不正当的、不合法的地方想。有人认为某些人的成功，是讨好别人的结果，而自己的成功是凭本事获得的，甚至公开宣称：“我从来不搞关系，我现在的一切，完全是凭真本事得来的。”事实上，如果一个人毫无能力，是无法完全依靠人际关系而成功的。但是，即便有高超的能力，如果缺乏良好的人际关系，也不可能成功。在家靠父母，出外靠朋友，说的就是人际关系的重要性。

当然，我们也看到某些人用心营造不正常、不正当的关系，然后用来营私舞弊，祸国殃民。这种人际关系所带来的弊端的确使得许多人为之心寒，以致他们认为人际关系只有害而没有利，因此不重视也不研究人际关系。还有一些人，由于自己不擅长建立人际关系，眼见他人因人际关系而获利，出于嫉妒或不满的心理，对人际关系产生强烈的反感。

于是，很多人把“搞关系”看成负面的东西，似乎好人从不搞关系，只有心术不正的人才搞关系。这样看待人际关系，自然形成偏激的态度。人际关系本身是中性的，运用得恰当，便是良好的人际关系；用错了，当然产生不好的影响。行为正当的人对拉关系、套关系，实在不必过分敏感。往好的方面想，反而容易获得良好的效果，何乐而不为?

\section{人际关系需要伦理道德规范}

中国人的伦理观念比其他民族发展得都早，而且最完善。孟子说：“使契为司徒，教以人伦：父子有亲，君臣有义，夫妇有别，长幼有序，朋友有信。”就是说，父子之间有骨肉之亲，君臣之间有礼义之道，夫妻之间挚爱而又内外有别，老少之间有尊卑之序，朋友之间有诚信之德，这是处理人际关系的行为准则。

自古以来，我们制定了形形色色的准则，无非是为了加强对个人的约束，提醒我们除了自己以外，还有各种有关系的人，因而自己的一言一行都要格外谨慎。其中，君臣、父子、夫妇、兄弟、朋友这五伦，是人生不可或缺的，对中国人的言行产生重要的影响。

既然伦理是处理人际关系的准则，那么我们的人际关系势必打上伦理的烙印。自古以来，我们所建立的是一种罕见的人伦关系，中国人根本没有什么人际关系。而我们一直错误地想把西方的人际关系移植过来，弄得大家的关系愈来愈紧张，愈来愈败坏。

人伦关系和人际关系最主要的差异在于，前者重视“合理的不平等”，而不是后者所主张的“平等”。西方人认为“人生而平等”，于是发展出一套平等的人际关系。中国人认为“人一出世，就不平等”，就算同一家庭、同一父母所生的子女，在资质方面也不相同，加上出生时家庭的环境、父母的年龄与社会地位也不一定一样，子女们怎么可能平等呢?先天不平等，后天也不可能平等，?多经过合理的调整，做到合理的不平等。

伦理就是合理的不平等，父父、子子，应该各如其分；君君、臣臣，用现代的话说是上司、下属各自扮演不同的角色。彼此之间，必须维持“合理不平等”的分寸，而不是“平起平坐”的“不合理的平等”。二者之间的差异，必须慎重拿捏，才不致发生差错，避免自己不明不白地遭受“平等”的祸害。只要把“合理的不平等”这个观念端正起来，很多问题都会解决。我们现在满脑子追求平等，弄得大家很不幸福。人有胖瘦之分，两个人坐在一起，胖的人占的位置自然就大一点，瘦的人就委屈一点。两个人吃一样东西，喜欢吃的人就多吃一点，不喜欢吃的人就少吃一点，为什么非要一人一半呢？全世界都是相对的平等，没有绝对平等，只是表现的方法不同而已，没有对错，没有好坏。

中国人十分重视做人、做事的道理，人际关系如何反映出做人的效果。只会做事而不会做人，人际关系搞

不好，得罪了许多人，又怎么能够好好做事呢?透过好好做人，来好好做事，通常比较有效。

人际关系既是做人的道理，也可以说是做人的技巧。做人讲求技巧，免不了有一些权谋、圆滑、奸诈的味道，引起很多人的反感。这时候注入伦理道德，可以使权谋变成权宜应变、因时制宜，圆滑变成圆通，而奸诈也变成一种机警。

若是只学做人的技巧，而忽视做人的原则，不但没有成效，而且会被人嘲笑。做人不可以玩弄权谋，许多人误把圆通、应变看成讲求谋略，其实应该是策略才对。做人做事可以有策略，不可讲谋略。换句话说，一切要求应当正当合理，不应该有不正当的念头。

我们有三个根深蒂固的观念，很不容易改变：

第一，法是死的，人是活的。死的法需要活的人来加以合理运用，而不是不动脑筋地死守法律规章。

第二，天下事好像没有什么是不能变通的，若是变通不了，大多是因为找错了人。只要找对人，变通应该是没有问题的。

第三，法由少数人订立，由一个人修改。这种现象是中国社会自古迄今的一贯精神。一些可以控制别人的人，碰到别人稍有反对意见就加以恐吓、威胁，甚至杀一儆百。

对中国人而言，法不够用。因为，中国人不喜欢违法，不做违反规定的事情。但是中国人普遍喜欢动脑筋，做一些法律没有规定的事情。这种情况必须用伦理道德来弥补，才能收到预期的效果。只有大家凭良心有所不为，才有可能加以改变。夜不闭户不是靠法律来实现的，而是靠道德达成的。

总的来说，人伦关系的重点，在“公正”而“不平等”。对上要有礼貌，但不可以谄媚、讨好；对下不宜太严，也不能过分宽松、纵容；平行同事，不必太拘束，也不可以过分熟不拘礼。其中的轻重，必须因人、因时、因地、因事适当拿捏，这样，才称得上公正。只有用心体验，不断改善，才能达成良好的人伦关系。

为了叙述的方便，本书仍使用人际关系一词。

\mainmatter

\part{圆通的四个维度}

\chapter{言——模棱两可，言不由衷}
《孙子兵法》说，知己知彼，百战不殆。彼此都是中国人，建立人际关系时，就更需要确确实实了解对方在说些什么、想些什么。

俗话说，言为心声。中国人的心是难以捉摸的，所以中国人的话常常模棱两可。我们主张“逢人只说三分话”，同样主张“知无不言，言无不尽”。“逢人只说三分话”是对交情不深、关系不够的人而言的，因为人心隔肚皮，知人知面不知心，当然应该小心试探。“知无不言，言无不尽”是对交情深厚、关系密切的人而言的，既然大家亲如一家，也就不需要互相隐瞒。

其实，中国人说“逢人只说三分话”的时候，已经含有“知无不言，言无不尽”的意思。况且中国人所说的“三分”，既可以是“三分流水七分尘”的“三分”，也可以是“天下只有三分月色”的“三分”，就看到时如何拿捏了。当彼此尚不熟悉时，当然“未可全抛一片心”，等到互相信赖了，完全可以知无不言，言无不尽。

同样，中国人说“言无不尽”的时候，也不要忘记“逢人只说三分话”，因为彼此虽然关系密切，但是有的话可能会伤害对方的自尊心，或者引起他的嫉妒，所以必须有适当的保留，说三分留七分，那七分就心照不宣了。

\section{“言”的维度}

不了解中国人的人觉得我们很难捉摸：“我明明听懂了他的话，他怎么还是不高兴？”因为中国人说的话通常包含很多意思，听懂了表面意思却常常听不懂言外之意。有时候，中国人不说话，只是一个眼神、一个动作就包含了很多意思，这当然需要彼此的默契。如果没有默契，要搞清楚中国人到底在说什么，确实很难。下面这个小故事就说明了这一问题。

苏东坡被贬至黄州后，一天傍晚，和好友佛印和尚泛舟长江。忽然，苏东坡用手往岸上一指，笑而不语。佛印顺势望去，只见岸边有一只黄狗正在啃骨头，顿有所悟，便将自己手中题有苏东坡诗句的蒲扇抛入水中。两人心领神会，不禁相视而笑。

原来，这是一副哑联。苏东坡的上联是“狗啃河上（和尚）骨”；佛印的下联是“水流东坡尸（东坡诗）”。

当然，要达到苏东坡与佛印这种心灵相通的程度，除了默契之外，还要有较高的智商才行。

中国人南腔北调，就算大家都说普通话，也很难保证彼此听得懂。勉强听懂，也未必弄得清楚他的真正用意。

指鹿为马的故事大家都熟悉，人们把“指鹿为马”理解为“不明是非、颠倒黑白”，其实很不恰当。赵高只是借此试探一下究竟自己在朝廷的权势怎样，其真正的话意是：“你们是服从秦二世，还是服从我?”我们嘲笑群臣不分黑白，其实他们才是真正听得懂话意的人。

有一天，我搭出租车到某地。由于大路发生事故，所以司机改走小路。但小路蜿蜒曲折，司机不太熟，越走越觉得没有把握，便停下来，问路旁边一位老先生：“请问我要到某地去，该怎么走?”

老先生气定神闲，不慌不忙地回答：“有路就可以走，多问几次就会到。”

这两句话，叫人听了觉得十分有道理，同时又觉得摸不着头脑。

司机表示感谢，很有信心地向前驶去。

我觉得很纳闷，问他：“你知道怎么走了?”

他说：“知道。有路就可以走，表示我走的路是对的。如果我已经走错了，他会把手一扬，然后指向正确的方向。现在我走对了，他不必举手，所以说有路就可以走，告诉我顺着这条路一直走下去。多问几次就会到，意思是后面会有几个比较复杂的岔路口，那时候一定要问路，不要乱闯。”

经他这么一解释，我才恍然大悟。原来中国话如此简单明了，两句话就可以交代清楚。但听者必须动脑筋才能听懂。现在有些人只听不想，以致听不懂中国人所说的话，实在是一种遗憾。

中国话很不容易听，才是我们真正的难处。中国人要人家不要“听”话，其中含有“中国话不可以用耳朵听”的意思，必须特别小心。“不要听他的”，包含“不要听他的话”，也包含“不要单凭耳朵听他的话”的意思。中国人很少说：“听他说什么。”反而常常告诫我们：“看他怎么说。”也就是说，中国话不适合单用耳朵听，应该配合眼睛看。中国话听起来含含糊糊，“看”起来清清楚楚。中国人不喜欢啰里啰唆讲一大堆，只喜欢简单明了，短短一两句话，含意很深，所以“看”了之后，还要多想。如果不用心想，还是弄不清楚中国人的话意。也就是说，“看”话不能单凭一双眼睛去看，还要动用“心眼”，才能够真正看清楚，才能领悟“话中的话”，以及“话外的话”。

“心眼”要大，才听得真实。“心眼”太小，成了“小心眼儿”，就会“以小人之心，度君子之腹”。如果对方有难言之隐，有说不出来的苦衷，有说出来反而彼此难过的事情，千万不要用不正当的心思去曲解。

一句“你看着办吧”，究竟是“全权委托你”，还是“猜猜我的用意”，甚至“居然搞成这样子，你自己收拾烂摊子吧”？短短五个字，足够让别人思前想后了。凡是耳朵听不懂的时候，就要用眼睛看，还要动脑筋，结果呢?你看着办吧!

\section{“随便”并不随便}

中国人常把“随便”挂在嘴边上。请客吃饭时，总是说：“没有什么好菜，随便吃点。”实际上菜肴十分丰盛。反过来说，客人真的随便吃喝，丝毫都不客气，即使宾主交情甚深，主人看在眼里，也会不太舒服：又不是只有你我两人，当着我的妻儿，你也未免太随便了!

“随便”的意思有两种：一是“随意、任意”，一是“不拘束、不认真”。但是，在日常生活当中，使用起来却有三种含义：

1.“随便”代表“看看你的诚意”。人家问我要什么?我说“随便”，意思就是说：衡量你自己的能力，可以提供什么。中国人认为，我说“随便”是不想为难你，如果你真的“随便”，就是轻视我。如果我到你家做客，你问我喝什么，我说“咖啡”，你没有，岂不是很尴尬，所以我只能说“随便”，给你留了面子。你虽没有咖啡，但是给我端上一杯好茶，我自然也很高兴，因为你没有“随便”应付我。但是，你真要“随便”给我一杯白开水，那我肯定不高兴：“如果只有白开水，何必问我喝什么?”

你请我吃饭，问到哪里用餐?我当然不能直截了当地提议上豪华酒楼，万一你认为那样不值得，我岂非自讨没趣?不过我也不愿意自贬身价，一开口就选择普通餐厅，非但显得土气，对方也未必领情。最好的办法，还是说“随便”。至少可以了解你认为拿什么招待我最合适，进而了解我在你心目中的地位，以便调整自己所应表现的态度。你经济拮据，只能够请我上普通馆子，我照样吃得高兴，因为你够诚意，我不在乎吃什么；如果你手头宽裕，豪华酒楼也请得起，却只请我上普通的馆子，我就知道你不够意思。

“随便”绝对不是含糊，而是“在和谐中找到合理”的一个代名词。中国人如果真的随随便便，一定没有前途。

甲和乙是好朋友，一天，乙到甲家里做客，甲热情招呼，顺口问他：“喝点什么东西？”乙回答：“随便，随便。”

甲当然心里有数：家里确实有好酒，是留给上司丙的，现在当然不能拿给乙喝。衡量与乙的关系，决定泡一壶好茶招待他。

乙见甲并没有敷衍他，自然很高兴，也知道自己在甲心中的分量。

正在此时，丙不期而至，明显就是来喝那瓶好酒的。此时，甲该不该将好酒拿出来？如果拿出来就会得罪乙，不拿出来又会得罪丙。这种两难的情况却也难不倒深谙圆通之道的甲。

甲对他太太大声说：“我刚才找了半天，你到底把那瓶好酒藏到哪里去了？”

甲的太太明白甲的意思，也大声地回答：“我昨天收拾屋子，怕把它弄脏，特地藏起来了。”

话音未落，甲的太太就拿着好酒出来了，并准备了丰盛的下酒菜。乙和丙都很高兴，当然，甲将危机巧妙地化解于无形，更是高兴。

用高声说话来向别人传达自己的想法，是中国人的绝招，在人际关系的运作中，具有转危为安的决定性力量。甲很清楚，丙提前来就是想喝自己的那瓶好酒，可是乙来了自己没拿出来，等丙来了才拿出，乙心里肯定不高兴。可是不拿出来，就会引起丙的不满。拿也不是，不拿也不是，但是高声说话就解决了问题：“刚才找了半天”就是告诉乙“我本来想请你喝，但是没找到”；而丙明白，就算乙是甲的好朋友，还是没我的面子大，哪怕是“特地藏起来了”，到了紧要关头，也得拿出来。

知识很重要，但是知识之外，人际技巧也很重要，若想圆满解决问题，在知识之外，还需要一些艺术气氛。

2.“随便”表示具有“单凭物质不足以表达全部的敬意，必须拿精神来补助”的用意。以请客吃饭为例，主人即使准备很多佳肴，仍然说“随便吃点”，意思是这些菜虽然很好，但总觉得应该有更好的，才足以表达主人对贵宾的敬意。如果客人认为过于丰盛，就会说“太破费了”；如果客人觉得不过尔尔，既然主人说了“随便吃点”，那客人也不会不满。

如果东西非常好，中国人也会轻描淡写地说“随便买的”，这不是谦虚，而是中国人认为精神重于物质，“千里送鹅毛，礼轻情义重”。同时，也希望对方不会有“受之有愧”的负担。如果东西并不好，一句“随便”，表示“我已经尽力，希望你能够谅解”。中国人讲究“尽心尽力”，只要尽力而为，对方多半是会体谅的。

3.“随便”暗示“我有我的意见，只是不便说出来”。上司征求下属的意见，下属绝不敢说“随便”，而是请上司做主。因为说“随便”，就意味着下属有自己的意见，只是不方便说而已。

“随便”正是孔子的“无可无不可”的体现。人与人之间的关系是很微妙的，时时有变化，如果丝毫不加考虑，冒冒失失地说出自己的意见，很可能使对方为难，不如将心比心，先说“随便”，好让彼此有个商量的余地。

既然“无可无不可”，“随便”就不是对“任何人、任何事、任何地”都通用的。有的时候，随便说“随便”也会给对方造成困扰。例如，晚辈对长辈最好别说“随便”，否则就显得太随便。朋友之间也有不宜说“随便”的情况，同是请客吃饭，说“随便”可以让主人根据自己的情况准备，但是如果主人境况不好，要保全主人的“面子”，就不宜说“随便”，最好说：“我最近肠胃不好，很怕油腻，就在附近这家饭馆吃点素菜就行。”

有人托你帮忙，要请你吃饭，你就要暗自盘算一下，“吃人家的嘴软”，背负人情太重，哪里敢说“随便”，赶忙借故推辞，心领为上。

“随便”并不是“差不多”，而是以“合理就好”为原则，应该说“随便”的时候，才可以说，不应该说的时候，就不可以说。别人在对你说“随便”的时候，如果他不是随便说的，那就等于说：“你自己想想，怎样才合理；只要合理，我当然就随你的便!”

“随便”绝不是“马马虎虎”，我们将心比心，既不希望人家“马马虎虎”待我，当然也就不可以“马马虎虎”对待别人。该“随便”才能随便，不该“随便”绝对不可以随便。

以“不随便”的态度来“随便”，才能符合合理的标准。我们一方面不可以随便说“随便”，另一方面也不能够随便理会或处置别人的“随便”，因为一个有修养、有分寸的中国人，是不随便说“随便”的。

\section{“不随便”并非“随便”}

中国人的基本立场是不偏不倚，秉持既不反对也不赞成的态度，在征求中国人意见的时候，中国人基本上说的都是既不赞成也不反对的话。因为事情总是不尽如人意，所以只能赞成一部分而反对另一部分。任何事情，总是不断地演变：原本可以赞成的事情，演变到最后，令人不得不反对；而原来应该反对的，也可能愈变愈应该赞成。这些都是中国人不愿意公开说赞成还是反对的原因。

当把一切弄清楚，而且保证不再改变的时候，中国人自会当机立断，给你一个准确的答复。但是，就算如此，警觉性特别高的中国人，也尽量避免公开表态。一旦你明确表示赞成，那些不希望你赞成的人就会给你施加压力，相反，你反对的话，那些不希望你反对的人也会频频阻挠，徒增许多麻烦。所以，最后养成中国人凡事都含含糊糊的，“这件事嘛，呵呵呵……”。

中国人之所以这样，是因为没有安全感，赞成或反对就像下赌注一样，万一押错宝，后果惨重。只要保证他的安全，要中国人表明赞成或反对的立场，其实并不难。古时候，大臣要说一些可能冒犯皇帝的话时，总会先请皇帝饶他不死。这就是在寻求一种安的状态。

如果有人直接问你这种问题：“你赞成领导关于这件事的处理吗?”那他可能是不怀好意，挖了陷阱让你跳。你说“赞成”或“反对”，都会对自己不利，不是被利用，便是被嘲笑。

碰到这种情况，中国人不可能直截了当地表明态度，而是会反问：“你认为如何?”然后兼顾赞成及反对双方面的看法，适当地表述自己的意见。而被反问的人也是同一态度，或者采用“我也不反对”或“我也不赞成”来应付，因为中国人明白，“赞成”、“反对”不是“二分法”，“不赞成”并不代表“反对”，“不反对”也并不代表“赞成”。

也许有人认为中国人不实在，见风使舵，虚情假意，而且缺乏胆识，不敢担当。其实，我们说既不赞成也不反对，事实上是赞成之中有反对，而反对之中也有赞成，并不是不分是非，糊里糊涂或者怕惹事端。先赞成后反对，先反对后赞成，对方比较容易接受，这种“攻心”术是中国人的独特手法。例如：“我毕业后马上出国留学，你赞成吗？”如果你不赞成，最好回答：“我不赞成你马上出国，但是如果你准备得十分周全，知道自己所要学的是什么，将来学成之后要做什么，我当然不会反对。”

既不说赞成的话也不说反对的话，一方面保证自己的安全，另一方面可以增强对方的责任，凸显听者的自主性，使其更加重视自律和自动。一旦表示赞成，听者受到很大的鼓励和支持，可能大意失荆州，造成阴沟里翻船的惨剧。而一旦表示反对，听者受到挫折，可能因而放弃，或者缺乏信心。不赞成也不反对，听者才会面对现实，用心地研究判断，自己做最后的决定。

中国人相信，公开表示赞成或反对，事实上都不一定可靠，不如采取观察、试探、测试、迂回打听等方式来加以判断。若是因为怕得罪人，吞吞吐吐，既不敢赞成也不敢反对，那就是心术不正，终将被别人厌恶，自己也不会有什么好前程。合理的赞成加上合理的反对，才是正当的行为。

当然了，在人际交往中，圆通的人也不应该直接问别人是否赞成或反对，以免给人别有用心之嫌，一般只问“你有什么看法”之类的语气较弱的问题。

看清楚之后，合理地表达反对或赞成的意见；尚未看清之前，千万不要冒冒失失地说出“我赞成”或“我反对”，以免被利用，或被人看不起，令自己骑虎难下。


\chapter{行——谨慎小心，反求诸己}
“合理就好”的理念使得中国人的一言一行都要把握好分寸。谨慎的人在人际交往中，就像林黛玉进贾府一样，不肯轻易多说一句话，多行一步路。

\section{“谨慎小心”并非“随便”}

中国人在人际交往中，首先要弄清楚对方到底是谁。中国人认为“有人才有事”，而且“事在人为”，很难“对事不对人”，所以常常把人和事联系在一起。在中国社会，每听到一句话，如果不清楚是谁说的，就很难判断它究竟是对的还是错的、是真的还是假的。可见弄清楚对方是谁，乃是开展人际交往的第一步。

中国人比较倾向于“差别性待遇”，以不同的标准来对待不同身份的人。如果你的职位比我高，那你说什么都应该是比较正确的。如果你和我职位平等，那就以“来而不往非礼也”的态度对待你。如果你的职位比我低，我不会以大欺小，却绝不容许你以“下”犯“上”。

这样看来，中国人似乎没有是非观念，但事实绝非如此。在中国社会，职位低的人是不宜反驳职位高的人的。举例说明，如果你的上司冤枉了你，你该怎么办？据理力争的话，他会明白是他错而非你错，但这又如何呢？他身为上司，竟然失察而冤枉了你，自然觉得相当没面子。

中国人没有面子的时候，最要紧的，便是设法找回面子。如何挽回他的颜面？很简单，一心一意地找你的毛病，只要被他抓着了，他的面子便全回来了。他一心一意找你的差错，你真的插翅难逃，因为“人非圣贤，孰能无过”，你迟早会被他逮个正着。

但是，如果上司冤枉你后，你保持沉默，一句话都不讲，表面上看你是忍辱负重，其实不然。你的上司看到你居然一言不发，就会觉得奇怪：“这个人怎么搞的?难道我屈说了他？”于是，他自然想着解开谜团，结果发现自己确实冤枉了你，并由于自己内心愧疚而善待你。

一般来说，上司冤枉你，纯属偶然，很少有上司故意颠倒是非，存心而为之。但是你若不幸碰上这种上司，足以证明上司早已容不下你，在这种情况下，你据理力争又有何用？不如另谋高就，“此处不留人，自有留人处”。如果你没有别的门路，最好忍气吞声，说不定上司见你一直逆来顺受，网开一面，不再针对你。

如果你不幸遇到的是个迷糊的上司，这种情形倒是常见，因为无能者却身居高位并不是稀罕事。遇到这种上司，你说得再对，他都可能斥之为狡辩，于事无补，不如安静下来，好好做事，以免被他抓住把柄。

迷糊的上司不好，但过于是非分明的上司也很难相处。因为太过分明，以致刚愎自用的人，总是认为自己的看法都是对的。说你错，你就错，你再说多少话他都不会更改原有的判断。你据理力争的话，只会弄得面红耳赤。

如果你的上司不是以上三种人，那他就是无心犯错。这种无心的过失，是应该谅解的。没有必要得理不饶人，让无心错怪你的上司难堪，所以应当用沉默来暗示他有错误，使上司自己察觉，自行校正。

如果有人和你职位相当，却指出你的错误，而且到处宣扬，那你就会积极找他的过失。这不是面子问题，也不是心胸狭窄的表现。如果他发现你有差错，当面规劝的话，虽然你一时可能无法接受，但是他既然说的是事实，出发点又是为你好，你终究会心生感激，又怎么会记恨他呢？但他却没有和你说，反而到处宣扬，分明是让你难堪。中国人一向是交互主义，你对我好，我也对你好，相反，你不仁就别怪我不义。所以，你也只有全力找他的差错，照样宣扬一番，让他尝尝同样的滋味。

如果有人职位比你低，却敢在背后指摘你的过失，你通常不会有任何顾虑地去整一下；而且“当年别人教诲我，如今我也应该教诲别人”的想法，很容易变成理直气壮的借口，整他，是让他明白做人的道理：有话最好当面说，不要背后胡扯。

所以，对方不同，你采取的措施也应该不同。先问清楚是谁说的，再做定夺，这就是一种“经”，如何应变，则是个别的“权”。

中国人既然以人为主，那一切事都离不开人，也就是离不开人际关系。何况我们一直重视伦理，对于人的身份地位十分关心。看见或听说一个人，总要进一步追问是什么样的人，并且依据身份地位做出不一样的反应，才算合理。这样才能建立良好的关系，以便进一步达成预期的企图。这不是势利眼，只要保持合理的程度，不要前倨后恭，没有什么不好。

相传，清代大书法家郑板桥去一个寺院游玩，并去拜见方丈。方丈见他衣着俭朴，以为他是一般俗客，就冷淡地说了句“坐”，又对小和尚喊“茶”。一经交谈，方丈顿感此人谈吐非凡，就将郑板桥引进厢房，一面说“请坐”，一面吩咐小和尚“敬茶”。又经细谈，得知来人是赫赫有名的扬州八怪之一的郑板桥时，方丈急忙将其请到雅洁清静的方丈室，连声说“请上坐”，并吩咐小和尚“敬香茶。”最后，这个方丈再三恳求郑板桥题词留念，郑板桥思忖了一下，挥笔写了一副对联。上联是“坐，请坐，请上坐”；下联是“茶，敬茶，敬香茶”。方丈一看，羞愧满面，连连向郑板桥施礼，以示歉意。

这个故事通常用来讽刺方丈势利眼，其实方丈是根据对方的身份，采取了相应的方式和态度，只不过一开始就没弄清楚、看走眼罢了。

\section{“谨慎小心”并非“随便”}

西方人在同别人打交道时，不论干什么，都要通过法律途径约束他人来防止自己上当，而且对于欺骗别人的人是相当鄙视的。中国人不会要求别人怎样，只会提醒自己“防人之心不可无”。我们很少去笑那些骗子，反而多少有些羡慕：“他真高明，居然骗了这么多人！”而对被骗的人，在同情之外，还有几分不屑：“你怎么这么笨，这种当也会上？”好像责任不在骗子，而在被骗的人。

其实，嘲笑被骗的人与鄙视骗子的目的完全相同，都是遏止骗人事件的发生。

西方人鄙视骗子，使其备受压力，抬不起头来。而中国人认为人人应各自小心，不要上当，使那些骗子无从下手。我们用“嘲笑被骗的人”的形式，来提醒大家，千万不要被骗，否则不但损失惨重，还要惹人嘲笑。所以在中国社会，那些被骗的人，反而不敢以正面示人。

西方人的想法很对，只要大家齐心，共同指责骗子，便可以减少骗人的行为。而中国人更实际一些，要求自己岂不比要求别人更容易一些？“求人不如求己”，要求别人不要欺骗，远不如自己提高警觉更可靠。防人之心其实就是让人小心不要上当，因此“反求诸己”更胜一筹。

试想，如果你和别人初次见面，就老实地对别人说：“我是个老实人，人家说什么我都信，所以请不要骗我。”那对方会怎么想？当然是：“他最容易上当，不骗他骗谁?”如果蛮横地要求对方：“你不要骗我，否则我要你好看！”对方就想：“骗你又怎样？你既然被我骗了，说明你不如我，你又有什么能力要我好看？”结果照骗不误。

一般中国人是不会这么做的，如果有人真的这么做，他往往会有拆穿骗局的能力。比如，老板对干部说：“我不会亏待你，但你千万不要骗我。你若是骗我的话，我绝不轻饶。”笨的干部就以为老板是个好骗的人，一开始虽然没想欺骗他，久而久之，就忍不住做些欺上瞒下的事。结果很快就被老板发现了：“我警告过你，千万不要骗我，可是你太让我失望了，我只好开除你。”

如果认为中国人专门欺负弱者，谁吃亏就笑谁，谁倒霉就笑谁，没有同情心，那就错了。中国人只有对自己人，对熟悉的人，或者有利害关系的人，才会“痛心”地嘲笑他，目的在于“加深他的印象”，使他深切体会“人家在你面前同情你，实际上背后都在笑你”，因而决心“自己小心不要上当”。对于陌生人，我们根本无从笑起，因为彼此没有“关系”，产生不出任何联想。对于认识而交情不够深厚的人，我们只会背后嘲笑他，但当面尽量不提，万一对方自己说出来，我们会表示支持，痛骂欺骗他的人。

但是“小心不要上当”，并不是“不要相信别人”，而是不给坏人可乘之机。

正因为小心，才会时时用心。可以相信的时候，“疑人不用，用人不疑”；不可以相信的时候，“知人知面不知心”，“人心善变，不可不防”。这些话都是中国人常说的，看起来互相矛盾，却是因“时”而制宜。

在人际交往中，你可以选择相信别人，也可以选择不相信别人。你相信别人，万一被骗，大家就会嘲笑你，“三两句话，把他骗得团团转”，结论是你“缺乏判断力”。你不相信别人，又怎能与其相处？在信与不信之间，需要自己拿捏。

在正常的情况下，当然应该相信别人，但是只能相信到合理的地步；遇到不合理的地方，就不应该相信他。站在不相信的立场来相信，才不至于一相信就上当。对任何人都相信，受骗的几率就会大幅度增加。你无法要求别人绝对诚实，因为没人能做得到，所以只能自己小心。社会上骗子并不多，但是喜欢被骗的人太多了，才会发生那么多骗人的事件。骗人的人当然不对，而被骗的人也未尝全对，谁叫他不小心?“防人之心不可无”，这句话永远不会过时。

人类是群居的动物，彼此互相合作才能生存。人应该相信别人，这是天经地义的事情。但是过分相信别人，同样引起别人欺骗的兴趣，以致上当，也是不争的事实。吃亏上当，人人不喜欢，却又经常发生这种不愉快的事情，主要是人们喜欢占小便宜，才造成因小失大。提防上当，最有效的办法，便是切记不要贪小便宜。

\section{凡事都追求合理}

中国人的观念是凡事追求合理。换句话说，把事情做到合理的地步，人们才会接受。我们虽然重视典章制度，却明白典章制度容易僵化而不合时宜，因此在典章制度的范围内，多半中国人喜欢权宜应变，以求合理变通。凡是不会变通的人，都会被人斥之为“死心眼”、“死脑筋”。

中国人认为，既然人与人之间是相互影响的，在人际交往的过程中，就应该“反求诸己”，即“要求他人合理之前，先求自己合理”，以自己的合理来影响他人，使他人亦能合理。

自己不合理，却希望他人以合理待我，结果经常不如意，这时怨天尤人也是枉然。自己先求合理，再来期望他人以合理待我，才是合理的态度，也比较容易实现。但是合理不合理，各人的标准未必一致，有时候你认为合理，我就认为不合理，因此这个“合理”也会引起很多争执，产生很多不愉快！

在中国社会，凡事自己求合理，就容易获得别人的好感。即使你的“合理”和他的“合理”会有些误差，但你的出发点是好的，起码会问心无愧。至于是否会引起争执，听天由命吧。虽然这听起来比较消极，但总比表现出不合理的行为，使大家心里不高兴，甚至造成误解的好。

自己合理不合理要通过他人的反应来判断。他人的反应就像一面镜子，会如实地反映你是否合理。当对方表现出不合理的行为时，不要指责对方的缺失，而是应该反省自己是不是不合理。如果是的话，赶快先把自己的缺失调整过来，对方也可能跟着表现出合理的反应。要改变对方，最有效的方法是先改变我们自己。

举个例子，中国人认为，朋友的朋友是朋友，敌人的敌人也是朋友，所以三个人之间的关系很复杂。正因为如此，一个人在不清楚对方的意图的时候，是不会承认自己和另外一个人的关系的。在这种情况下，贸然问起别人之间的关系，很难得到正确的答案，这就是自己的行为不合理造成的。

乙和丙是熟人。一天，甲和乙聊天时，甲突然提到丙，问：“你认识丙吗？”乙的第一反应是“不认识”，因为他不知道甲有何意图。

一句“不认识”，虽然只有三个字，却包含好几种不同的意思：“真的不认识”、“虽然认识，但是并无交情”、“认识是认识，跟不认识差不多”、“你有什么事情，要问我认不认识”以及“你少打我的主意，我认识不认识根本与你无关”。

说“不认识”有很多好处：第一，减少许多风险，省却许多口舌，避免许多麻烦。如果乙在不清楚甲的意图时就说认识，甲接下来很可能说：“那真是太好了，我正好有一件事要找丙帮忙，麻烦您帮忙引见一下好吗？”这不是自讨苦吃吗？这时候说“好”，会增加不少麻烦；说“不好”，等于驳了甲的颜面，令甲不满：“这点小忙都不愿意帮，还说是朋友呢？”

当丙身居高位时，乙如果先说“认识”，事情来了，再推说“不认识”或者坦白说明自己不愿意帮忙，甚至直接指称对方根本没有权利提出要求，都是伤感情的做法。在中国社会，得罪一个人的后果是很严重的，尤其不能得罪小人。

第二，可以探听信息。如果丙对乙心存不满，忍不住对甲抱怨了几句，甲如果不清楚乙和丙的关系，多半不敢直截了当地传话过去，必然先问：“你认识丙吗？”假若乙说“认识”，甲就不会再说什么，以免有搬弄是非之嫌，而乙也徒然失去获得信息的机会。若乙说“不认识”，甲才会放心地把丙的话告诉乙，乙也能了解丙对他的不满，以便采取适当的对策。

明明认识却装成不认识，这种不合理的行为完全是因为甲不合理在先。打听别人之间的关系，一开始就应该把目的说清楚。比如：“你认识丙吧？他托我给你带了点礼物。”乙一定回答“认识”。如果不能一开始就做到合理，直接问：“你认识丙吗？”得到“不认识”的回答后，就要调整：“哦，他说和你很熟，还想约你吃个饭呢。”那乙也会向合理的方向调整。不要以为从“不认识”调整到“认识”很难，中国人很会“颠倒是非”，拥有足够的智慧，可以进退自如，因为凡事早已留有余地。

乙只需作恍然大悟状：“哦，你说的是他啊，抱歉抱歉，刚才没听清，还以为你在说谁呢。”轻轻松松就兜回来了。

先说明自己的目的，再去问对方，以便给对方斟酌的空间，这才算合理的行为。不要责怪别人不诚实，不要认为一切不合理，都是“我”以外的中国人造成的。殊不知种种缺失，实际上都与“我”密切相关。

可惜，很多人不明白这一点，不反省自己是不是诚实，是不是考虑周到，是不是明白透彻，便大胆地指责中国人不诚实，实在不公平。


\chapter{心——追求圆通，善于自保}
中国人常将人分成两种——君子和小人。孔子对君子与小人有过许多论述，比如，“君子怀德，小人怀土；君子怀刑，小人怀惠”、“君子喻于义，小人喻于利”、“君子和而不同，小人同而不和”，等等。人们历来崇尚君子，但事实上君子斗不过小人的情况比比皆是。因为小人非常重视人际技巧，很容易打动别人的心，进而抓住别人的心。君子往往觉得自己凭良心，不做亏心事，不必讲求什么人际技巧，因而时常得罪人，引起别人的怨恨和报复。这一点是君子不及小人的地方。那君子为什么不改变一下，同样注重人际技巧，增强自己的竞争力呢?

\section{圆通而不圆滑}

人们常常认为小人圆滑，因而对他充满鄙视。君子当然不应该圆滑，但是要足够圆通，否则就常常使自己在与小人的斗争中处于劣势。中国人最讨厌圆滑，任何人只要给人一种“滑头”的感觉，便成为别人心目中“狡猾的人”，注定没有前途可言。

而圆通是我们十分喜欢的，但可悲的是，现代人脑筋逐渐扁平化，缺乏深度，看事情不能“入木三分”，所以“凡是圆通都看成圆滑”。圆通和圆滑虽然长得一模一样，表面上看起来，几乎不分轩轾，都是“推、拖、拉”，一副打太极拳的模样，但是二者的本质不同，主要有以下几点差异：第一，过程完全相同而结果截然不同。“推、拖、拉”到最后，把问题圆满解决掉，便是圆通；“推、拖、拉”到最后，没有解决问题，或者解决得不够圆满，那当然是圆滑。第二，圆通的人善于利用“推、拖、拉”的短暂时间来充分思考，寻求此时、此地合理的行动方案，以便减少阻力，使大事化小，小事化了；圆滑的人只想利用“推、拖、拉”把时间拖延过去，以便逃避当前的问题，根本不想解决问题。第三，圆通的人用“推、拖、拉”来降低竞争的压力，使大家觉得经过“推、拖、拉”之后得到的答案，应该可以接受；圆滑的人，则冀望于不经由竞争，便能够获得胜算，或者不愿意应付挑战，悄悄地逃之夭夭。

圆滑与圆通有着本质的不同，但是很多人盲目排斥“推、拖、拉”，视“推、拖、拉”为落伍、陈旧、腐败的手法，把圆通和圆滑都看成可恶的行为，因此失去圆通的本事，真是得不偿失。其实，合理地“推、拖、拉”，把“推、拖、拉”的功夫发挥到出神入化的地步，这就是圆通。如果一味地认为“推、拖、拉”是坏事，那就会事事看不惯，甚至整天不愉快。

圆通就是“面对现实”，“负起责任”，另外加上“不伤害面子”。不够圆通，很容易伤害自己的面子，也很容易伤害别人的面子。

譬如介绍别人的时候，经常听见当老师的人说“这是我的学生”。而被介绍的人，却似乎并不认定介绍人是老师，反而淡淡地说“我们现在合作一些事情”，这岂不是很尴尬？当老师的人，如果介绍“这是我的朋友”，对方紧接着说：“不，我只是教授的学生而已。”是不是就能避开尴尬的场面？让别人承认自己，而非自己盲目乱认，这就圆通多了。

“面对现实”和“负起责任”可以用理智和科学来处理；“不伤害面子”却需要以感性和艺术来协调。圆通就是科学加艺术，理智加感性，所以很难做到。

我们承认“圆通是高难度的素养”，一般人很难精于此道，何况一不小心，就会流于圆滑而遭人唾弃。但是，难学，还是要学，因为只有学会圆通，才能在人际交往中游刃有余。

圆通和圆滑仅仅一字之差，但是差之毫厘，谬以千里。比如，听到别人的夸奖，要针对不同的人回复不同的话：对西方人要说“谢谢”，对日本人要说“请多多指教”，对中国人要谦逊地说“哪里，哪里”。这种“见人说人话，见鬼说鬼话”的行为，就是圆通。如果对西方人说“哪里，哪里”，相信对方一定莫名其妙；如果对中国人说“谢谢”，对方会不以为然：“我不过说说而已，你还当真了。”

不可否认，圆通有圆滑的成分，但也含有不圆滑的成分。我们先接受圆通的概念，再观察圆通的事实，分析圆通的要素，才可以学得圆通的精髓。

第一，将心比心是圆通的先决条件。多以欣赏的眼光来体会他人的圆通，比较容易吸取他人的经验，迅速成长。常以厌恶的心情来批评他人的圆滑，结果便只有生气的份儿，失去学习他人宝贵经验的机会。

第二，不要完全排斥或放弃“推、拖、拉”，也不能够凡事都“推、拖、拉”，以免一不小心，就变成令人厌恶的圆滑。

不过，中国人自己“推、拖、拉”的时候，往往不认为自己正在做这种令人厌恶的动作，却自满于“事缓则圆”的解说。但是，看到别人“推、拖、拉”，却又立即觉察出来而心生反感。

圆通说起来相当简单，做起来难。因为圆通与圆滑的差距十分细微，稍不留意，就会“失之毫厘，谬以千里”。

\section{遇事先求自保}

“明哲保身”是中国人普遍秉持的人生哲学。朱子说：“明是明理之明，哲乃了然于心，保身则是依道理进退而处世，于是罪刑、灾祸不至。”

中国人凡事先求自保。遇到突发状况，就会遵循明哲保身的哲学分析形势，当进则进，当退则退。形势不利时，就要及早抽身，“留得青山在，不怕没柴烧”，何必逞匹夫之勇呢？很多人反对“明哲保身”，认为善于明哲保身的人怕死、虚伪、消极，其实不然，一个人连自己都保护不了的话，其他的都没有意义。人们常说，“泥菩萨过河，自身难保”，神仙都要先保全自己，何况凡人呢？中国人明哲保身主要表现在以下几个方面：


\section{深藏不露}

深藏不露实际上是中国人守身哲学的应用，也是我们奠定“立于不败之地”的良好基础。如果有人把自己的实力毫不隐瞒地展现无遗，很容易遭人算计，岂不是自讨苦吃？中国人深知“木秀于林，风必摧之”的道理，宁可韬光养晦，安然度日。

这也是中国人与众不同的心理。日本人就算没有本领，也要装成很能干的样子；美国人有多少本领，都要找机会表现出来；只有中国人深藏不露，静观其变，再衡量情势的变化，以对自己最有利的做法，做到“不鸣则已，一鸣惊人”，立于不败之地。

中国有句俗话“树大招风”，因此，中国人有能力并不表现，还喜欢伪装成没有本领的样子。例如，会写字的人不轻易宣扬自己擅长书法，会画画的人不愿意让人家晓得自己善于作画。这样做比较省力省事，否则请字求画的人络绎不绝，只能疲于应付，稍有不慎，厚此薄彼，就会得罪人。另外，自己深藏不露，还可以留下力气嘲笑爱表现的人。中国人不讲能力，却十分重视本事。只有能力而没有本事的人，往往一表现就遭殃。有本事的人，能够把能力表现得恰到好处，受人欢迎。

比如，大家在打球，招呼你也参加，就算你天天练习，也要谦逊地说：“好久没有打了，生疏啦。”大家让来让去，你才登场，不费吹灰之力，把对手一一打败。就算如此，也没有人觉得奇怪，因为大家明白你刚才不过是深藏不露的一种应用罢了。

喜欢保留实力、不强求出头并不说明中国人奸诈、不实在。相反，谦虚会给人一种高深莫测的感觉。

深藏不露有三个要点：

第一，人怕出名，猪怕壮。人愈出名，遭遇的麻烦愈多。越有能力，承担的责任越重。中国人比较喜欢打击有能力的人，同情没有能力的人；能者必须多劳，没有人会感激，而且容易招人嫉妒。因此，不如隐藏实力，凡事量力而为。

第二，不要夜郎自大。中国人对那些自我膨胀，把自己看得很伟大的人，大多抱着敬而远之的态度。背后则嘲笑他们夜郎自大，不自量力，十分看不起。

第三，慎防人上有人。即使自己的本领高强，也要慎防人上有人，强中更有强中手。所以要深藏不露，以免泄露自己的实力，招来对手的攻击，对自己不利。

深藏不露并不是不露，而是站在不露的立场，来求取合理的露，以免露得过分或不及，对自己有害。有能力必须合理地表现，凡事量力而为，才能恰到好处。中国人主张“不可以随便表现”，并没有反对合理的表现。不能不表现的时候，务必表现到合理的地步，不可任意逾越合理的范围，以免留下后患。

\section{以让代争}

人类社会，彼此难免要竞争。实际上，了解中国人性格的人，就知道中国人太喜欢争了，而且一争起来往往不择手段，所以不鼓励中国人竞争，认为能够不争的，大家都不要争，不能够不争的，便以让代争，这才是中国人推崇的君子之争。

争是必要的，舍也是必要的。什么都要争，最后就会迷失自己，不知道自己争的是什么，为什么要争。而且给人一种恶劣的印象，使大家提高警觉，甚至于联合起来，一点机会都不给你。因此太会争、太喜欢争的人，经常什么都争不到。只想到舍，同样会迷失自己，不知道自己要的是什么，不要的又是什么，而且给人一种消极、不愿意负责任的不良印象。

只有把“争”和“舍”融合在一起，该争的才争，不该争的就舍弃，争到好像没有争一样，才是圆满的做法。以让代争便是兼顾“舍”与“争”的权宜措施，站在“不争”的立场来“争”，才不会乱争，才可能争得恰到好处。

但是“让”并非消极的“让”，也不是让给谁都行，而是在冷静思考之后，让最合理的人来做最合理的事，才不会令众人失望。千万不要心里不情愿，却打着谦让的幌子，那是虚伪。让来让去，觉得自己最合适，这时候就应当仁不让。关键在于把握“当仁”的尺度，要从事情的性质、轻重、缓急、大小来判断。

当然，当仁不让也要以“半推半就”的形式表现出来。“半推半就”就是指确实“当仁”，而且推辞不掉，实在没有办法，才勉强为之。“半推半就”必须诚心诚意，另一方面表示“我虽然勉强接受，仍然随时准备交给比我更合适的人”，一方面让没有当选的人也有面子。比如你被选为部门主管，必须要表明你不想当，极力推辞，说自己能力不行，经验不足，体力也不好，做事情没有耐心，把各种理由都列出来，最后盛情难却，只好“勉强为之”。这样做也是为自己留条后路，但千万不要以此为借口偷懒。一旦你接受了，就必须负责到底，才不辜负大家对你的期望。

\section{秘不外传}

中国人喜欢将自己的绝技保密，就算传授给徒弟也要“留一手”。这一直是受人诟病的一种作风，就因为“留一手”，造成很多技艺失传。其实，这原本是中国人的一种自保心态，并非完全没有道理。“留一手”本身并没有好坏，关键在于“留得合理不合理”。

不可否认，从古至今，凡是能够发扬光大的，中国人并不主张“留一手”，反而毫无保留地让全世界的人来分享。但是，有些东西，中国人之所以“留一手”，有一些特别的道理。

技艺失传似乎要归咎于师父的自私，不肯把自己的真实本领倾囊传授，以致一代不如一代。但是却不能忽略两种情况，一种是教会徒弟，饿死师父，另一种是教会徒弟，杀死师父。年富力强的师父凭着威势和丰富的经验，或许能够制服徒弟；而年迈体衰的师父，如果不幸遇到徒弟翻脸不认人，就会悔之晚矣。“年老慎择徒”就是这个道理。在这两种情况下，预先留一手，也是一种自保的方法。如果遇到忠良可靠的徒弟，师父自会倾囊相授，哪里会留一手?

中国人很重视拜师礼，凡欲拜师者必须恭恭敬敬地向师父叩首，以显其诚意，要不然师父怎敢把功夫传给他？

儒家一再希望中国人“反求诸己”，徒弟发现师父秘而不传，实在不应该埋怨师父，而是应该自我检讨，为什么自己在师父眼中这么不可靠，弄得师父必须痛心地留一手？而为师者，应尽可能将技艺发扬光大，并想办法找到合适的传人，教他、导正他，让他刚正不阿，这才是秘传的真义。

事实上，“留一手”跟“不留一手”，基本上没有什么差异。想学的人，师父留一手照样学得很像；不想学的人，师父不留一手照样学不好。“留一手”不过是学得不像、学得不好的人常用的一种借口。依据中国人的观点，用“师父留一手”来骗别人，给自己留下面子，尚属可行。若是连自己都骗，真的把责任都推给师父，那就实在太不长进了。“留一手”是说给别人听的，不是说给自己听的。

现在，除个别门类外，磕头拜师的情形越来越少了，但是在工作中，师父带徒弟的形式却很常见。这就要求人们在“留一手”和“不留一手”中做好抉择，既要有足够的能力自保，又要将自己的技术发扬光大。

\chapter{性——爱占便宜，死要面子}
中国人的本性是什么？简单来说有三点：第一怕吃亏，中国人天不怕、地不怕，就怕吃亏上当；第二喜欢占小便宜，中国人一般不奢望去占大便宜，专占小便宜；第三，自私又爱面子。

\section{就怕吃亏上当}

说中国人怕吃亏上当，并不是贬低中国人，凡是有这种想法的人都不太了解中国人。中国人重视合理，合理地怕吃亏有什么不对？资源是有限的，人的生命也是有限的，不断地吃亏，岂不是跟自己过不去？其实每一个人都有这个特性，凡是宣称钱财乃身外之物的人，多半都很小气。因为一个人如果心里没有某个弱点的话，就不会强调它。例如一个人缺乏信用，就会强调自己重信用，通过不断地强调，才能让别人上当。

就因为中国人怕吃亏，所以和中国人打交道时，千万不要有欺骗的意图，以互利互惠为原则，才能皆大欢喜。一开始就想欺骗中国人，往往会“偷鸡不着蚀把米”，得不偿失。

每个人都怕吃亏，但是过分地怕吃亏就变成精打细算了，其实人算总是不如天算，“机关算尽太聪明，反误了卿卿性命”。所以中国人发展出一套“差不多主义”，差不多就是刚刚好，就是恰到好处，就是合理。

因为怕吃亏，中国人处处小心谨慎，这就是中国人所说的“防人之心不可无”。但是，如果防不胜防，真的吃亏了，中国人也不以为意，反而相信吃亏是福，“吃一堑，长一智”。

\section{只爱占小便宜}

如果说怕吃亏是被动的，那喜欢占小便宜就是主动的行为。占小便宜，关键词不是“占”而是“小”，只占“小”便宜。一个人爱占“小”便宜，就表示这个人不贪，贪心的人只想着占大便宜。中国人爱占“小”便宜，不贪心，很守本分，为了过日子，占一点小便宜，没什么可耻的。而且对中国人来说，占小便宜只是手段，不是目的。中国人之所以这样做，是在试探对方。和别人做生意也好，交朋友也好，首先要了解对方，看对方是不是值得打交道的人。

中国人比较在乎舍不舍得，换句话说，中国人不会盲目地为一个人去拼命，我们会考虑值不值得。要判断值不值得，就要看对方舍不舍得，你舍得我就认为值得，你不舍得我就认为不值得。这就是中国人占小便宜的本意。

在人际交往的时候，首先要看对方的度量大不大，没有人愿意跟一个小肚鸡肠的人做朋友。看对方舍不舍得，舍到什么程度，可以判断你在对方心目中是何分量。如果对方连鸡毛蒜皮的事都要和你计较，那他也就没把你放在心上。

《战国策》中有这样一个故事：

齐人有冯谖者，贫乏不能自存，使人属孟尝君，愿寄食门下。孟尝君曰：“客何好？”曰：“客无好也。”曰：“客何能？”曰：“客无能也。”孟尝君笑而受之曰：“诺。”左右以君贱之也，食以草具。

居有顷，倚柱弹其剑，歌曰：“长铗归来乎！食无鱼。”左右以告。孟尝君曰：“食之，比门下之鱼客。”居有顷，复弹其铗，歌曰：“长铗归来乎！出无车。”左右皆笑之，以告。孟尝

君曰：“为之驾，比门下之车客。”于是乘其车，揭其剑，过其友，曰：“孟尝君客我。”后有顷，复弹其剑铗，歌曰：“长铗归来乎！无以为家。”左右皆恶之，以为贪而不知足。孟尝君闻：“冯公有亲乎？”对曰：“有老母。”孟尝君使人给其食用，无使乏。于是冯谖不复歌。

这个故事是说，冯谖投靠到孟尝君门下，整天无所事事，却嚷着要锦衣玉食、宝马香车。不要以为冯谖是看孟尝君好说话，想多从他那里捞点好处。其实，冯谖一直在考验孟尝君，看他是不是值得效忠之人。所幸孟尝君通过了考验，后来，“孟尝君为相数十年，无纤介之祸者，冯谖之计也”。

但是，中国人不会轻易去占大便宜，一方面是因为中国人相信天上不会掉馅饼，如果有大便宜可占，中国人首先会想到会不会有问题，会不会上当。另一方面，中国人讲究无功不受禄，要等价交换，得到大便宜必须要付出大的代价才行。

相比之下，平原君就没有孟尝君那么聪明了。当时，秦国派大将白起攻打韩国，占领了韩国的一块土地野王。在野王邻近有另一块土地上党，其官员看到野王轻易地就被秦军攻下，怕上党也守不住，就写信给赵国，表示愿意归顺，希望得到赵国的庇护。结果平原君认为，不费一兵一卒，就可以得到上党这块“肥肉”，为什么不要呢？虽然有人反对说：“就是因为不花力气得到好处，轻易要了，恐怕会招来大祸。”但是，赵王和平原君却不听，把上党划为自己的领地。秦国知道后，认为赵国存心和自己作对，就命令白起率大军去攻打赵国。结果赵国的四十万大军全部被秦军歼灭，国都邯郸也被围困，险些被灭国。

这就是前车之鉴。所以，和?国人打交道，首先要明白，在合理的范围内，中国人比较喜欢占小便宜，但是也明白“礼下于人，必有所求”的道理，因此想要中国人帮忙，只要施以小恩小惠即可。如果施以“大恩大惠”的话，会让中国人心生警惕。

\section{自私又爱面子}

怕吃亏和喜欢占小便宜是出于自爱的心理，很多人据此说中国人自私。确实中国文化强调自爱，很少讲爱人，不像西方人一开口就是爱人，而不讲自爱。不自爱的人，心中没有爱，拿什么去爱别人？一个人要先爱自己，让自己的心中充满爱，然后才能去爱别人。西方人讲“爱人”只是嘴上说说而已，中国人才是发自内心的。但中国人不会爱所有的人，我们首先要判断这个人值不值得爱，值得，才去爱。合理的自私就是自爱。

老实讲，中国人比较重视面子。凡是说“面子不重要”的人，是因为他不了解人性。一个人活得连面子都没有，那活着是为了什么？人与动物的最大不同，就是人有面子，动物没面子。有人认为，中国人爱面子是缺点。我从不认为爱面子是缺点，一个人连面子都不爱，那他还会爱什么？面子就是羞耻之心，一个人连羞耻之心都没有，那就无药可救了。

爱面子，往好的方面发展，乃是重视荣誉的表现，没有什么不好。若往坏的方面发展，则是爱慕虚荣；更严重的，就是爱面子爱到不要脸的地步，那就是本末倒置。“面子”是“情”，“脸”则是“理”。中国人讲情理，以理为本，视情为末。换句话说，爱面子要爱到不丢脸的程度，不丢脸就是合理；爱面子爱到合理的地步，才是合情合理。

有些人为了自己的面子常常说一些谎话，比如家里穷的人，只能吃得起青菜豆腐，就会宣称吃素食有益健康，所以自己只好常吃青菜豆腐。钱锺书先生在其所著的小说《围城》中，就形象地描写了中国人这种好面子的本性：“汪先生得意地长叹道：‘这算得什么呢！我有点东西，这一次全丢了。两位没看见我南京的房子——房子总算没给日本人烧掉，里面的收藏陈设都不知下落了。幸亏我是个达观的人，否则真要伤心死呢。’这类的话，他们近来不但听熟，并且自己也说惯了。这次兵灾当然使许多有钱、有房子的人流落做穷光蛋，同时也让不知多少穷光蛋有机会追溯自己为过去的富翁。日本人烧了许多空中楼阁的房子，占领了许多乌托邦的产业，破坏了许多单相思的姻缘。譬如陆子潇就常常流露出来，战前有两三个女人抢着嫁他，‘现在当然谈不到了’。李梅亭在上海闸北，忽然补筑一所洋房，如今呢？可惜得很！该死的日本人放火烧了，损失简直没法估计。方鸿渐也把沦陷的故乡里那所老宅放大了好几倍，妙在房子扩充而并不会侵略邻舍的地。赵辛楣住在租界里，不能变房子的戏法，自信一表人才，不必惆怅从前有多少女人看中他，只说假如战争不发生，交涉使公署不撤退，他的官还可以做下去——不，做上去。”

类似的话，相信每个中国人都听过，但只要无伤大雅，完全可以听之任之。只要不存心害人，或者对别人构成不便，用谎话来维护自己的面子，大家很容易谅解。见到有人适当地往自己的脸上贴金，你大可一笑置之，不必拆穿他的谎言，否则他丢了面子，会恨你一辈子。但是如此的“大话”要在合理的范围内，否则就成了夜郎自大，惹人嘲笑，那就真成了“不要脸”了。中国人常说任何事情都不可以过与不及，可见合理的爱面子才属正当；一旦过分，害处之多，为害之大，便不是三言两语所能说尽的。

保留别人的颜面，对自己有百利而无一害，因为对方会心存感激，因而会想方设法报答你。有些时候，给别人留面子，比给他金银珠宝更有用。

战国时，有一次，楚庄王战胜邻国后大宴群臣，席间叫妃子们代他向大臣们敬酒。其中有一将军酒后失态，趁年轻貌美的宠妃敬酒之机调戏她，恰巧一阵大风吹熄了所有的灯，黑暗中宠妃机智地摘下了那人头盔上的红缨以为凭证，哭着请求大王处理非礼者，并交上红缨为据。但是楚庄王并没有追究，而是让所有的将军都摘下帽上的红缨，并让太监收起来。灯再次亮起的时候，就分辨不出调戏者是谁了。而那位将军明白，楚庄王之所以这么做，是免得他当众出丑。后来，在一次战斗中，

楚军被许多的敌兵围困，陷于欲进不能、欲退无路的境地，楚庄王大急，就在此时，旁边有一员大将杀出，威猛无比，奋不顾身，杀得周围的敌兵四处溃散，最终挽救了这场战斗。而这员猛将正是那天那位被宽容的将军！

上司问下属：“交代你写的计划书写好没有？”下属心里可能想：“糟糕，晚上忘了写了。”但是碍于面子问题，嘴上却说：“写好了，只是早上急着要准时上班，忘带了。”上司明白，这是下属在找借口，如果说破的话，下属面子上过不去，以后可能处处都要与自己唱反调，所以只是淡淡地说：“哦，那明天别忘了带过来。”

下属因说谎暂时保住了自己的颜面，下班后自然会赶快把它写好，放在公文包里，第二天一上班，就把它交上去。下属不但完成了工作，而且吸取了教训：这一次差点惹麻烦，幸亏平时信用还不错，勉强抵挡过去，下一次不可以再犯，以免被拆穿了，不但难堪还可能受罚。

这样，既可以保证计划书尽早完成，又不会伤了彼此的和气，何乐而不为？

说谎话，找借口，有时候是为了保留对方的面子，促使其好好表现反省改进。有时则是为了保留自己的颜面，促使自己好好检讨，加倍努力，以求表现得更好。

说谎话骗人又保全自己的面子，一般来说，无可厚非。因为骗骗别人，只要不伤害到对方，对中国人而言，原本是平常事。最要紧的，便是暂时保留颜面之后，千万不要忘记赶快检讨反省，以便及时改善，寻求合理的补救。不能骗别人之后，自己也愈听愈相信，那就会使自己不知反省，不求上进，反而害苦自己。比如，有人问：“你们是同学，而且上学的时候他还不如你，怎么现在他升了经理，你却只是个普通职员啊？”“他呀，擅长阿谀奉承，又会钻营，当然升得快。”用这种话抵挡一番，先保住自己的面子，然后再好好搞清楚原因，以便切实改进，迎头赶上，才算合理。你骗别人，信不信，是他的事，理应由他自己承担所有的后果。骗自己，岂不自作自受？根本没有人会同情，反而惹人笑话，当然要极力避免。

和爱面子的人打交道，只要注意不让他丢面子就行了，当然，如果想办法让对方有面子，那就更好了。适当地恭维一下别人，多说一些对方爱听的话，让对方觉得面子十足，他就会变得很好说话。切记，使一个中国人觉得没有面子，吃亏的并不是他，而是我们自己。

合理地怕吃亏、合理地占小便宜、合理地自私、合理地爱面子没有什么不好，如果不承认这些的话，就是假道学。

\part{处世的基本心法}

\chapter{正视自己}
自己想做什么样的人，就会建立什么样的人际关系。现代的价值观是，尊重个人的价值取向，把每一个人依自己的特长、志趣而从事的活动，都视为正当。做正当的事，就是正人君子。

一切从自己做起，才有成功的希望。偏偏现代人“多要求别人却很少要求自己”，大家所注意的，是别人如何如何，很少反省检讨自己有没有做得不好的地方?

中国人最讨厌彼此骗来骗去，也不承认自己会欺骗别人，但是实际情况却显示中国人常常骗来骗去。相信你一定知道这种情况：

电话响了，小王拿起电话，电话那边传来一个男人的声音：“您好！我姓张，要找李先生，请问他在吗？”小王用手捂住话筒，对李先生说：“李先生，有你的电话，是一位张先生打来的。”

李先生一听，连忙告诉小王：“你告诉他，就说我不在。”

这种话不是欺骗，而是中国人惯用的一种“设计”，如果此事发生在你身上，你就要自己检讨，为什么做人做到这种地步，明明人在这里，却不愿意和你见面？

人的一生是自己创造出来的，希望成为什么样的人，就应该一步一步按照自己的希望，认真地做好自己。这是儒家“反求诸己”的主张，也是“自己负起完全责任”的具体表现。

一表人才能成为十大要领之首，就说明第一印象的重要性。人与人之间，只要彼此来往过几次，就会互相产生某种评价。人们大多习惯于利用过去的经验来判断其他的人，同时，也习惯于和熟悉的人打交道。

第一、二次信用良好，就可能被对方视为有信用的人。所以中国人常说，好借好还，再借不难。中国古代的商帮将信用视为生命也是这个道理。开过一两次玩笑，以后所说的话，常常被当做笑话。我们从小就是听“狼来了”的故事长大的，对这种情况再熟悉不过。

每一个人迟早都会被人贴上一张看不见的标签，清清楚楚地写着他的性格和人品。许多人更是一见面，就喜欢论断对方究竟是什么样的人，所以第一印象非常重要。

\section{确定自己要做什么样的人}

我们要替自己负起责任来，因为我们有权力决定以后的路要怎么走，我们要对自己的所作所为负责，因此人际关系要从自己做起。在开展人际关系前，首先要想想自己要做什么样的人。

人大致上分成三种：第一种是忠义之士，为正义而牺牲在所不惜，这种人又叫做硬汉。别人不敢说的他说，别人不敢做的他做，这就是硬汉。我很佩服这样的人，但我不希望人们都做这样的人，除非他自愿。做硬汉，要承受一些痛苦。比如，你要受很多的苦难，否则怎么知道你是硬汉？此外，你不能变节，做硬汉最可怕的就是变节，即做到一半不做了。做人应该有正义感，但是小心被利用。如果你所知有限，又充满正义感，就很容易被别人拿来当工具。每一个时代都有硬汉，如果说他们真的拿捏合理，就会流芳百世，否则就是无谓的牺牲。人各有志，不能勉强，我们应该尊重每个人对自己的选择，不鼓励也不反对，因为社会需要这样的人。

第二种人是顺民，即唯唯诺诺、听话、保平安、少惹事。老实说，不是硬汉，非要他当硬汉，他很辛苦；不是顺民，非要强制他当顺民，他也很痛苦。在一个组织里面，一定要有顺民，也一定要有硬汉。

如果这两种人你都不想做，那你还可以走第三条路，做隐士，即不闻不问免得生气，反正都是别人的事。这种人既不会有什么贡献，也不会惹什么麻烦。

在这三种人中，做硬汉很辛苦，做顺民划不来，做隐士也不容易。一个人要做到完全不闻不问是很难的。因为我们是社会的一分子，总希望参与社会事务，总需要有点作为，所以隐士其实不容易当。

那怎么办？好在中国人从来没有走投无路的时候，既然三种人哪个都不好当，那我就三合一。在中国文化中，除了“孝”以外，还有个重要的观念，就是“合”。中国人之所以聪明，就是把所有的东西都合在一起考虑。我们尽量当顺民，但是偶尔也会当硬汉，必要的时候又可以做隐士，这种人就是随遇而安。领导看得起我，那我就当硬汉；领导看不起我，我就当隐士，反正我讲的领导也不听；领导对我一般，那我就当顺民。中国人常常是看你对我怎么样，才决定我如何对你。同样一个人，有时候他什么都不管，有时候他很认真、很积极，而有时候他又好像很听话的样子。

该当硬汉的时候，就当硬汉。当年史可法不可以逃走吗？当然不是，带兵打仗的人想逃太容易了，但是他一跑，这辈子就白活了，如果不跑，就可以万古流芳。人迟早是要死的，要死得其所，要死得其时，这是你可以选择的，所以史可法选择战死沙场。岳飞也是一样，如果他不想回去，别说十二道金牌，就算二十四道金牌也没用。结果他回来了，身遭不测却流芳百世。

人生不过是四个字而已——心想事成。我们不要把“心想事成”当做祝语，它本来就是事实。但心想事成不意味着你一直说想成功就能成功，老天听不懂人的话，整个宇宙都是信息场，老天只接受你的信息，但听不懂你的话。当你发出“我要成功”的口头信息时，老天所接收的信息是“这个人没有成功”，然后就如愿以偿让你不能成功。

什么时候你能掌握宇宙信息场的变化，你就可以心想事成。现在有一种说法是心灵管理，即学会怎样去控制对方的心灵。我相信有一种力量，是到目前为止科学所不能解释的，叫做念力。宇宙非常大，我们已知的部分很小，如果已知的部分是科学的话，未知的部分就是神，如此而已。

每个人要建立自己的价值观，把自己的价值观发散出去造成一种磁场，物以类聚、人以群分，频率相同的人就是志同道合的人，这个就是所谓的缘分。缘是机会，分就是关系，有缘无分就是说，两个人有机会，但始终无法建立关系。

人与人之间要保持适当的距离，否则物极必反，过分亲密就容易吵架。“君子之交淡如水”是有道理的，因为两个人就算是唇亡齿寒的关系，牙齿迟早也会咬到嘴唇的。

\section{处理好自己和自己的关系}

要想和别人搞好关系，首先要同自己搞好关系，要能够接受自己。只有接受自己的人，才能使自己的身心得到充分的发展，因而获得和谐的人际关系。你跟自己相处不好，就不会有人跟你处得好；你跟自己处得好，别人才会跟你处得好。所以人要首先了解自己，跟自己相处好，而不是只看到别人。每当面对着镜子的时候，问问自己对镜中人究竟是爱是恨?一个不喜欢自己的人，大概很少有人会喜欢你。可事实上就是有大部分人不喜欢自己，总觉得自己这儿长得不好看，那儿也长得不好看。要改变别人只有一个办法，就是改变你自己，而改变自己不是说去整容，而是接受你自己。当你哪一天接受了自己，就跟自己相处好了，那时你会发现，所有人慢慢地都跟你处得很好，都开始接受你了。

憎恶自己的人，必然也憎恶别人。不能接纳自己的人，在情绪上常常显得很不稳定，不是有意表现优越，便是相当自卑。这种内心的摩擦，使得不能接受自己的人，同样也会憎恶他人。

如果发现自己的人际关系并不好，不妨反省一下自己和自己的关系如何。先调整自我关系，然后改善人际关系，才是有效的途径。人具有相当程度的自主性，这是人类和其他动物最主要的差异。

既然能够自主，一切由自己决定，当然要由自己承担所有的责任。换句话说，一切的言行，事实上都要先通过“自己”这一关。自己认可的，才说得出来；自己认同的，才做得出来。所有接触的对象，也由自己来决定。人，最先也最多接触的，应该就是自己。

接受自己必须在合理的范围内，过分地爱自己，有时会成为可怕的自恋狂。自恋狂最大的特征是，完全以自我为中心。这种人不但不能欣赏别人的优点，而且从来不怀疑自己很可能具有某些缺点，以致自认为十分可爱。我们最好冷静下来，客观地反省自己究竟有哪些优点，又有哪些缺点。发现自己有缺点，只要实实在在努力去改善就可以了，不必过分憎恶甚至嫌恶自己，造成自我否定，反而容易成为坏人。

善于发现自己的缺点，就能够把它变成优点。接受自己有某些缺点的事实，尽量不要让它发展，进而将它转变为自己的优点，才是最好的办法。贝多芬后来变成聋子，照样谱写传世名曲。人都不是完人，有缺点是十分自然的现象。人没必要变成完人，“完人”表示把人做完了。对于自己的缺点，固然不可忽视，却也不必紧张。因为紧张无济于事，并不能解决问题。

把缺点变成优点，听起来有一些奇怪，怎么会这样?其实任何言行，配合时空的变迁，调整到合理的地步，便是优点。离开了具体的时空，本来就没有什么优劣可言。

尽量减少缺点增加优点，并不意味着让所有人都说你好，想得到所有人的赞美本身就是贪心的想法。大家都说好，未必真的好；大家都说坏，也未必真的坏。特别是多元化社会，同样一件事，有一个人说好，就可能有五个人说坏。我们不能够单凭人家的论断来判定这件事情的好坏，却应该看看赞成的是哪些人，反对的又是哪些人，然后再谨慎地进一步判断。

不要太过在乎别人的批评，听到有人说自己的坏话，要看看说的人是谁。希望每一个人都说自己好的人，往往过分害怕得罪人而成了乡愿，并不是大家所喜欢的对象。

因此，听到不同的声音，要先明辨是非，因为不同的人立场常常不一致。凭良心，不讨好，也不刻意为了凸显自己的好而得罪小人，这才是上策。

\section{以实实在在为做人的准则}

中国人从小就被教育“做人要实实在在，做事要规规矩矩”，这是中国人安身立命的基本原则。中国人虽然处世追求圆通，但始终以务实为修身之本。中华文化一直强调“君子务实”，务实的具体表现为，做人重诚信，不伤害其他人，并且以诚恳的态度对待别人。

知错能改，善莫大焉，一个人能认识自己的过错，进而改变自己的行为，正是务实、务本的实践，也是对待自己的最佳途径。

从古至今，获得成功的人不外乎两个途径：一是正道，一是偏锋。以务实的方式来建立人际关系，属于正道；用欺诈的方式来建立人际关系，即为偏锋。循正道成功，才是实至名归，值得敬重。因偏锋而成功，不过是欺世盗名，并无多大价值，徒然惹人背后耻笑。中国人厌恶权术而欣赏艺术，便是由于艺术才是务实、务本的方式与技巧。

成功离不开实实在在做人、规规矩矩做事，但只是实实在在做人、规规矩矩做事并不一定能成功，这只是基础工程。就像打地基一样，地基稳，才能在上面建造高楼大厦，但地基并不等于高楼大厦。所以务实只是本分，守本分之外，需要进一步持经达变，培养自己的随机应变能力。

如何培养自己的随机应变能力？最好是多看、多听，先了解环境，再适应环境，然后才动脑筋改造环境，最后才有能力合理地创造环境。多看、多听、多问，并不一定只限于正面的、好的东西，对于那些负面的、不好的东西也要了解一下，这样有助于防患于未然。记住，此处只是教你多看、多听、多问，但要少说，正所谓言多必失。贸然说出一些话来，固然痛快，却也很快就要承受某些痛苦。但少说并不意味着不说，否则就是矫枉过正。当你看准了、想明白了再说话，才会言必有中，每一句话都合理，这样比较妥当，也会受欢迎。胡言乱语，不但让别人看不起，而且降低了自己的信用。

在不忘本的情况下权宜应变，才不致乱变。人际关系是不进则退的，就好像一株幼苗，需要时时浇灌才会茁壮成长，否则的话就会枯死。如果不能随时注意调整，久而久之，人际关系只会转坏而不可能转好。培养自己的应变能力，因时、因地而制宜，才会增进与别人的良好人际关系。

务实的同时必须适当地调整，使根本稳固。因为所有的事物都是时时刻刻在变化的，必须富于改善意识，运用精锐的眼光，发挥自己的智慧，不断寻求改善。我们常说随机应变，就是说，任何事情都需要因人、因地、因事、因时而制宜，不可以一成不变，但是“不可乱变”，不可以为了求新求变而忘本。有所变还要有所不变，变来变去都能够务实，才是以不变应万变。

人际关系的建立还有赖于自己长期地培养声望。同样一句话，声望不够的人说出来，便是人微言轻，很少有人加以理会；换一位有声望的人说出来，马上显得有分量。

养望相当困难，不是一朝一夕所能达成的。培养自己的沟通力，是比较可行的途径。人际关系有许多地方离不开沟通，良好的沟通力，正是建立和谐人际关系的主要方法。沟通能力强，就是说话说得让对方听得进去，让对方乐于接受，能够引起对方的共鸣，进而引发共同的行为。

沟通力良好，才能在和谐的气氛中彼此协调。这时候更需要重视人伦的道理，使大家觉得这种尊重他人的态度，值得敬重，也能够给予相当的信任，久而久之，声望就建立起来了。这是急不得的，必须经过时间的考验。在中国社会，遇到争执的时候，经常会请几位德高望重之人出面调停。难道德高望重之人就一定公平吗？未必，但是大家对他们有信心，知道他们一定会公正，所以将其视为沟通的桥梁。不管他们的评判有没有道理，对立双方都比较方便下台阶。


\chapter{善待他人}
人是群居的动物，从一出生起，就避免不了同各种各样的人打交道，尤其是中国社会，人与人的关系很复杂，没有人能逃得过这张巨大而无形的关系网。那我们应该如何维系这张网，即我们如何对待他人呢？

\section{待人需要诚心}

朋友相交，贵乎交心，也就是说，知己朋友要心意相通、坦诚相待。诚信是自古以来人们推崇的美德，在人际交往中，如果没有诚信，就会变成尔虞我诈、钩心斗角，那还有什么意思？

人与人之间要相互了解才能产生信任，知己知彼才能立于不败之地，这是千古不变的道理。但是，要求不败，除了自己知己之外，还要让好人知己，不让坏人知己。然而，好人坏人很难分得清楚，所以不得不制造若干假象，来掩盖自己的实情。中国人通常会让自己的眼光收敛一些，不流露出真正的感情。

戴上面具会拉开人与人之间的距离，完全没有面具等于赤裸裸地站在众人的面前，反而容易引诱别人欺诈或利用你。所以我们必须小心，有时把弱点装成优点，有时又要把优点装成弱点，因为防人之心不可无，在没有弄清楚对方之前，以及弄清楚对方根本不怀好意之后，必须谨慎地保全自己。否则一直诚恳待人，反而为人所害，势必使自己从此对诚恳失去信心，结果变成一个奸诈阴险的人，岂非害了自己?

中国人的自保心理很重，做起事来难免“遮遮掩掩”。明明是个中高手，却装成一窍不通；喜欢的东西，假装不喜欢；想要的东西，先让给别人；不想来说“要来”，说“不来”却来了。中国人在虚虚实实之间，有很多变化。有些人说中国人骗来骗去，实际上中国人多数重视诚信，不过是因为亏吃多了，自然懂得遮遮掩掩的好处。有时候中国人并不骗人，只是暂时蒙一下，瞒一时，一方面明哲保身，一方面静观其变。出发点是保护自己，而不是欺骗别人，所以也不算是什么善意的欺骗，因而丝毫没有欺骗的意思。

人与人相处，用诚恳的态度来感动对方，固然是不易的原则，但是实际应用起来，这种单一方式很难奏效。有时候需要“设计”一点小花样。这些花样是建立人际关系必不可少的润滑油，可以使彼此之间减轻摩擦，促进互动。

耍花样的人，很容易惹人厌恶，引起反感。但是，完全没有花样，一切直来直往也不容易受人欢迎。举个例子，有人过生日，你直接送他礼物，他会感谢，但未必感动。如果你设计一些花样，让他体验到从未体验过的事，他可能会很感动。为什么大多数女人都喜欢浪漫，浪漫不过是在适当时间、适当的地点制造一些花样而已。

心正的人，就算花样一大堆，也称为艺术。心不正的人，如果花样一大堆，便是玩弄权术。艺术和权术只一字之差，但性质不同，艺术的基础是诚恳，否则便成为权术，成为破坏人际关系的杀手。前者以心正为出发点，后者则因心不正而表现出损人利己的行为。

对待一些人，可能正常的方法行不通，那就要想点小花样，但千万不可以耍花样害对方。比如，我们常说，请将不成，何不激将。激将法早已有之，即用恩威并施、刚柔相并的方法达到目的。运用得当，可以收到奇佳的效果。

对于自负的人，我们越是诚意请求，他就越发得意，不是托故拒绝，便是推三阻四，此时运用反激的方法，会屡试不爽。激将的要领，在于一切为他着想，例如担心他受害，恐怕他不能胜任，唯恐他受人摆布，激起他的自尊心，使他非要表现一下不可。

对于老于世故的人，反激若是没有把握，可以先把他捧得高高的，再暗示他此事非同小可，而且非他莫属，让他自告奋勇。这种正激的方法，使他很难推诿。这对于有能力而退缩一旁的人，具有奇效。

这些小花样，我们不应该排斥，否则很难彼此和谐。这些花样是人际技巧而非伎俩，与以诚相待并不矛盾，因为你的出发点是好的。

不过，权术和艺术，从形式上、表面上看起来，几乎一模一样，以致很多人分不清楚，也看不明白。有人把权术当做艺术，惹来很多反感。有人则把艺术看成权术，极力加以排斥，使得自己寸步难行，还被人视为不近人情。

我们既不可以盲目排斥花样，认为凡是有花样的都属权术，也不能够心术不正玩弄权术，把别人当做傻瓜，自以为聪明地耍手段、玩花样。因为这两种极端的态度，都不得其法，对自己都没有好处。

\section{待人需要友善}

对待别人只有诚心还不够，还要有善意。

子路曰：“人善我，我亦善之；人不善我，我不善之。”子贡曰：“人善我，我亦善之；人不善我，我则引之进退而已耳。”颜回曰：“人善我，我亦善之；人不善我，我亦善之。”三子所持各异，问于夫子。夫子曰：“由之所持，蛮貊之言也；赐之所言，朋友之言也；回之所言，亲属之言也。”

这段话的意思是：

子路、子贡、颜回在一起谈论待人之道。子路说：“别人以善意待我，我也用善意待他；别人用不善待我，我也用不善待他。”

孔子评价道：“这是没有道德礼义的夷狄之间的做法。”

子贡说：“别人用善意待我，我也用善意待他；别人用不善待我，我就引导他向善。”

孔子评价道：“这是朋友之间应该有的做法。”

颜回说：“别人以善意待我，我也用善意待他；别人用不善待我，我也以善意待他。”

孔子评价道：“这是亲属之间应该有的做法。”

人性原本有善有恶，事实上也可善可恶。一般人偏向性善，“人之初，性本善”，自己抱持善意，期待对方善意的响应，多半能够心想事成。因为中国人是交互的，你对他好，他没有理由不对你好。同样的道理，一开始就认定对方缺乏诚意，对方十分敏感，一下子就看出来，当然不会诚恳地响应我们，这是十分自然的事情。表现自己的善意，使大家对我们产生比较良好而深刻的印象，才能进一步建立关系。别人注意我们，却产生不良的印象，不如不让他注意，反而更好。所以，我们不但需要引起他人的注意，还需要造成良好的观感。

不仅自己向善，还要导人向善。朋友的美德，应该替他宣扬；朋友有缺点，要帮助他改正。闻过则喜是一种美德，可惜做得到的人委实太少，那我们应该怎么办？子曰：“忠告而善道之，不可则止，无自辱焉。”意思说，当朋友有过错的时候，要忠诚地劝告他，要好好地开导他，如果他不接受的话，那就要停止，不要去自取其辱啊！

指出朋友的缺点或错误时，要先赞美他，然后批评他，最后再赞美他，这样他才乐于接受。也可以借用别人的话来批评对方，譬如，“有人说你相当不近人情，但是我的感觉完全不是这样”或者“我看你整天忙于工作，偏偏有人还在批评你偷懒……”，这样就可以弱化批评的语气。

想改掉朋友的坏毛病，最好以身作则，让其耳濡目染，慢慢自己调整。因为朋友之间是互相影响的，我们常常告诫别人，谨防交友不慎，就是因为不好的朋友会给我们带来不好的影响。我们把这种朋友斥之为狐朋狗友，而把那些能带给我们好的影响的朋友称为益友。我们要影响对方，必须先了解他的背景和立场，主动给予配合，才能让他乐意接受。尽可能发现彼此相似的地方，最好能找出他的优点，逐渐喜欢他、关心他。当这些反应传给他以后，他也会有所回馈。彼此都伸出友谊之手，产生交互作用。先承认他言之有理，让他了解我们并不是存心要改变他。然后晓以道理，使他觉得如果换一个角度看也有另外一番道理。当他身心放松的时候，趁机让他自己改变，通常比较容易收效。

表达善意要多替他人着想，站在他人的立场，别人才可能依据交互精神，同样站在我们的立场来合理地回应。大家都将心比心，交集的范围加大，彼此有共识，当然容易沟通而建立良好的人际关系。任何人凡事只想到自己而不考虑他人，必然被视为自私自利而得不到他人的支持和欢迎。

想想自己，也想想别人，才能合理兼顾各方面的利益，合乎将心比心的原则。有一句话很通俗，“前半夜想自己，后半夜想别人”。不过有些人前半夜想自己，想着想着就睡着了，从来没有想过别人。

\section{待人需要礼貌}

中华文化非常重视“礼”。《论语》上说：“不学礼，无以应。”民间也有“礼多人不怪”的说法，可见，在人际关系中，礼貌是不可或缺的。

司马牛忧曰：“人皆有兄弟，我独亡！”子夏曰：“商闻之矣：‘死生有命，富贵在天。’君子敬而无失，与人恭而有礼；四海之内，皆兄弟也。君子何患乎无兄弟也？”

这段话的意思是，司马牛忧伤地说：“人人都有好兄弟，就单单我没有！”子夏说：“我（商）曾经听人说过，人的生死，冥冥中自有天定，人的富贵也完全是由上天安排。君子只要恭敬地修养自己而不犯过错，对待他人谦恭有礼，那么在四海之内的所有人，都是我的好兄弟了。君子又何必忧愁没有好的兄弟呢？”

礼貌的作用，一使大家相处愉快；二使大家在一起时，秩序井然；三为表现合乎自己的身份；四为发挥以让代争的精神。做一个谦恭有礼的人，别人自然会喜欢你而与你亲近，你又何必做司马牛之叹？

对人有礼貌，不见得马上有实际的效益。有些人急功近利，认为有没有礼貌，人家又不能把我怎样，何必约束自己，处处讲求礼貌?有些人注重形式化的礼貌，却丝毫不关心人，令人觉得此人很虚伪。就像美国人那样，他们惯有的礼貌只是一种虚假的表象，并不含有任何情感因素。

\section{待人需要差别}

虽然与我们打交道的人众多，但我们却不会一视同仁。中华文化中形容交情的词有很多：点头之交、泛泛之交、患难之交、刎颈之交、莫逆之交、生死之交……这些词多以交情深浅而划分，这就告诉我们，对待交情不同的人应该采取亲疏有别、差别对待的形式。

对只是认识但并不熟悉的一般人，我们仅限于礼貌性地打招呼，谈不上什么信任不信任的问题。逐渐互有交往，增进彼此的了解，成为熟悉的朋友之后，虽然也交换一些意见，但多是闲聊，毕竟交浅不言深，不可能涉及不足为外人道的题材。熟悉的朋友当中，有一些信得过的，慢慢变成好友，再进一步成为密友，这才放心商量一些私人的问题。密友经得起再三的考验，切实可以交心的，才会变成知己。

比如《三国演义》中的刘备，虽然将关羽、张飞、赵云都视为兄弟，但他对三人的感情并不同：关羽、张飞是他的结拜兄弟，三人自桃园结义开始，便亲如一家，刘备所说的“妻子如衣服、兄弟如手足”指的也就是关羽、张飞二人。而赵云的地位则不同，关羽曾说：“翼德吾弟也；……子龙久随吾兄，即吾弟也：位与吾相并，可也。”一句话点明了赵云的地位，不过是因为他战功卓著，加上忠心耿耿，所以从情分上视为兄弟而已。

其实每个人都如此，都会谨慎区分自己的知己、朋友、熟人，我们如此对待别人，别人也如此对待我们，这都无可厚非。所以，在人际交往中，我们还要掂量一下自己在别人心里是何分量。如果你和别人是知己，完全可以无话不谈，他有什么过错，你都可以委婉指出；如果你和别人只是熟人而已，则交浅不可言深，以免犯了别人的忌讳，埋下祸根。

要想交到知心的朋友，离不开感情的投入，即从对方的立场去考虑其想法，以了解其感受、要求和苦恼。这种感情投入是自发的，不应计较产出，你自然会得到许多朋友。如果你是为了获得回报才投入的，总是想着“我对你好，你一定得对我好”，那你肯定得不到他人的回馈。

人和人打交道不应该点到即止，如果你认识的人都只是泛泛之交，而无一个知心朋友，那又怎会有乐趣呢？我们认识别人，并不会只是抱着认识的目的，而是想通过交往熟悉起来，成为朋友。

\section{三大禁忌}

人际关系有三个破坏因素：锋芒毕露、自负和轻率。如果你有这三个毛病，最好马上改掉。

一个人锋芒毕露，其人际关系不可能好。很多年轻人都有这个毛病，到哪里都要变成焦点，别人讲话他要插嘴，对什么事情都有意见……这些都是锋芒毕露的表现。锋芒毕露不一定会出人头地，因为所有的人都会找机会把你的锋芒除掉。很多人年轻时有棱有角，后来却变得很圆滑，就是因为受了很多打击。

自负也是破坏人际关系的主要因素。很多人刚上大学的时候，眼睛长在头顶上，只看到天，没把别人放在眼里。上了二年级后，挨老师的批评多了，眼睛才慢慢长到下面来。

尤其现在把大学分等级后对年轻人的影响更大。不是一流大学毕业的学生，底气都不足。但是社会风气就是这样，而且全世界都这样，比如日本，只要一个人上了东京大学，他这辈子什么问题都解决了，职业有人给他找，老婆有人给他找，房子有人给他找，他什么都不用操心。这就造成了一些年轻人进入名牌大学以后，目空一切。结果在工作时，很难与别人合作，管理他也非常困难，除非你都听他的，而他又不成熟。所以，企业只好让他做可以独立完成的小项目。

西方人一直强调人要对自己有信心，但我认为，人应该对老天有信心。什么叫老天？老天就是自然，人是自然的一部分，必须服从自然的规律。

中国人讲究天人合一，伏羲氏当年就是从自然的景象体会出人应该怎样的。伏羲说，天、地、人，曰三才。意思是说，只要三个人在一起，其中有一个人像天，有一个人像人，有一个人像地，像天的人就要领导其余两人。皇帝自称天子，是因为只有他才能够代表天；地是很踏实的，它要承受所有东西，负荷力最强；人要承上启下、顶天立地。三个人在一起，角色不同，互相配合，才能相处愉快。

这样的话，每一个人都不能有锋芒，一切要顺乎自然，这是中国人一向的主张。花长得美，就会被摘掉；树长得直，就会被砍掉；人有才能，就会遭打击。所以一个人要修炼，要有内涵，不露锋芒，就会有前途；一个人很自负，自己认为了不起，那就叫做无法无天。孔子有一个学生，叫颜回，成绩最好，也是孔子最喜欢的学生，但他英年早逝，什么东西也没有留下，反而其他成绩差一些的学生做出了很多功绩。

轻率，就是做什么事情都不经过大脑。自以为很聪明，实际上就是轻举妄动。遇事不要立即反应，因为每次发生的事情都很特别，历史虽会重演，但从来都是大同小异，所以谁也没把握处理好。

很多人经常说要按照惯例办事，但每次的事都不同，何谈惯例？所以，同样的事结果往往不一样，有走运的也有不走运的。有些人开车闯红灯被抓，第一反应是倒霉，从来没有人主动承认“我违规了”。被抓的人肯定想：“前面的车闯红灯警察没看到，后面的车闯红灯警察没来得及抓，所以只有我倒霉。”

一般人遇到刺激就会乱说话，事后才后悔，而闭上嘴后，这个刺激会迫使大脑飞快地运转，积极思考问题，这叫做谋定而后动。千万不要急于显示自己很聪明而立即反应，否则你会后悔的。我们都不是神仙，每一个人都有一些缺点，最要紧的就是告诉自己：“我很有能力，但是我不要随便表现。”

因为我们不是西方人，西方人一生只做两个字而已，一个叫Show，Show就是作秀的意思，西方人有什么才能，一定要表现出来。另一个叫做Tell，Tell就是告诉别人我表现了什么。西方人不断地表现，而且到处去告诉别人，自己能做什么。这种人在中国是非常吃亏的。中国人遇到事情都先说“我不会”，然后看看大家都不会，才说“我来试试看”，结果做得非常好。

你问一个美国人最近在干什么，他什么都告诉你；你问一个中国人最近在干什么，他什么都不会告诉你，因为他的第一反应是“你问这个干什么”。你问一个美国人“你会打乒乓球吗”，他会就说“会”，不会就说“不会”；你问一个中国人会打乒乓球吗，他一定说“不会”。中国人不敢说“会”，否则对方会很认真地跟你打，把你打得惨败，然后还笑话你：“你这样就算会啊？那谁不会呢？”这不是自取其辱吗？

中国人为什么会立于不败之地？就是因为我们很会保护自己。比如，有人问：“对这件事你有什么看法？”中国人肯定说：“我不懂。”这样，就没有任何责任了。然后他听听大家的看法，听来听去觉得还没有自己懂得多，于是就站起来说：“我刚刚想到怎么办，让我试试看。”结果不言而喻。中国人没有确保安全的时候，不会轻易出手。

情势不停地在变，我们要特别机警，有任何风吹草动，先闭上嘴，赶快动脑筋，找到当前的合理点。全世界最受老天关爱的人有两种：一种是犹太人，一种是中国人，这两种人所受的磨难最多，但是最聪明。

全世界公认犹太人最会做生意，但有一个犹太人跟我说：“全世界最会做生意的是中国人，不是犹太人。中国人把‘最会做生意’的名号安在犹太人头上，人们看到犹太人就很警惕，结果犹太人根本赚不到钱，而中国人一直说‘我不会做生意’，结果赚得盆满钵满。”

所以说，一个有能力就表现出来的人，是没有什么前途的。


\part{人际交往“十要诀”}

\chapter{一表人才}
中国人从小就被培养建立正确的人际关系的意识，中国人的一切都是从人际关系开始的。成功的人际关系，是中国人共同的愿望。中国人普遍相信一条定律：在家靠父母，出外靠朋友。“靠”的意思，当然是“依靠关系”。于是靠得住或靠不住，便成为中国人十分重视的变量。靠不靠得住，必须经过比较长时间的考验。而朋友众多，选择靠得住的朋友的几率也比较大。中国人主张“四海之内皆兄弟也”，便是希望广泛地结交朋友，然后从中寻找知己。就算不能成为知己，朋友也总比陌生人要好得多。

很多人用其一生来寻觅若干具有良好人际关系的知己，以期成大功、立大业。如果说中国人一生都在努力建立靠得住的关系，大概并不为过。伯牙、子期虽然谱写了《高山流水》的千古绝唱，但在现代社会，像伯牙这样的人却不值得推崇。知己难逢，我们最好扩大交友的范围，然后由疏而亲，由浅及深，在众多的朋友中寻觅知己。

为了建立良好的人际关系，我们需要掌握一些技巧。庆幸的是，中国人善于把复杂的东西简单化。人际关系的建立技巧可以简化为十大要领，分别是一表人才、两套西装、三杯酒量、四圈麻将、五方交游、六出祁山、七术打马、八口吹牛、九分努力、十分忍耐。这十个要领是建立良好人际关系的要诀，缺一不可，必须合理应用，以确保成效。十大要领通俗易懂，学会以后一辈子都不会忘记。要领虽然只有十个，但内涵丰富，不可局限于字面上的意思。按照这十点去做，基本不会产生太大的冲突。

十个要领可以归纳成五个方面：人际关系的起点、人际关系的媒介、人际关系的交往、人际关系的技巧、人际关系的修养。这五个方面，中国人用一句话就可以概括，即一阴一阳之为道。这十个要点就是道，归纳成的五个方面都是一阴一阳、两两相对的，例如，一表人才是指先天的，两套西装是指后天的。

一表人才有三个重要的含义：第一，一表人才通常叫做第一印象。当你看到一个人的时候，首先会对他有个初步的印象，这就是第一印象。在现实生活中，千万不要用第一印象来论断一个人，这是很危险的做法。因为外表看起来很忠厚的人，内心可能非常阴险；外表看起来很聪明的人，实际上很可能是个草包。外表不能够代表人的全部，有很多人就吃过这个亏，只凭第一印象判断别人，这种做法太武断了。

第二，很多人都会凭第一印象来论断你。这就是一个事情的两个方面：你不要犯这个错误，但是你无法禁止别人犯这个错误。太多的人都是凭第一印象来决定他喜不喜欢你，甚至很多年轻人凭第一印象来找自己的终身伴侣，即所谓的一见钟情，结婚后才发现彼此根本不适合，结果只能草草结束婚姻。

第三，外表有利也有弊。并不是说人长成什么样，就一定有什么样的成就。人的脸孔是经常改变的，这是自然的转变，而不是要做整容手术。所谓相随心转，就是说心在改变，相貌就跟着改变。所以你跟别人见面时，不要再讲这种老话：“好久不见，你一点都没有变。”这是骂人家不长进。从生理的角度讲，人体一直在进行新陈代谢，经过一个周期，人体的大部分细胞都会更新，所以人时时刻刻都在变。

\section{不要以貌取人}

我们都知道，不要以第一印象来论断别人，因为人不可貌相。一眼就要看穿别人，实在非常不可靠。但是，别人常常以第一印象来论断我们，往往在第一次见面时，就表现出欣赏或不屑的样子。

人际关系与人的长相并没有必然的联系，不是说人长得漂亮，人际关系就好，人长得丑，人际关系就差。但不可否认，有人靠长相，在处理人际关系时占了很大的便宜，同时也有人因长相吃了很大的亏。举个例子，在《三国演义》中，最出色、最具传奇色彩的恐怕要属诸葛亮了。但不要忘了，当时还有一个与诸葛亮齐名的人，就是庞统，“伏龙、凤雏，得一人可得天下”，可以说二人的实力差不多。可是二人的际遇却大不相同：刘备三顾茅庐，才请得诸葛亮出山，并奉诸葛亮为军师，以师礼待之。而庞统虽然巧献连环计，帮了刘备一个大忙，却需要拿着诸葛亮的荐书去谋个一官半职，最后刘备只派他当个耒阳县令，要不是张飞误打误撞，使庞统展现出过人的才华，估计他一辈子就是“县太爷”的命了。诸葛亮成为蜀国的丞相，帮助刘备三分天下，而庞统可以说寸功未立，就惨死于落凤坡。

为什么实力相当的两个人，际遇如此不同？是庞统浪得虚名吗？不是，我们再看书中对诸葛亮的相貌描写：“孔明身长八尺，面如冠玉，头戴纶巾，身披鹤氅，飘飘然有神仙之概。”再看庞统：“（孙）权见其人浓眉掀鼻，黑面短髯，形容古怪，心中不喜。”“玄德见统貌陋，心中亦不悦。”

这恐怕就是关键所在。就因为庞统长得难看，估计又形成恃才傲物的性格，所以孙权不喜欢他，刘备也不喜欢他。我想?他初见刘备之时，就算拿出诸葛亮的荐书，恐怕刘备也只是碍于情面，封他一官半职，而非真正欣赏他的才华。这就是第一印象的重要性，也是以貌取人的代价。

我们不能以貌取人，但是也不能忽视别人的长相。一表人才也提醒我们仔细端详别人的长相，因为相貌在一定程度上反映了人的内心世界，即“相由心生”。但是，如果完全凭一个人的相貌来判断他的心，就容易犯致命的错误。

人是很善于伪装的，所以要特别小心，看人的时候要多看几遍，要多来往几次。

看一个人，首先要看他的眼睛，眼睛是心灵的窗口：一个人目露凶光，说明他很凶狠；目光暗淡，说明他没有什么能力；双目无神，说明这个人的运气不好。当一个人看到你，目光就闪烁不定，根本不敢跟你对视时，要看他是男人还是女人。一个男人想骗你的时候，他根本不敢看你的眼睛，而是会躲来躲去；女人则刚好相反，她在盯着你的眼睛的时候才开始骗你。这就是男女有别，不要上当。

其次，看他的鼻子。

……

在人际交往中，几乎每个人都带着一个面具。随着人际关系的加强，你要使他主动拿掉面具，达到真诚交往的目的。但是你不能要求每个人一见面就很真诚地对你，这是不太可能的事情。

真诚是一个人重要的美德，但是原本很真诚的人，会变得非常不真诚。很可能是因为他真诚待人，上过几次当以后，就对真诚失去了信心，变得比谁都不真诚了。所以，我们培养人际关系的时候应该记住，要谨慎地保护自己，使自己永远对真诚有信心。只有步步为营的人，才会少吃亏，而少吃亏，才能保持真诚的态度。

其实中国人的信任感是很低的，一开始都不太相信别人，因为历史往往证明，轻信别人肯定吃亏。所以中国人的谨慎小心，不是为了耍别人，而是为了保持自己对别人的信赖感。

\section{提升第一印象}

既然你无法阻止别人以貌取人，那么，在外貌方面你该讲究的还是要讲究。为了防止别人以第一印象论断你，避免受到鄙视或排斥，丧失建立人际关系的大好机会，你必须用心调整自己，表现出“一表人才”的样子，使人一见面就产生比较良好的印象，对建立及开拓人际关系，有相当程度的帮助。

有人认为外在美不重要，内在美才是最重要的，因此理直气壮地忽视外在美。殊不知，如果不先展现自己的外在美，谁会有耐心慢慢发掘你的内在美？

当然，你无需因别人错误的要求，就扭曲或压抑自己的形象，虚伪做作是不会讨人喜欢的；也无需因为完美主义作祟，而压迫自己原有的形象，过分追求外在美，把自己打扮得油头粉面的。在与人相处的过程中，应该尽量表现出真正的自我形象，没有虚饰，也没有伪装。当然，每个人都应该努力使自己的真实形象变得更美好，美国总统林肯说过，人四十岁以后要对自己的相貌负责。就是说，人的相貌是天生的，但是成年之后，我们的经历、我们的体验、我们的情感会在自己的脸上留下痕迹。到四十岁以后，人的先天相貌已不重要，重要的是人在后天所形成的相貌。有的人四十岁时仍然精神焕发，有的人则老态龙钟。至于变得如何，心态很重要。

要想提升自己的形象，我的建议是：第一，要保持身心健康。一个身心健康的人总是容光焕发的，只要健康，你长成什么样已经变得不重要了。你长得再漂亮，脸色灰灰暗暗，目光无神，也不会引起别人的好感。

要想做到“一表人才”，最好的办法是靠内心的修炼，即真心诚意，不胡思乱想，慢慢的，你的相貌就会变得端正。好不好看是另外一回事，最起码要相貌端正，使别人不会对你产生提防心理，不会看到你就害怕，然后会乐意跟你接近，喜欢跟你沟通。能做到这一点，你的人际关系就会发展得比较快。

所以，人每天早上起来，不是看自己长得怎么样，而是看以下几个方面：

鼻梁正不正。鼻梁不正，代表心不正。你没有办法把鼻梁扭过来，只有调整心态，哪一天心态端正了，鼻梁就会正了。所以照镜子是看鼻子的两边对不对称，不对称的话，就说明你的心思歪了。

看自己的气色好不好。气色不好，就说明你的身体有问题，而且给人的感觉也不好。

再看自己的心情好不好。一个人际关系良好的人，一定要有好心情，心情不好看到谁都不顺眼。当你心情不好的时候，要尽量掩饰，但是掩饰不意味着装笑脸，不要给人皮笑肉不笑的感觉，只要不表现出愤怒、悲伤等情绪就好。

第二，提高自身的学识修养。

一表人才，主要来自一个人的学识修养。先天相貌由父母的遗传来决定，后天相貌可以经由自己的学习，透过“相随心转”的运作，来加以改变。学识丰富，内心充实，行为端庄，加上仪容整齐，不就是一表人才了吗?心一改变，外表也跟着改变，不妨试试看。

\section{正确认识自己}

每一个人长得都不一样，即使是双胞胎，也有微小的区别，这叫做个别差异。可是你想过没有，为什么老天要费那么多心思让我们每一个人都长得不一样？这表明上天不喜欢我们一模一样。而今天的教育要把每个人都培养成一样的人，是违反天理的。社会需要各种各样的人，“不拘一格降人才”，“天生我材必有用”，是“用”在不同的地方。所以，我们即使长得不好看，也不必介怀，生下来是什么样子，就要接受这个样子。明白这一点，就比较容易了解自己。一个人照镜子不是为了看自己美不美，而是看自己到底是什么样子，想想自己是来干什么的。有的人永远都不清楚自己这辈子到底要做什么。

随着时间的变化，每个人给别人的印象也会发生变化。人大概可以归纳成四种：第一种，第一眼看上去会觉得他很了不起，然而越看越觉得他没有什么；第二种，第一眼看上去觉得他很了不起，然后越看越觉得他了不起；第三种，一开始就觉得他不起眼，事后证明他的确很不起眼；第四种，一开始觉得他没什么，越看越觉得他了不起。通常是第四种人将来成就最大，因为他不会引人注目，不会遭受别人的打击。凡是长得很漂亮的人，都会引起别人的戒心。所以很多中国人聪明却不喜欢外露，就是这个道理。

每当我们照镜子的时候，都要反省一下：自己现在做的事正确吗？

我们在安排自己人生的时候，要记住适可而止。不要奢望达到人生的顶点，这是错误的观念，爬过山的人都知道，一旦爬到了顶点，就开始走下坡路。没有达到人生的顶点，你是很快乐的；一旦达到人生的顶点，你就会苦恼不堪，高处不胜寒，但这是你自找的。人最怕的就是盲目地要往上爬，真是没有必要，人应该不断地充实自己，而不是盲目地向上升迁。在处理人际关系的时候也要记住这个观点：你不断地充实自己，这样很多人都喜欢跟你在一起；一旦你不断地要往上升，你必定要排挤别人，必定要踩在别人头上，别人看到你就有危机感。

人要有上进心，但并非要不断地升迁。找到你最快乐的位置，找到你最喜欢做的事情，才有意义。如果你有这样的人生观，你的人际关系自然会改善，你的“一表人才”会自动显现。不管你长得怎么样，很多人都会喜欢你。很多人是由不了解到慢慢了解，由了解到彼此相交很深；也有很多人是因不了解而彼此交往很深，因了解而分开的。


\chapter{两套西装}
两套西装与一表人才说的都是外表，但是二者的重点不同。一表人才指的是人天生就具有的东西，比如，你的头发长成什么样，你的额头高不高，你的鼻子直不直，你的嘴巴大不大……而两套西装是指后天附加的东西，比如，你戴什么帽子，戴什么眼镜，穿什么衣服，打什么领带……这些都是两套西装的范围。

一表人才与两套西装是人际关系的起点。你的人际关系好不好，首先看你的本钱。你是不是一表人才，是不是会打扮自己，这就是本钱。

\section{两套西装与一表人才}

两套西装是指外在的、后天的东西。在这里，两套西装也有三层含义：第一，就是佛要金装，人要衣装。穿衣服有两个作用，一方面是保暖，另一方面是尊重别人。什么场合穿什么衣服，不能由你的兴趣决定，因为人是要合群的，要合群就必须要彼此尊重。比如，你要去游泳池，就要换上泳衣，倘若你穿戴整齐地走进去，会被当成怪物；一旦你离开游泳池，就不能再穿泳衣，哪怕是短暂的离开也不应该。你没有换好衣服就是你的不对，就是看不起别人，就是给别人难堪。我们一定要有这样的观念。

文学家但丁有一回受邀参加国宴，故意穿得破旧邋遢，丝毫不起眼，一进宫门，立刻被安排在宴会厅的偏僻角落，没有被盛情款待。等他再接受国王邀请时，改以华服美饰装扮出现，结果立刻被邀至国王身旁的贵宾席位。

宴会进行中，但丁把美酒倒在衣服上，把美食往身上涂，国王以及宾客们都看傻了眼。这时，但丁说，这次所以被礼遇是因为这套衣服，可见被邀请的是衣服，而不是他本人。

看完这个故事，人们可能有不同的感想，如果你认为国王的仆人势利眼，对华冠丽服的人另眼看待，也没错。可是按照中国人反求诸己的想法，其实错在但丁。他不修边幅，被安排在偏僻角落，无可厚非。他故意如此，反过来责备别人，有故意使人难堪之嫌。

第二，什么身份穿什么衣服。有的人经常被别人当做不三不四的人，那也没办法，谁叫他打扮得不三不四的。不合乎身份的打扮，对你是非常不利的。我曾亲耳听见一个老板说：“某人各方面的才能都很好，就是他老穿那套衣服，我实在没有办法提升他。”因为衣服丢了前程，太不应该。不要讽刺别人“只认衣衫不认人”，其实错在自己。

不止是衣服，你的发型，你的眼镜……你身上所有一切外加的东西，都要符合场合。

第三，穿着要合禁忌。就是说，你的打扮不要犯了忌讳。在中国，喜庆穿红，丧事穿素，如果你硬要反过来，就是存心找别扭。

\section{人的外表应符合社会文化背景}

西方人对别人充满了好奇，所以他们经常改变自己，以吸引别人的注意。中国人追求长远，希望慢慢地互相了解。你该变的部分变化了，别人都能接受；你不该变的部分总是在变的话，别人就会接受不了。当你变到让人家对你失去信心的时候，你就被视为不定时的炸弹，别人会对你敬而远之。

一个不会给别人难堪的人，人际关系会比较好。如果一个人平常都很支持你，突然间会当面给你难堪，你以后就会疏远他。西方人可以当面说：“我不赞成你的意见。”中国人一般不会这样做。

文化不同，人们的表现也不同。假如一个老板大声斥责员工：“你去死！”员工该怎么做？如果员工同样怒不可遏，对老板说：“你去死！”试问，这样的员工谁敢用？当然，员工也不可能真的去死。他应该说：“好的，我这就去死。”然后，该干什么还干什么。中国人犯不着当面去反驳别人，但背后可以不照他的意思去做。西方人说“我去死”是真的要去死，而中国人说“我去死”，不去死又怎样？我们有太多的自由，干吗学西方人呢？

凡是后天学习的东西，千万要符合文化背景。一个西方人说“我要退休了”，就会真的退休。一个中国人经常说“我要退休了”，结果十年都没有退休。西方人说了不做，就会失去信用，受到别人的质疑；中国人说：“我一直都想退休，领导不让我退休，我有什么办法？”照样生活得很好，别人也会见怪不怪。就单纯拿穿衣服来说，中国人结婚时新娘子喜欢穿红衣服，越红越喜气；西方人结婚时新娘子会穿白色婚纱，象征圣洁。文化不同，习俗也不同。

人际关系离不开社会的文化背景。有些事情在西方国家可以做，但在中国还是少做为妙。不过有很多事情，我们都接受了错误的信息。没有去过外国的人，都是通过电影、电视了解外国人的，然而，电影、电视所传达的信息往往有误差。

比如，在电影上看到的美国人在家里都穿拖鞋，据我所知，美国人从来不穿拖鞋。就算家里铺着雪白的地毯，客人照样穿着皮鞋走来走去。主人没有权利叫客人换鞋，客人也不需要?。

美国人家里干干净净，如果家庭主妇们需要到外面取信、晾衣服，也都是光着脚出去的，根本不用穿鞋。美国人的观念是外面和里面一样干净，为什么要穿鞋？所以不要相信道听途说的东西。

每一个人都应该有基本的标准，不能乱来，规规矩矩的，人际关系才会好。比如，很多人吃自助餐，在取食物的时候，还聊着天，把口水都喷到食物上了，这是吃自助餐的大忌。还有的人，看见前面的人正在夹菜，等不及，就用自己用过的筷子去夹。

像这些外在的表现，我都归入两套西装的范畴。也就是说，凡是外加的东西，都要符合你的身份，符合场所，符合文化背景。每一个人要选择自己合适的东西，不要盲目地跟潮流。

所以，建立人际关系的时候，要特别小心。一个人如果给人不伦不类的感觉，就会降低人们的信任度。

一个人穿惯了休闲装以后，穿正装会觉得不舒服。人一旦习惯了放松的状态后，再想绷紧会痛苦不堪。所以，不要让自己的身体太舒适，一旦懒散就很难严谨起来，要去磨炼它。在精神方面，人越舒适越自由越好，但肉体方面要多加磨炼，要让自己吃一点苦，才能走更长远的路。

所以，只要你站得住，就不要坐下；只要你坐得住，就不要躺下。人活着就是要承受磨炼，否则就无法形成必要的抗压性。现在，人的压力越来越大，如果没有抗压性，很容易走到崩溃的边缘。

一表人才与两套西装，都是告诉你，要随时准备好与别人见面。你必须接受自己的面貌，哪怕天生有缺陷也要坦然接受。我曾经遇到一个盲人，令我很佩服，他说：“正常人只有两只眼睛，而我全身都是眼睛。”

天生我材必有用，不管你缺乏什么，都可以想办法补全，这才是自信的表现。不要把信心建立在成功上面，总想着“我一定要成功”并没有什么用。“我既然活下来就不会饿死”，这才是我们的信心。“只要我动脑筋我一定有办法”，“只要我喜欢自己，别人一定会喜欢我”，要有这样的信念，然后再去调理自己，实现有形和无形的变化。这样，你很快就能建立令自己满意的人际关系。


\chapter{三杯酒量}
三杯酒量并非只是指喝酒而已，抽烟、喝茶等能吃（喝）进自己的肚子里的东西，都属于三杯酒量的范围。人与人之间的互动需要一些社交活动，喝酒、品茶、喝咖啡、吃饭等并非只是吃吃喝喝，填饱肚子，而是必要的社交方式，能传达很多意思。这些吃吃喝喝的活动，是增进彼此关系的手段，完全不予理会，并非良策。适时、适地、适质也适量地参与，彼此愉快，而且能确保安全。

\section{爱好并非是坏事}

当一个人不断地充实自己，使得自己具备很好的条件时，别人看到你就会喜欢，就会想办法来跟你互动，你就是一表人才；加上你的衣着打扮恰如其分，你就已经在人际关系方面站到了非常好的起点上。可是，你有这样的优势，却舍不得跟别人交往，把自己封闭起来，那就叫做独善其身，或者叫做孤芳自赏，人际关系还是没有办法开展。所以一个人条件再好，也要有诚心去交朋友。

三杯酒量就是人际关系开展的媒介。为什么要谈酒？中国人最清楚，在家靠父母，出外靠朋友。朋友的“朋”字，不是两个月亮，而是两块肉的意思，而且这两块肉还是臭的，如果是香的就变成两条狗了（狗肉通常被称为香肉）。因此说，朋友就是臭味相投的人。你喜欢打麻将，我也喜欢打麻将，我们就会很自然地凑在一起；你喜欢钓鱼，我也喜欢钓鱼，我们很快就会志同道合。所以当你要结交一个朋友的时候，最好的办法就是打听他到底喜欢什么，投其所好。

当一个人没有什么名气的时候，你要认识他很容易，他也很乐意认识你。可是一个人有了名誉、地位以后，就很怕陌生人去打搅他。你登门拜访，他不见你；打电话，他通常会挂掉；你托人介绍，他把介绍人也拒之门外。这时，最好的办法是打听到他喜欢什么。比如，你要向一个人推销汽车，打听到他喜欢钓鱼，那你就赶快去练习钓鱼。然后打听他喜欢在哪个地方钓鱼，你也去钓。你千万不要一开始就坐在他旁边，就跟他说你喜欢钓鱼，这样他会马上提高警觉，认为你不怀好意。你只管钓你的鱼，并且假装没有看到他，你越不在乎，他就越不防备。慢慢地，你们就可以很自然地凑在一起，他如果问你是做什么的，千万不要马上把名片拿给他，这样就前功尽弃了。你只有若无其事地钓鱼，他才会觉得你的动机很纯正，对他没有什么不良的企图。通常在这种情况下，他会告诉你自己是干什么的，而你什么也不要说，欲速则不达。

等他非要了解你是干什么的，你才可以告诉他：“我是卖汽车的。”他可能会说：“你怎么不早说？我前两天刚买了一辆，你早说的话我就向你买了。”你应该说：“没有关系，我来这里主要是钓鱼，生意可以另找时间去做。”你这样说，他反而会主动为你提供信息：“我的朋友还要买汽车，我打电话叫他买你的。”

这样一来，你就掌握了主动权，不要总去求别人，要做到让别人来求你。

一个人要懂得布局，也就是说设下圈套，让别人掉进来。不要以为这是奸诈、耍阴谋，任何事情，你往好处想它就是好的，你往坏处想它就是坏的。两个人有共同的乐趣，才容易情投意合，进而志同道合。可是要了解一个人很不容易，因为中国人有一句话，叫做“防人之心不可无”，人心隔肚皮，不能不防。在西方，你吃亏上当，别人会同情你，因为他们有同理心。在中国，没有人同情你，我们没有同理心，如果有的话，那就容易相处了，大家遵守同样的道理就行了。可是跟中国人讲道理，非常困难，因为我们只相信自己的道理，从来不相信别人的道理，我们很固执，经常自以为是。但这不是缺点，我从不认为中国人有什么缺点，一个人的缺点就是他的优点，一个人的优点正好是他的缺点。要想建立人际关系，就不要有缺点、优点的观念。

一个人喜欢喝酒好不好？没有什么好不好。人与人之间隔着一层假面具，很难彼此诚心诚意地沟通。要使一个人对你不设防，有两个办法：其一是请他去泡温泉，这时大家都赤裸相对，会暂时摘下假面具。其二是请他去喝酒，酒一喝多，就会把持不住。

\section{凡事要适可而止}

据说，我国各地出土了大量的酒器，有些酒器的年代非常久远，这说明中国人很早就懂得喝酒。为什么要喝酒？为了酒后吐真言。交朋友很容易，要变成知心朋友就非常困难。人生最快乐的，莫过于有几个知心的朋友。如何彻底了解一个人？最好的方法是把他灌醉，看他醉后是什么样子，你就可以判断这个朋友可交不可交。

喝酒的本意就在于此，但事实经常是你把对方灌醉了，却不知道他是什么样子，因为你也喝醉了。会喝酒的人能让别人醉，而自己保持不醉。

酒这个东西，很容易使人茫茫然，控制不了自己，然后很自然地流露出自己的本性。我五六岁的时候，我的祖父就让我喝酒。后来我就问我祖父：“我那么小，你怎么就让我喝酒？”?说：“我得知道你醉了以后是什么样子，才有办法教导你。”

人就怕醉了以后乱讲话。因为我们每一个人都有很多秘密，坦白地说，我们是靠合理的秘密在过日子的，人与人之间几乎不可能完全没有秘密。因为人们的立场不同，对一件事情的看法就不一样，所以关系很难协调。适当地隐瞒一些，可以维持正常的关系。

这种隐瞒，如果按西方人的标准，就叫欺骗；如果按中国人的标准，就算不上欺骗，因为我们的出发点只是不想让对方生气。中国人是很诚实的，绝不欺骗，只是没有讲实话而已。这个标准在西方是不成立的，西方人认为，你没有讲实话就是欺骗。

吸烟只有坏处，没有好处。但是酒少喝有益，多喝有害。所以我们只说“三杯酒量”，没有说“两盒香烟”。全世界的人都知道吸烟有害健康，但是还有大量的人在吸烟，可见吸烟也有它的道理，因为社交场合所有人都吸烟，你不吸的话，就融不进他们的圈子。但为健康考虑，还是不吸为妙。

记住，人是主宰，烟酒只是工具，你不能让它们反过来控制你。所有的社交工具、社交媒介都是为我们所用的，我们一定要抓住主控权，绝不做物质的奴隶。一个人如果变成物质的奴隶，非怎么样不可，是很凄惨的。任何活动只要你不沉迷，不上瘾，都没有多大坏处。

可是千万记住，所有的媒介一定要有正当性。酒和色经常是纠缠在一起的，而色情是高度危险的东西。男女之间，一定要坚守最后的防线，一旦攻破了你就全完了。但是很多人都陷了进去，从来没有觉悟过。已婚男女之间只能有友情，不可以发生爱情，这是三杯酒量的一个重要原则。

\section{酒桌交际的技巧}

很多时候，人都是在迫不得已的情况下喝酒的。没有一个人，一入席就把杯子翻过来，说：“我们今天不喝酒。”这种人是不受欢迎的。有的人喜欢喝，你凭什么说“今天不喝酒”？或者你说：“我今天不喝酒。”这样做，你就会跟别人格格不入。

我的酒量并不好，但是我从来没有醉过。第一，我从来不说“我不会喝酒”，这种话一说出来，就会打草惊蛇，所有的人都想把你灌醉，因为你最容易醉，不灌你灌谁？这是自找麻烦。第二，我不主动去敬酒。第三，任何人敬我，我都毫不推脱，然后很快拿起杯子，猛喝一口，结果一点也没喝进去，我做了一点手脚。

在喝酒的场合是不能讲道理的，当然你最后还有一招，就是把医生给你开的证明拿出来摆在桌上，说“再喝就会死”。但是这样的话，所有人都会很快跟你疏远了，你就失去很多交往的机会。

所以三杯酒量要求，你要有起码的社交活动的本事：乒乓球打得不好，但是最起码要能应付；桥牌打得不好，最起码知道怎么跟别人配合；麻将打得不好，最起码不要讨厌打麻将的人。

做人要坚持原则，但态度要随和，这是孔子所讲的“和而不同”，不要标榜自己跟别人不一样。要随和，但是也要坚持原则：要喝酒就喝，但绝不喝醉；要打麻将就打，但绝不赌博；要钓鱼就钓，但绝不沉迷。人永远要做自己的主宰，不要变成任何人或者任何神的奴隶。

总之，三杯酒量要求我们，不要以自我为中心，否则你一个朋友都交不到；也不能沉迷，否则后果堪忧。


\chapter{四圈麻将}
这里的四圈麻将与上面的三杯酒量一样范围很广，打球、跳舞等不会吃（喝）进肚子里的东西都属于此范畴。很多人对打麻将非常不屑，其实如果把它当做一种社交活动的话，无可厚非，但沉迷于赌博就不好了。

三杯酒量与四圈麻将都是人际关系的媒介。也就是说，要想建立人际关系，就要投入这些有形无形的活动中，这样才能同他人加强交往，才会维持人际关系。

打麻将不一定完全有害无利，不必对麻将深恶痛绝，但要适可而止，不影响健康，也不妨害正常生活，才能宾主尽欢。

我们当然不会鼓励别人都去打麻将，但是我们也不鼓励大家对打麻将深恶痛绝。因为世间一切事务，离不开“自作自受”的因果法则，一人做事，最后必须由自己承担所有的后果。既然跑不掉，打麻将的人自己应当会替自己想才对!

\section{赢到后来总是输}

为什么要讲“四圈”麻将？因为上台容易下台难。上台时，声明只打四圈，输赢不论，可是等到四圈打完了，那个输家绝对不肯放你走。所以开始很愉快，最后都会吵架。为什么会这样？原因只有一个，大家都想赢。为了避免吵架，最好的办法是，先问你自己要输多少，而不是问要赢多少。因为打麻将最主要的目的是促进友谊，如果你输200块钱对你影响不大，那输掉200块，却赢了一些朋友不是很好吗？如果总想着要赢200块，结果所有朋友都跑光了，岂不是得不偿失？如果你常常输钱，别人就喜欢和你打麻将；如果你每次都赢，那渐渐地就不会有人跟你打了。千万记住，无关紧要的事情宁可输不要赢。

假如你很喜欢下象棋，你的老板也喜欢下象棋，有一天，他把你叫过去杀几盘，这对你来说是千载难逢的好机会，因为你从来没有这么近距离地与老板接触过。但是，你却连赢他三盘，以后他再也不会请你下棋了，除非他是白痴。天底下最笨的人，就是跟老板下棋居然连赢他三盘的人。但是，如果你连输他三盘，他以后也不会叫你来下棋了。你赢，他不高兴；你输，他觉得没有兴趣。那有赢有输岂不好？非也，你只要赢他一盘，他就不高兴了。那如何讨他欢心？只有一个办法，让他感觉到跟你下象棋很吃力，而到紧要关头他却能赢你，他就会兴致勃勃，每次都会找你。明着强调下棋是不能让的，实际上你不让他才怪。

打麻将也是一样，你去跟老板打麻将，拼命赢他，他不高兴；你拼命输，他会怀疑你有什么企图。记住，中国人是有高度警觉性的，有任何风吹草动都会想半天。

下棋也好，打麻将也好，该赢的时候才赢，该输的时候一定要输，这不是靠技术，而是看情况。你无法说这是游戏，因为这根本不是游戏，而是促进或者败坏人际关系的手段。

\section{经营人际关系}

世界上一切的事物都不是静止的，人际关系也是一样：不是越来越亲密，就是越来越疏远；不是越来越好，就是越来越坏；不是越来越信赖，就是越来越猜疑。总之，始终是动态的。我们常说维持现状，那是自欺欺人，凡事都不可能维持现状。所以人际关系是要经营的，经营人际关系与经营事业一样，只要一段时间不经营，就会变冷淡了。

尤其是当今社会，变动非常快，比如电话，一阵子不联系的话，就可能找不到人了，所以我建议人们要做好几件事情：

第一，每年把电话簿更新一次。最好在元旦的时候，一是比较有空，可以把电话簿拿出来，按照上面的号码一个一个地打；二是可以有个借口，说：“我知道你平时很忙，不敢打搅你，现在是元旦假期，问候你一下。”碰到打不通的电话，就赶快把电话号码划掉。电话簿每年要更新，它才有用。

第二，要把电话号码按地域分区。这样做，是为了你有机会去某地时，可以给当地的各位朋友都打个电话。不然的话，他们会责怪你：“来到我的地方都不和我打个招呼，你眼里还有我这个朋友吗？”我在深圳有几个朋友，每次到了深圳我不去拜访他们，因为怕打搅他们，但是总会给他们一一打个电话，打不通算你幸运，这样既打了招呼，又不会给彼此添麻烦。朋友要经常联系，但不要互相打搅。这样可以加强对方对你的好印象，使对方对你越来越有好感。

胡雪岩出身穷苦，但他的事业越做越大，主要原因是他的人际关系非常好。他有一个诀窍：“走官场先拜宝眷，会同行先说油水。”

到生疏的地方创立事业，最要紧的是建立人际关系，有人好办事。找人，找什么人呢?找官场的人，可以有个靠山；找同行的人，能够打听行情。官场、同行都有人协助，对事业的建立和发展，必定有很好的助益。

找做官的人之前先去找他的眷属，不要先去找他本人。因为做官的人最怕惹事，对陌生人的警觉性很高，生怕会看错人使自己惹上麻烦。而且碍于面子，不好意思直截了当地谈论涉及金钱的事情。他们打惯了官腔，摆惯了官架子，说惯了冠冕堂皇的官话，一下子调整不过来。倒不如想办法找官夫人，送礼、谈回扣都比较方便。出了状况，做官的可以用一句话就把责任推掉：“怎么有这回事，我完全不知道呀!我家里的人从来不插手公务，不可能有这种事情。”从官夫人处打通关系，再见官谈事情，通常很有效。

会见同行就要开门见山地说明给对方多少好处。大家都是内行人，用不着拐弯抹角，不妨开门见山，把如何分享油水说清楚，免得互相猜疑，反而增加沟通的障碍。油水分享得宜，其他就比较好谈。见面不说油水，对方以为不付费白请教，当然不肯供应情报；说不定还会提供一些不正确的信息，混淆视听。该讲的话一句不可少，不该讲的话一句不可多，这就是人际关系拿捏得恰到好处的表现。

有一次左宗棠问胡雪岩：“你一个小小商人，出手这么大方，献给我这么多米，为什么？”胡雪岩心里好笑：你这样问，叫我怎么回答？我当然不会告诉你，这些米本来是要卖给太平军的，因为我要赚钱，不小心被你们劫到，我只能说是献给官兵的。胡雪岩很聪明，他回答：“国家兴亡，匹夫有责，我只是略尽绵力，作为一个老百姓，本来就应该这样。”这样一说，滴水不漏。

有的人叫别人替他打牌时交代：“赢了算你的，输了算我的。”做人要有这个气魄，如果你说：“你替我打，只许赢，不许输。”那谁还敢替你打？很多人就是因为输不起，才破坏了人际关系。能不赢最好不要赢，这样，你的人际关系才会好；如果时时刻刻想赢，最后一定会输得精光。

一个人要多多参与社交活动，但要把握分寸。一个人会唱歌好不好？有一个人，本来很得老板的赏识，但是公司连续几次提升干部都没有他的份儿，原因很简单：每次公司组织唱卡拉OK，因为老板歌唱得不好，所以每次老板唱歌他都助唱。他以为这是帮老板壮声色，殊不知，本来老板唱得不好，一般人还没有感觉出来，他这一帮忙，人们立即听出高下来了。老板心里十分害怕出丑，偏偏他不识相。所以我觉得，很多事情都不能怪别人，千万要回头想想自己有什么做错的地方，天下并没有怀才不遇这回事。

当领导的一定要记住，平常不要主动去拉别人参加什么活动。我们中国人是以领导为中心的：领导喜欢打高尔夫球，下属马上就去练习高尔夫球；领导喜欢钓鱼，下属就练习钓鱼。因为中国人的眼睛是往上长的，下属很自然会去揣摩领导喜欢什么。所以当领导的人根本不用倡导，只要以身做则，下面的人自然就会响应。但是，所有的人际关系媒介都应正常化、正当化，凡是跟色情赌博有关的，千万不要沾染。

有时候，领导也会拿这一点来测试下属，做下属的一定要提高警觉。如果领导问：“这附近有没有什么消遣的地方？”有的人积极推荐，要唱歌到哪里去，要跳舞到哪里去……领导一听就知道，这个人平时不务正业，只想着吃喝玩乐。从另一个角度讲，做业务的人，如果连附近的娱乐场所都不清楚，那他就不太适合做业务，因为他们要经常与客户联络感情，离不开这些场所。但是，做业务的人沉迷此道也不对，跟客户来往，不需要进行不正当的活动，因为做生意，条条大道通罗马。正当性、正常性的活动照样可以联络感情，如果超越了这个限度，最好要有防备，否则有一天会吃大亏。


\chapter{五方交游}
“五方”指的是东南西北中，也就是说，你要结交各方面的朋友，将来有困难时，才会有人来帮助你。人是很复杂的，你认为能救你的人将来可能会害你，你认为与你毫无关系的人，可能会成为你的死敌。所以，人只有广交朋友，才能保证在你需要时，有人会假以援手。最可怕的是只与自己的学业或工作相关的人交往，而不接触其他的人。五方交游，意思是不要自我设限，尽量扩大交友的范围，与三教九流的人都可以放心地交朋友。

\section{多个朋友多条路}

中国人深知“山不转人转”的道理，一方面力求不得罪人，以免冤家路窄；一方面则广结善缘，以便随时、随地可以找到熟人，比较方便办事。同时，结交各行各业的朋友，不但可以扩大见闻，增长知识，而且能够随时请教，不致请问无门。

我一直认为，有人赏识你、提拔你，比你自己努力奋斗更重要。这个观点也许有人不同意，他们宁可选择自己奋斗，我宁愿有人提拔我、赏识我，实在没有办法时，我才会自己去奋斗。但是话讲回来，一个人自己不努力、不奋斗的话，没有人会赏识你，也没有人会提拔你，这就叫互为因果。

一个人只知自己奋斗，而无人赏识，进步会很慢。所以你除了要懂得很多社交的媒介、掌握很多社交原则外，还要多方去尝试。五方交游就是说把你的触角伸得广阔一点，说不定哪一天你就会碰到赏识你的人，中国人通常称这种人为贵人。你若是真能得到一两个贵人相助，很快就能出人头地。

一个人不断提高自己的能力，加上有眼光，有朝一日巧觅有力靠山，彼此互利互惠，就有如神助了。先充实自己，再巧觅靠山，才有发挥的余地。不充实自己，一味等待靠山的提携，是本末倒置，不可能成功。

胡雪岩就是一个例子，他一生有两个贵人，王有龄和左宗棠。当年王有龄穷困潦倒，根本没有人看得起他，但是胡雪岩送给他五百两银子，结果他当了官，又正好派回到杭州。那最受益的就是胡雪岩了。

求人不如求己。一切靠自己，总比样样依赖别人要安全、可靠得多。但是，有人相助，要比单打独斗来得轻松愉快，也是不争的事实。何况当今世界，单打独斗能够成功的机会愈来愈少。有贵人协助，成功的几率必然提高。

胡雪岩先得王有龄的帮助，再获左宗棠的支持，才能够青云直上，快速地开展事业。遇贵人，自然有福气；有实力，贵人的力量才能发挥得出来。

一个人想要赚大钱、立大业，单凭自己努力，终究不如有人提携。好比爬山，靠自己一阶一石向上攀登，哪里比得上有人从上面放下一条绳索，把你拉上去来得轻松、省力而愉快!

“有人好办事”这句话，可以分成上、中、下三个层次来看。对上，有人提拔，使自己扶摇直上。平行，有人依靠，替自己分忧解劳。对下，有人跑腿，把自己吩咐的事情圆满完成。这样，还有什么办不好的事情?

这上、中、下三层次的人都是贵人。空有上面的贵人提携，缺乏左右和下面的辅助、支持，成不了大业。得到下面和左右的支持，找不到上面那一根提拔的绳索，辛苦一辈子，成就了了。

上、中、下都有贵人，自己当然也就成了贵人。但是，贵人毕竟可遇不可求，而如何“遇”，就看你的修为如何了。

不要存心做什么事，否则不会有好结果。记住一句话，有心种花花不开，无心插柳柳成荫。在《易经》中，讲感应的卦，不叫感卦，而叫咸卦。就是说，你要感动别人，一定不要别有用心。存心去感动别人，是感动不了任何人的。

一般人都爱凑热闹，看谁官运亨通就围着谁，殊不知，风水是轮流转的。平常要存好心，照顾那些被冷落的人，这是你最好的机会。要认识一个人，最好是在他默默无闻的时候接近他；帮助一个人，要在他需要的时候伸出援手；赏识一个人，要在他还没有发挥潜力的时候去赏识他。现在一些企业设立奖学金，资助贫苦或成绩优异的学生，就是这个道理。如果企业不注重对人才的培养，专门高薪挖墙脚，那么你挖来的人才可能很快就另谋高就。总之，一句话，人才是要靠自己培养的，而不是去挖现成的。

一个人四面八方都有朋友，将来做什么事情都很方便。但这时候要小心，不要认为太方便就投机取巧。建立人际关系最大的忌讳就是投机取巧。你可以随机应变，但是随机应变过度，就很可能变成投机取巧。一个人最难能可贵的就是会自己节制、自己防备，没有后遗症。

贵人的脸上从来不写“贵人”二字，所以我们分不清到底谁是我们的贵人。偏偏在我们的生活中，你认为是贵人的，事后才发现他其实是害你的人。你认为对你没有好处的人，才是真的贵人。

世事都有两面性，有一句话叫“恩生于?”，就是说，从某种角度讲，所有害你的人都是你的大恩人；所有对你有恩惠的人，其实也在害你。对年轻人来讲，客气不是福气。老板对你要求越苛刻，你获利越多；老板对你越客气，你受害越深。不要以为老板对你太好了，是你上辈子修来的。一个人没有经过磨炼，是无法成功的。就像父母教育子女一样，中国人有句话，叫“棍棒出孝子”，当然，父母不可以乱打孩子，但是一定要有适当的惩罚。正所谓，爱他就是要限制他。一个老板给下属诸多限制，就是希望他们能长进。父母真的爱子女，就要告诉他们这个不可以，那个不可以。

有个人还小的时候，母亲对他极为宠爱，凡事都由着他，他干了坏事母亲也从不指责。有一次，他偷了别人一件东西，回家交给母亲。母亲问他偷东西时是否被人看到，他说没有，母亲就心安理得地收了，毫无责备之意。从此，他的胆子越来越大。真个是“小时偷针，大时偷金”，长大后，他因盗窃罪被判极刑。临刑前，他只有一个要求，要见他母亲一面。看管他的人觉得他很有孝心，就网开一面，让他们母子相见。结果，这个犯人对他母亲说：“我从小喝你的奶水长大，现在我快死了，我只有一个要求，想再吃一口奶水。”母亲信以为真，结果被儿子凶狠地把乳头咬掉了。

这个犯人之所以这样做，是因为他恨母亲！要不是母亲在他小时候一味地纵容娇惯他，他绝对不会成为死刑囚犯。“养不教，父之过；教不严，师之惰”，如果当初教育得当，及时制止他的错误，如今他走的也许就是另外一条光明的人生之路了。

很多事情都要多方面考虑：好心经常做坏事，坏心经常做好事。这是不争的事实。

朋友之间要互相勉励、互相择善，要彼此规劝而不是同流合污。朋友有错，要苦口婆心地进行规劝。有时候，家人反而很难规劝。夫子不择善，朋友要规劝，中国人易子而教，就这个道理。

\section{扩大交友范围}

一个人扩大自己的交友范围，确实存在着风险。因为隔行如隔山，你不可能很清楚别人所处的行业是怎么回事。但是，广交朋友，也可以得到很多好处。我比较喜欢科系多、规模大的综合大学，而不喜欢专业性大学。原因很简单，比如医大的学生，他们在学校五六年时间，听到的、看到的都是和医学有关的。综合大学的学生就不一样了，有学农的、学工的、学商的，等等，大家可以探讨不同门类的东西，彼此能学到很多知识，丰富自己的见闻。

专业只能让你有饭吃，而不能让你过好日子。人要吃饭就不得不有专业，可是有了专业以后，生活的乐趣就大大减少了。一辈子从事一个行当，一辈子只懂这一行，有什么乐趣？

如果你涉猎的知识很多，别人讲“商”你听得懂，讲“工”你听得懂，讲“农”你也很在行，岂不是很愉快？你永远不知道谁是自己的贵人，永远不知道将来你会怎样，所以一个人一定要多方面准备才有安全感。尤其是现在，科技发展日新月异，说不定哪一天，你所从事的工作突然失去存在的必要，那你该怎么办？

一个人在扩大交友范围时要多请教少发表意见。有很多人，虽然有很多朋友，但是他没有长进。就是因为他看到谁都拼命宣讲自己的那一套理论，这样做的话，有再多的朋友也没有用。比如，你好不容易找到一个医生做朋友，经常咨询他医学方面的知识，这样你会受益良多。如果一见面你就拼命讲自己的专业，你将一无所得。交朋友是要增加见闻的，而不是吹嘘自己。你讲得越多，人际关系越坏，因为大家好不容易聚在一起，只听你一个人滔滔不绝，那别人还有什么意思？一个人际关系良好的人会多关心别人，让别人多说话。中国有一个广泛流传的说法，说小孩子有耳朵，有眼睛，没有嘴巴。就是说，小孩子要多听多看，少说话，才能增长智慧。

中国人，第一次请你提意见是客套，不要当真，他说“有意见请说”，其实心里想：“你还要说什么？”有时候领导会对下属说：“不要客气，我们是一家人，有话请说。”下属说了就会对自己不利。

以喝茶为例，西方人请你喝茶你就喝，没什么好客气的。中国人自己要喝茶时，也会说“请喝茶”，其实这根本不是在请你喝，只是给你一点面子，客套一下而已。所以，很多事情都要推让两三次以后，你才会清楚他是真的还是假的。

再举个例子，你到别人家里去，正赶上他们家在吃晚餐，主人一定招呼你一起用餐，这时你应不应该过去？换成外国人，就会直接过去吃饭，而中国人就不行，虽然此时你肚子饿得要命，也要说“我刚吃饱”。如果你像外国人一样真的过去用餐，你会发现桌上菜也没有了，汤也没有了，你这不是自寻烦恼吗？主人会想：“你还真的过来一起吃啊，我只是给你面子，招呼你一下。”所以当中国人第一次招呼你用餐时，你要说“我刚吃饱”；他说：“不要客气，过来吃点。”你要说“我才吃了三碗”。等他第三次招呼你，而你看见饭菜很丰盛，就不必客气了。此时你再吃三碗也没什么，没有人会笨到问你“怎么一顿饭要吃六碗”。

很多时候，中国人有自己特殊的习惯，不能以常理来判断。下面这个小测验可以测试出你的人际关系技巧如何（仅限于男性）：假如你烟瘾发作，你发现烟盒里就剩下一支烟而周围有很多熟人，请问你怎么办？

烟这个东西，你只要伸手递给别人，人家随手就拿走了。但是你躲起来抽也不是办法，你不躲起来没有人会注意你，一躲起来所有人都会看到你。这时，人家会说：“一支烟也不值多少钱，你还躲起来抽，像话吗？”

碰到这种情况，只有一种做法：你要让所有人都知道，“下午才买的烟怎么就剩下一支呢”，即告诉所有人你只有一支烟。然后照样请别人抽烟，当然没有人会接过来，那你就可以堂而皇之地自己抽。我还有个更聪明的做法：首先让别人知道你只剩下一支烟，然后对别人说：“来，我们一人一半。”没有人会要你这半支烟的，这样你既可以抽烟，又能大得人心，广结善缘。

不要以为这种人花样太多。什么是文化？文化就是花样。美国人的一套花样就叫美国文化；日本人的一套花样就叫日本文化；我们中国人的一套花样就叫中国文化。端午节划龙舟就是花样，清明节大家去扫墓也是花样。


\chapter{六出祁山}
六出祁山讲的是诸葛亮的故事，意在说明，人要明知不可为而为之。凡事难免遭遇困难，若是遇到挫折便心灰意懒，何以成大事?交友过程中，有时会引起误会，如果因此而垂头丧气，如何培植深厚的友谊?必须坚定信心，以愈挫愈勇的精神，排除万难，而且要表现在实际行动上，大家才会相信，拉近彼此的关系。

\section{明知不可为而为之}

现在企业管理都讲求可行性分析，这种做法看似科学，我却不以为然。人又不是神仙，怎么可能预知未来？现实中有很多歪打正着的事，如果不去试的话，怎么可能断言不会成功？科学只能接近真理，永远无法达到真理。

诸葛亮为什么要六出祁山？像他这么聪明的人，掐指一算就知道后果如何，他为什么还一次次地不辞辛苦地带兵打仗？为的就是让后代人明白，他真的是为蜀国鞠躬尽瘁，死而后已。

西方人重利害，凡事计算来计算去，无利可图就会放弃。中国人重势利，无利可图也会做。五方交游与六出祁山讲的是人际关系中的相互交往，一方面要扩大范围，一方面要拓展深度。就是说，好不容易交到一个朋友，要深入地去交往，这就是六出祁山的含义。

祁山是个地名，在三国时期是个非常重要的地方。当时，魏蜀吴三分天下，蜀国要想攻打魏国，一定要经过祁山。看过《三国演义》的人都知道，如果没有诸葛亮，或者诸葛亮不出山，曹操会很快统一全国。可以说，曹操最大的克星就是诸葛亮。一个人的力量可以改变天下的局势，这是事实。但是话说回来，诸葛亮出山又怎样？用水镜先生司马徽的话说，诸葛亮“得其主而不得其时”，难道诸葛亮不明白这一点吗？他当然清楚，但是既然刘备这么有诚意，只好努力去做。明知不可为而为之，这是诸葛亮给我们的启示：一个人不要过于重视结果。

今天很多人都奉行结果论，认为“成则王侯败则寇”。但是中国有句老话——不以成败论英雄，结果不是人力所能控制的，人能控制的只是过程。所谓“谋事在人，成事在天”，你一定要认真谋划，要全心全意投入，至于结果怎么样，由上天来决定。同理，我们能把握的只有自己，至于跟别人互动的部分，永远没有把握。很多时候、很多事情都是听天由命，这不是宿命论，也不是消极思想在作怪，我只是强调，既然你可以控制过程，就要全力以赴，至于结果，不要考虑得太多，不要因为结果可能不理想，就束手束脚。凡事“尽人事，听天命”就可以了。

\section{朋友之间可以通财}

通财只是朋友互动的一种形式，不要只局限于字面意思，朋友之间的各种互助形式都在通财的范围内。

我们一定要想办法深入地跟一些人交往，否则的话，就是“朋友满天下，知心无一人”，那你交那么多朋友有什么意义呢？一般的朋友叫做点头之交、泛泛之交，在紧要关头是帮不上忙的。我们要有三五个知心朋友，当你危急的时候，他们真的会仗义相助。

很多人平常不在乎有没有朋友，在危急的时候才意识到，朋友是很宝贵的。可是等到危急的时候再去交朋友已经来不及了。临渴掘井是愚蠢的行为，我们一定要未雨绸缪。宁可一辈子不用他们，也不要用的时候找不着人。假如有一天，你急需钱用，你会找谁？你只要向三个人开口却无功而返的话，你就会心灰意冷了。

当你不需要钱的时候，要试着跟你的朋友去借钱，看看能不能借到。一旦真的需要钱的时候，才知道找谁借才不会落空。如果你有一大群朋友，就要自己去规划：这几个朋友，我跟他开口借钱没有用；这几个朋友，我根本开不了口；这几个朋友上次跟我借钱，我没借给他们，如果再跟他们开口岂不是笑话吗？去掉完全不可能的，然后看看还剩下哪几个。这时，你可以拿起电话跟他们借钱，看他们如何回答。你一开口他就说不行，那这个人你是借不到钱的。你一开口，他满口答应，这个人你也是借不到钱的。中国人在“没有问题”后面常常加上四个字“从此不提”。如果对方说：“你是要现金还是要支票？”那这个人就会借你钱，不然他不会多此一问，这样的朋友就值得交。但事情还没有结束，等对方把钱送来以后，你要不动声色地拖几天，看对方是什么反应。如果对方第二天就急着问你什么时候还钱，那这个朋友还是很麻烦的。如果过几天后他还若无其事，你就可以把钱还他，并致谢。以后，你真的需要用钱的时候，就可以找他。中国人是救急不救穷，穷人没有办法救，因为那是无底洞。连“急”都不救的话，那要这种朋友干什么？平常要跟你的朋友多进行互动，你才知道他哪些方面是值得你信赖的，哪些方面是他力所不能及的。

在人际交往中，千万不要留下遗憾。人生最要紧的，就是做到自己不后悔，如此而已。好坏是一回事，成败是一回事，只求不后悔。老实讲有时候失败的人，反而比成功的人日子更好过。我相信各?也听过这样的故事：

楼上楼下有两家人，楼上住着大富翁，整天笑也笑不出来。楼下住的是穷人，他整天唱着歌，快乐得不得了。楼上的那个富翁就很感慨：为什么人家那么穷，还整天开心，我这么有钱，却整天苦恼？他的朋友告诉他：“你想要快乐吗？那你就把钱送给楼下的穷人好了。”富翁真的这样做了。从此楼下的人整天愁眉苦脸，不知道拿这些钱怎么办，而富翁却开始笑了。

钱有很多好处，但钱不能解决所有问题。任何事情都有好坏两个方面，所以，人无论相信什么，只相信到差不多就可以了。“差不多”真正的意思是不能差太多，不能差太多就是刚刚好。凡事只要差不多就好。


\chapter{七术打马}
人际关系其实是自己跟自己的关系，不完全是自己跟别人的关系。一个人先要接受自己，别人才会接纳你。一表人才，应该由内而外。我们现在过分地重视外在，而忽略了内在的修养，这是错误的方向。内在的东西才是主导，外在的一切都是根据内在来改变的。所以当你觉得对外面不满意的时候，应该从里面去修正。

给别人的第一印象好还不够，还要做好各种准备，这就是两套西装所讲的。如果你要去参加会议，事先打听一下，应该穿什么衣服，带什么用具，然后再确定怎么做。和别人一模一样，不好；和别人完全不同，也不好。人与人之间，最好大同小异。完全相同你就不显眼了，但不同到特殊的地步，同样对你很不利。

我们平常还要练习一些社交所需要的活动，一个样样都不会的人，是没法和别人打成一片的。没有办法跟别人互动，你有再好的条件也没有用。

所以一个人要有三杯酒量，要懂得四圈麻将。“三杯”的意思是，吃喝方面要适当地加以控制，过量的话，对自己身体有害，而且没有人会同情你。“四圈”是指积极参与社交活动，但不能过分。社交的媒介是用来促进人际关系的，一旦变成破坏性活动的话，就会适得其反。

另外，还要扩大交友的范围，交各方面的朋友，要多多地向他们学习。

\section{多说恭维话}

如果能做到上面这些，就可以进一步探讨一些人际交往的技巧，这就涉及七术打马的内容。七术打马就是俗话说的拍马屁，为什么不用拍马，用打马？因为中国人是不喜欢马屁精的。总有人认为中国人喜欢拍马屁，拍马屁就能成功，这是绝对错误的。这种说法是用来骗别人的，千万不要用来骗自己。一旦相信拍马屁能成功的话，就会走上歧途。

西方人不反对你骗自己，只反对你骗别人。中国人相反，你骗不骗别人无所谓，千万不要骗自己。世界上没有一个人光靠拍马屁就能够得到提升的。很多管理者都非常讨厌马屁精，因为他们相信，马屁精会害死他们。

虽说这里的七术打马也是这个意思，但重点在于“术”。很简单，要是你有本事，拍马屁拍到好像没有拍一样，让上司、长辈对你的举动留下深刻而良好的印象，遇到机会就会主动提拔你一下。如果你一行动，别人就知道你在拍马屁，那就算了，你连拍马屁的本事都没有，还拍什么？

中国人最讨厌拍马屁的人，但是很喜欢享受“马屁”的味道。所以有本事的人，不去拍马屁，而是制造“马屁”的味道，这就叫七术打马。为什么魏徵那么出名？就是因为几千年只出了这一个，物以稀为贵。如果每个人都犯颜直谏，那魏徵就不会出名了。凡是当面说实话的人，大都下场悲惨，难得有一个魏徵，说了实话却没被砍头，更加说明唐太宗很了不起。

中国人要特别小心自己的耳朵，因为它只听悦耳的话，不愿意听逆耳的话。其实，我们并不了解自己的耳朵。一个人出生以后，耳朵就开始慢慢退化。职位低的时候，什么话都要听，不听就会挨骂；职位慢慢高了以后，就只听自己喜欢听的话；到最后，什么话都不听，唯我独尊，变成孤家寡人一个。

喜欢听好听的话，而不是真实的话，这是人最大的毛病。所以，我们讲话的时候，如果不加上适当的恭维，对方根本听不进去。讲话首先要让对方听得进去，否则就是白费唇舌。对方听不进去，你的话再对、再真实也没有用。

恭维话是不是等于奉承话?如果是的话，岂不是人人都是小人，喜欢被奉承？恭维话不等于奉承话，一般人把这两种话混在一起，分不清楚，以致产生严重的误会，认为吹牛、拍马屁是成功的必要条件。

成功的人，实际上不吹牛也不拍马屁，却装得像吹牛、拍马屁一样。换句话说，即用吹牛、拍马屁的形式，说出有根有据的实话。忠言逆耳，而奉承话又十分危险，容易被对方听成挖苦话。现在把忠言逆耳和奉承话统合起来，说出一番有事实根据，能让对方听得进去，而且听起来很受用的恭维话，才是有效的途径。

把事实的重点说出来，加上适当的形容词，叫做宣传。如果没有事实的依据，只是加上了很多形容词，就成为不实的广告。恭维话和奉承话的不同，大概和宣传与不实广告的区别差不多。

多说恭维话，少说奉承话，说得恰到好处，效果必定良好。

所以说，美国人的沟通方式，在中国根本行不通。因为美国人听话跟中国人听话的目的不一样。美国人听话，是听这句话是真是假，中国人听话，是听对方的立场如何。我们比较在乎对方同自己是不是同一立场，而不在乎你讲真话还是讲假话。美国人只在乎你讲话的内容，不在乎你怎么讲。而在中国人看来，形式大于内容，你讲什么都无所谓，就是不能“这样”讲。

讲话的态度，讲话的形式，在中国人眼里是非常重要的。我们对别人要说恭维话，就是因为我们都比较爱面子。

有些人受到西方文化的影响，提倡同理心，注重礼貌和道歉。其实不然，西方人做错事是要道歉的，在中国，很多人认为道歉没有用。如果别人打你一下，再对你说“对不起”，你会高兴吗？中国人一般是不太接受道歉的。

现在，美国的大型百货公司的售货员都会说日语，我到美国的时候，经常有售货员误以为我是日本人，就用日语招呼我。当我表明自己是中国人的时候，对方就会很诚恳地道歉。但我一点也不感动，我不相信中国人会因为别人的道歉而感动，我只是觉得好笑而已：这有什么好道歉的？如果这事发生在中国，情况会大不一样，中国人如果把别人当成了日本人，会觉得不好意思，但是嘴上绝对不会认错，反而会说：“你怎么长得像日本人呢？”意思是，明明是你长得像日本人，所以我才会跟你说日语。中国人嘴上不会认错，心里会认错，这是我们跟西方人不一样的地方。

不要强求中国人嘴上认错，因为道歉通常是假的。在中国的组织中，通常是职位高的人认错，职位低的人不要认错。假如开检讨会的时候，职位低的人开口就承认自己错了，他就惨了，所有人都会把错误推给他，那他会无缘无故地背很多黑锅。假如是老板承认错误，说：“这件事大家都没有错，是我一个人的错。”那么，所有的干部都会承认是他们做得不好。

如果干部没有反应，那老板会问：“李经理，你听了我的话有什么感想？”如果李经理还是无动于衷的话，老板就会说：“那你真的认为我错了，你没有错？”李经理绝对不敢点头说“是”。中国人的一言一行是高度艺术化的东西，这不是奸诈，只是诚心诚意地玩一些花样，而非假仁假义地玩弄权术，二者有着本质的不同。

哪个人特别爱面子，你就多对他说一些恭维话，让他觉得有面子，他就会心甘情愿地替你做事情，这才是最轻松的方法。建立人际关系的目的，实际上，就是你要从中得到一些好处，这不是丢脸的事情。交朋友就是当你碰到危机的时候好向别人寻求援助，如果紧要关头你的朋友都袖手旁观，那你就知道自己交错朋友了。

要别人帮忙的时候，你也不必求人，因为一开口相求，你自然矮人一等。你只需要会讲恭维话就可以了，千万不能拍马屁，那是死路一条，而讲恭维话才会有机会。

很多人吃亏，就在于说话不好听。下面这个小故事就说明了这一点：

从前有一个国王，他晚上做了一个梦，第二天上朝，他对大臣们说：“我昨天晚上做了一个梦，你们找一个人来圆一圆，这个梦是什么意思？”大臣们很快就找到一个人来圆梦，那个人说：“启禀国王，您这个梦是说，您的朋友一个一个都死掉了，最后只剩下您一个人。”这个国王越听越伤心，就把这个圆梦的人推出去杀掉了。

不久，大臣们又找了一个人来圆梦，那个人心里想，国王的梦预示的内容和第一个人讲的差不多，但是如果实话实说，就会落得像第一个人那样的下场，所以他说：“恭喜国王，贺喜国王，因为您是您所有朋友里面最长寿的。”国王听了非常高兴，赏赐给他大量的财物。

两个人说的意思其实是一样的：国王的朋友一个个都死掉了，那他就是最长寿的。但结果是，第一个人被砍头了，第二个人获得了丰厚的赏赐。那我们为什么不学学第二个人，把话说得好听一点呢？

一个人嘴巴要甜一点，才能受欢迎。心坏没有人知道，嘴巴坏的话，所有人都讨厌你。同样一句话，有正反两种说法，就看你自己如何选择了。

恭维人家也要适度，你把人家捧得太高，连他自己都不相信，那有什么用？说恭维话与拍马屁是不同的，举个例子，当你的主管讲完一句话，你马上说：“这样做最好，这真是明智的决定。”这不是拍马屁是什么？所有人都认定你是在拍马屁，连你的主管都会觉得不自在。但是，当你的主管讲完话，你稍微停一下，说：“这么一来，我所有的问题都解决了……”你这就是在恭维别人，因为你说的是事实。就算要拍马屁，也要拍得不让别人感到肉麻；你一拍马屁，所有人都坐立不安，全身发麻，那不如不拍。

\section{少去讨好人}

有人以为“中国人喜欢被讨好”，似乎只要肯用心去讨好中国人，自然左右逢源，什么事都办得通。其实，中国人最不容易讨好，因为我们的警觉性很高，遇到有人讨好，立即提高警觉：“他为什么对我这么好？”因而怀疑“他究竟安的是什么心”，以至“心里好笑”，处处加以防备。

历史上有很多事实，证明喜欢被讨好的人，最后被小人所包围，因而拖累了自己。喜欢讨好别人的人，由于讨好所有的人结果等于没有讨好任何人，势必采取押宝的方式，押对了固然可以得势一时，但终究会败下阵来。万一押错了，徒然费尽心机而毫无所得，亦将会悔恨不堪。

喜欢被讨好和喜欢讨好人，都得不到实质利益，中国人明白这番道理，当然不屑为之。

自古以来，没有人立志亲近小人，也没有人愿意成为小人，那又为什么会不断形成“小人当道”的局面？其关键在于，恭维话和奉承话很不容易区别，以致自己受害、社会不安，而国家也难以求治。

举个例子，下属主动替上司分忧解劳，算不算拍马屁？那要看情况：如果自己的工作做得乱七八糟，却经常跑到上司面前问“有什么事情要我帮忙的”，当然是拍马屁；如果分内的工作做得很好，还有多余的时间和精力为上司分忧解劳，谁敢说他拍马屁！相信大部分人都会认为这个人确实不简单，有机会升迁，非他莫属。

问题就出在现代人越来越搞不清楚这种判断标准。当上司的，看见下属热心分忧解劳，不管他本职工作有没有办妥，便认定他是好下属；做下属的，不管自己分内工作有没有做好，就厚着脸皮要为上司分忧解劳：结果都造成小人当道的局面。

不以讨好的方式，不抱讨好的态度，却能够得到他人的欢迎，在他人心目中占据牢固的位置，这才是人际关系的精髓所在。


\chapter{八口吹牛}
拍马屁的对象是别人，吹牛主要是吹自己。吹牛也包括很多意思：凡是自己表现得与众不同的地方，就将其放大，这就是吹牛；凡是把自己碰到的倒霉事都推给别人的，也是一种吹牛的方式。一个人要适当地奉承一下别人，适当地吹嘘一下自己，但一定要适可而止，不要吹破了。最容易吹破的就是“第一人”这三个字，很多人宣称自己是某方面的第一人，结果很快就被戳破了。又没有正式比赛，你凭什么就号称自己是第一人？

\section{要自己肯定自己}

七术打马和八口吹牛是人际关系的技巧，用得合理就能发挥很大功效，但太过分了，只会伤害自己。

一个人应该谦虚，却不可以看轻自己。这两者并不矛盾，可以相辅相成。看得起自己，是对自己有一份期待，是给自己一份激励，也是给自己一些鞭策。谦虚不自满，是对人的一种态度，让人家乐于亲近，更愿意帮忙。

对自己有信心，还要谦虚地不断充实自己，才是真正看得起自己。

\section{让别人来赞扬你}

有些话，只能说别人而不能说自己。比如，说别人“你是最棒的”，没关系；说“我是最棒的”，那你就完了。我们恭维别人可以，千万不要捧自己，最好是让别人来捧你。真正会吹牛的人，都是通过别人来吹，而不是用自己的嘴巴来吹。

有个老板，把所有干部召集起来开会，说：“我比你们辛苦，知道吗？你们以为老板好当？你们下了班还可以看电视，还可以陪小孩做作业，我没有时间。你们看电视时我在制订计划，你们陪小孩的时候我还在处理公务……”

可是，这话一说出来，所有的干部都不以为然，心里想：“既然你那么辛苦，干脆不要做了。你那么辛苦干什么？”甚至有人想：“你那么辛苦，那换我当老板好了，我不会像你那么辛苦的。”

聪明的老板要想表扬自己，要通过得力干部的嘴。比如在会上，一个得力的干部站起来说：“各位，我告诉你们，我偶尔有事情去找老板，不管是晚上几点，他都在处理公务，不像我们在家里看电视，照顾小孩，老板比我们辛苦多了。”而老板还要谦虚地说：“那只是偶尔一两次而已。”这样，所有人都说这个老板了不起。

红花要有绿叶配，自吹自擂毫无作用。有些话，别人说可以，自己打死了也不能说，说了只有反效果。朋友时常会起到代言人的作用，就像在家庭里面，爸爸妈妈常常互为代言人一样。如果爸爸对儿子说：“我这么辛苦都是为了你。”儿子听了一般不会感动，而妈妈说的反而有效：“你看我们可以轻松地聊天，而爸爸还在忙，爸爸这么辛苦是为谁啊？为我们。”爸爸对孩子说：“你看，妈妈整天在家，好像妈妈没有什么能力，实际上，妈妈处理家事比我办公还辛苦。”这样，小孩子会很感动。如果妈妈说：“我做家事很辛苦。”小孩子会想：又开始唠叨了，真烦。

一个会吹捧自己的人，往往通过自己的好朋友来达到目的。诸葛亮从来没说过自己料事如神，都是他的朋友帮忙宣传的。一个人要自己吹嘘自己，就说明他连一个朋友都没有。吹牛，一定要预计好程度，过头了就会产生反效果，那就闹笑话了。

如果你到别人家里，主人请你吃饭，你就不要老谈自己的事情，那表示你目中无人。吃饭的时候最起码讲两句恭维话：“哎呀，我本来是来谈事情的，想不到让嫂子多劳了。”这样，女主人就很开心。如果你说：“这个菜太好吃了，在饭店绝对吃不到。”这就有点像挖苦了。有时候，恭维话说得太过分，就会变成挖苦。

现在很多人一看到男的就喊“帅哥”，看到女的就喊“美女”，那就表示中国没有帅哥和美女。本来一个人长得不错，你说她美若天仙，那别人看了，会觉得不过如此。这就是期望太高产生的心理落差。

我们捧别人，要捧到大家都不会感觉到难过；捧自己，要使自己能够被别人所接受。

吹牛要适可而止，吹过头了，“牛皮”就会破掉。不要动不动就给自己加个“最”字，这种“牛皮”最容易戳破。


\chapter{九分努力}
九分努力就是不断地精进，按照字面的意思理解就可以，不用做过多解释。

\section{努力没有用}

有人曾问我，到底是运气重要，还是努力重要？我回答说：“当然运气比较重要。运气一来，挡都挡不住，否则你再怎么努力也没有用。可是话又说回来，如果你不努力，你怎么知道自己的运气好不好，所以还是努力比较重要。”这话听起来似乎不知所云，其实很简单：努力只是人的本分，努力并不能保证成功，否则人生太简单了。但是，努力还是一个必要的过程，所以我们还是要九分努力。

我每次给企业做培训的时候，都会说：“努力没有用。”当我说这句话时，最开心的就是企业的老板。企业老板认为，就是要强调努力没有用——员工都很努力，努力浪费时间，努力制造问题，努力打小报告……有什么用呢？

一个人要用心，而不是努力。什么是用心？每个人的感觉都不一样，没有办法一概而论，所以我只能强调努力。现在很流行一句话，叫做“年轻不留白”，我认为这句话害人不浅。

一个人不留白，就是整天忙碌，没有时间思索，学习知识也不消化。所以，人一定要留白。每天最少留20分钟给自己，什么都不做，只想想自己今天的得失，这样你才会有收获。现在，很多人一有时间就玩电脑，电脑是死的，人是活的，用活人去跟死的东西拼，那是愚蠢的做法。机器是让我们使用的，不是让我们跟它拼命的。每个人都要留白，这样才会开悟，智慧才会增长。留一些空白，你才有想象的空间，才能发挥智慧，才能开发潜力。

\section{运气很重要}

努力只是对自己的一种交代，不要期望它一定会产生怎样的效果。有很多人否定运气，因为他们不懂什么叫运气。人生下来就有一口气，气要去运，不运等于零，你怎么运你的气，就是你的运气。人生有很多苦恼，不要怕挫折，你要接受它们的磨炼，即好好地运你的气。

那运气怎么能运得好？很多人运气好，只有一个原因，就是祖上有德。什么叫做德？中国人很重视德，德其实就是“得”，就是说，跟你在一起的人，都得到一些好处，你就是有德。跟你在一起的人，都损失了一点，你就是缺德。千万记住，有钱没钱是一回事，当不当官是一回事，发不发财又是一回事，最根本的就是别人跟你在一起时，是有所得还是有所失。所以说，当领导的，最要紧的就是做到让你的下属，第一能赚到钱，第二能学到东西，第三能不断长进，这样，你就是有德。否则，你就是缺德。

别人为什么对你好？就是因为你祖上有德。中国人是崇拜祖先的，认为不努力就对不起祖宗。中国的家庭里小孩不听话时，家长会说：“你这样不听话，我倒没什么，不过你想想看，如果你爷爷泉下有知，会怎么想？”小孩会想，我这样做，我爷爷要是知道肯定会很伤心，下次他就不这样了。


\chapter{十分忍耐}
九分努力之后，还要十分忍耐。什么叫做忍？“忍”是心上插着一把刀。想想看，一把刀插在心上，你还若无其事，你就很会忍。一个人要想得到别人长期的拥戴，非学会忍不可。因为人们之间经常有暂时性的误会，暂时性的不了解，凡是忍耐过去的，后面才有好日子过。

\section{会忍才能赢}

成功者永远是忍耐力强的人。你的基础越好，表现越好，各方面给你的打击就越大。我觉得中国人是重道义的，你表现不好，别人就没有必要打击你。当很多人都打击你的时候，就说明你很了不起。有人把你当对手你才有价值，没有人把你当对手，就是根本不把你当一回事，这是最糟糕的。

为什么企业竞争到最后都演变成价格战？价格战的结果往往是同归于尽。你卖8块，我卖7块……结果企业利润尽失，只能走向灭亡。那怎样避免价格战？就是要忍耐。如果你的报价别人嫌贵也没办法，要坚持住，不能随便降价，否则就会陷入恶性循环。

某生产企业生产的零件一般市场价格在120元左右，由于市场竞争激烈，同类企业纷纷采取降价措施，以吸引买家，结果价格越来越低，利润越来越少。这家企业不想陷入降价的陷阱，一直把价格稳定在120元。当然了，买家每次来询价，都会因嫌贵而打退堂鼓，但这家企业坚持不降价。

过了一段时间，企业的销售人员会写一封信给买方，说：“您上次到我们这儿询价，距现在已经有两个星期了，您没有再来找我们，就说明您已经买到比我们更便宜的产品了，所以我们特别写信恭喜你。同时，我们内部也检讨，为什么我们的价格始终下不来？检讨的结果是我们没有办法降价，因为我们生产的产品材质与同类产品不一样，工艺也不一样……如果您买的产品适用，您就继续购买那种产品；如果有需要的话，欢迎您随时回来。”

买方收到信后往往会再次光顾这家企业。

在没有营销学这门学问的时候，生意都很好做；有了营销学以后，生意越来越难做。以前企业管理只是针对生产部门，比如怎样挑选材料，怎么控制流程，怎么确保质量等，对顾客都是以诚相待，没有其他的花样。现在企业在进行市场营销时，往往设计出很多花样，渐渐失去了顾客的信任。我有一次到某地去旅游，购买的景点门票附了一张抽奖券，我一看就知道是骗人的。营销学没有盛行以前，这种抽奖活动都是真的，多少让你得到一点奖品；营销学盛行之后，这种活动完全是噱头。果然不错，有人中了奖，奖品是一张购物的八折券。结果发现，没有使用八折券的人，用五折的钱就可以买到商品，使用八折券的人不但花钱更多，并且还限购商品，基本上是一般人不愿意购买的商品。

今天的顾客对一些企业失去信心，往往就是因为这些企业太热衷于这种“现代营销学”。顾客上过一次当以后，就不会再上当了。顾客并不是傻瓜，如果企业把顾客当傻瓜，那企业也就成了傻瓜。凡是把别人当做傻瓜的，自己才是真正的傻瓜。聪明和笨的差距，只有5分钟而已：聪明的人5分钟以前就想到了，笨的人5分钟以后才想到。玩任何噱头，最后都是自己倒霉，大家逐渐失去对你的信任感，你的路越走越窄，越走越艰难，完全是自作自受。

\section{欲速则不达}

要想忍耐很简单，先要想想你为什么急。今天整个社会有一个毛病，都强调快，说速度就是生命，任何事都要捷足先登，但我不这么认为，单纯求快，只会死得更快。方向绝对比速度优先，方向错了，越快就越倒霉；方向对了，才能放心地加快速度。

另外，急事要缓办，否则就会忙中出错，而且越错越忙，越忙越错。急是没有用的，做任何事情一定要有一个过程，欲速则不达。

当你把一件事情说到绝对的时候，就开始有错误了。因为我们是活在一个相对的世界里面，只要我们活着就没有“绝对”这回事。死了以后是绝对的，出生以前是绝对的，出生后就是相对的。爱因斯坦也只是提出了相对论，没有提出绝对论，因为一个是本，一个是道，本立而道生。每一句话多少都有一点道理，但不能强调，一强调就错。比如说当学生要努力，结果有的学生过分努力，只会毁了自己的身体。一个人不用功，你要告诉他用功一点；一个人很用功，你要让他放松一下。

所以一个人要懂得见人说人话，见鬼说鬼话，才不会错。对于不同的人要讲不同的话，这是高度的智慧。真理是永远说不清楚的，因为它范围太广，每个人只能探究到一部分。因此越聪明的人越需要忍耐，否则会活活被气死。真正高明的策略都是饱受攻击的，因为一般人看不懂。等到一般人都看懂了，已经时过境迁、毫无意义了。所以忍耐是对有才能的人来讲的，对于没有才能的人，忍耐与不忍耐根本没有什么不同。懂得越多的人越需要忍耐，一个懂得多的人，讲话不会铁口直断，因为凡事充满了变数，随时有改变的可能。


\part{沟通的门道}

\chapter{沟通的作用}
人际关系与沟通，彼此影响。二者可以互补，也能够相克。人际关系良好，沟通就比较顺畅；沟通良好，也能够促进人际关系的和谐。反过来说，人际关系不良，就会增加沟通的困难；沟通不良，就会促使人际关系变坏。

不善于沟通的人，最好加强人际关系，来弥补自己的缺失。人际关系不好的人，最好培养沟通的能力，以求改善人际关系。事实上，二者之一获得改善，对二者都有所助益。

什么叫人际关系？简单地说，我们每一个人都有能量，不同的能量相碰撞，就产生磁场，产生磁场后，彼此调整自己的频率，就形成人际关系。

所谓沟通就是对准频率，频率不对，就会听而不闻、视而不见、沟而不通。我们要调整自己的频率，而不是叫别人调整频率。要先观察对方的频率，把自己的调得跟他一样，他就很容易跟你沟通。

由于中国人的人际关系实际上是人伦关系，所以在沟通的时候，需要融合伦理的观念，忌讳没大没小，否则就会破坏人际关系。比如，打招呼这件简单的事情，你若是遇到熟悉的同事、同学，大可拍拍他的肩膀：“老张，你好啊！”若是遇到领导或长辈，还以此种方式打招呼，对方心里就不高兴。若是遇到你的下级或晚辈，你只需等着对方和你打招呼就行了。

中国社会特别重视关系，彼此的关系良好，就算偶尔说错话，也没有什么关系。若是关系不够，或者关系不好，那就会鸡蛋里挑骨头。不过，也有本来关系很好，一句话伤了和气，便老死不相往来的可能。这就更加促使我们重视沟通的重要性。

人际关系与沟通，可以简称为人际沟通。在我们日常生活当中，人际沟通是不可或缺的活动，必须勤加练习，多加磨炼，养成小心应对、用心体会、虚心检讨的良好习惯。一方面使自己的沟通能力不断提高，一方面促使自己的人际关系获得改善。与其讨好别人，不如用心保持和谐、互动、互助的良好状态，透过好好沟通来互相感应。若能心意相通，大家都愉快，那就是良好的人际关系。在愉快中把正当的事情办理妥当，则是我们共同努力的目标。

人际沟通注重和每一个人进行良性的互动。既不能够偏重某些人，使其他人受到冷落；也不应该只顾自己，想说什么就说什么，爱说什么便说什么。否则你只是在发表意见，根本不是在进行沟通。善于沟通的人，必须随时顾及别人的感受，以免无意中破坏了自己的人际关系，造成恶劣的沟通效果。


\chapter{沟通的现象}
中西方的沟通态度大不相同，西方人进行沟通的时候，彼此不计较身份地位，畅所欲言。中国人则不然，就算面对面坐在一起，也会拘泥于彼此的身份地位，不敢贸然开口。

人虽有言论的自由，却不意味着可以乱说，或者想说就说，而是要谨慎小心，防止一语不慎招来祸端。

我们虽然处于重视沟通的时代，但是愈沟愈难通。理论上，沟通可以解决大部分的问题，实际上，沟通往往带来更多头疼的问题。为什么？因为现在越来越多的中国人学习西方人的沟通方式，结果造成不沟通还好、一沟通就糟的局面。

“先说先死”和“不说也死”这两句意思相反的话要合在一起考虑，深切体会这两句话的含义，就能掌握沟通的奥秘了。既然先说先死、不说也是死，那就要求我们，明白“先说先死”才会“不说”，了解“不说也死”才会“说”。因而“站在不说的立场来说”才不至于乱说，却能够说得恰到好处，做到“说到不死”。

要“说到不死”，必须说得合理。很多人就是认为自己所要说的十分合理，这才理直气壮，结果惹出麻烦。有一个故事，说一家人生了一个孩子，亲友都来祝贺，有的说：“这个孩子将来一定能当大官。”有的说：“这孩子将来会发财的。”……都是好听的。唯独有一个人说了一句“这孩子会死的”，就被主人打了一顿，赶出了家门。可是细想一下，那些说孩子会升官发财的话都不可靠，唯独“孩子会死”是句真话，可是说真话的人却被痛打了一顿，就是因为他这么说不合时宜。


\section{先说先死}

先说为什么会先死呢？先说的人说出一番道理来，后说的人很容易站在相反的立场，说出另一番道理。虽然双方都说得头头是道，毕竟后说的人可以针对先说的人，做一番整理和修补，甚至大挖其漏洞，弄得先说的人好像相当没有学问似的。

先说的人站在亮处，人家把他的底细摸得很清楚。后说的人若是存心挑毛病，专门挑他的缺失，保证把他整得体无完肤。先说的人，说来说去顶多能说出道理的一部分或者大部分，总有一部分被遗漏掉；后说的人，就可以针对这些缺失来大做文章，表现得很内行的样子。

有时候，人的身份地位不同，先说先死的情形也不同。比如，下属先说，说错了就会受到上司的批评，从下属的角度说，上司批评下属很正常。但是万一上司说错了，下属指出其毛病，那上司就会很尴尬：发火的话，就显得自己没度量；如果不发火，面子实在不好看。

有一次，化工厂厂长带领一群客人参观工厂，经过仪表控制室，忽然看见仪表板上，有若干颜色不同的指示灯，有亮着的，也有不亮的。有一个指示灯，则是一闪一闪的。

有人问：“这个指示灯为什么会闪?”

厂长回答：“因为液体快到临界点了，如果到达临界点，它就不闪了。”听起来也蛮有道理。

想不到厂长刚刚说完，仪表工程师说：“不是的，那个灯坏了。”

结果厂长表情极为尴尬。

仪表工程师先说，厂长才可以责骂他，叫他“先死”。如今自己先开口，不幸又说错了，若是此时指责他，显得自己恼羞成怒，似乎不很得体。

先说先死，好像是必然的，不过有的当场死，有的以后死，所谓死有先后，时间不同而已。就因为先说先死，所以明白此道理的人，与别人一见面从不说正经话，专说一些没有用的闲话。中国人不是不喜欢说话，而是中国话多半不容易表达得很清楚，话本身已经相当暧昧，听的人又相当敏感，于是“言者无心，听者有意”，往往好话变坏话，无意成恶意，招来洗不清、挥不掉的烦恼，何苦来哉？所以中国人对闲聊很有兴趣，见面不谈正经话，专说一些没有用的，就怕先开口，露出自己的心意，让对方有机可乘，徒然增加自己的苦恼。这样做表面看起来是在浪费时间，其实，其目的是让对方先开口，使自己获得有利的形势。中国人明知“形势比人强”，时时不忘“造势”，而说一些废话，正可以在不知不觉中造成有利的形势，何乐而不为?更何况，言多必失，废话说多了，难免会说漏嘴，透露一些有用信息，这样就可以明白对方到底是怎样想的，然后采取相应的应付手段。

中国人擅长明哲保身，就是因“先说先死”的痛苦经历造成的。中国人说话一向含含糊糊，让对方不明白其真实意思，就算随便一句打招呼：“要到哪里去?”得到的多半是“随便走走”之类的回答。只有碰到熟悉的朋友，才会说“我要去……”。

中国人十分习惯于“不明言”，即“不说得清楚明白”，却喜欢“点到为止”，以免伤感情。不明言的态度，比较不容易先说先死。因为一部分是我们说的，一部分是别人自己猜的，大家都有面子。同时也不容易被别人抓住把柄。

但是，现代中国人受到西方文化的冲击，在领悟“先说先死”之前，便勇敢地“有话直说”，弄得自己灰头土脸，依然不知道毛病出在哪里，也从不检讨自己，反而怨天尤人。或者已经惹祸上身却不能自知，反而沾沾自喜，以为得计。

很多人“有意见也不一定说”，往往鼓励别人先说，然后见机行事。他若不同意别人的话，就大肆抨击，抓住别人说话的漏洞，添油加醋来进行陷害，使别人哑巴吃黄连，有苦说不出。他若同意的，也可能赞扬备至，把别人的话改头换面，当做自己的真知灼见。这种让别人站在明亮处，自己躲在黑暗处的作风，使得别人不敢开口讲话，造成很多沟通的障碍。

不知道“先说先死”的人，常常死得不明不白。只知道“先说先死”的人，会落得难以沟通的评价，对自己的前途非常不利。

老子说“不敢为天下先”，孔子也大骂“始作俑者，其无后乎”，难道真的一点道理也没有？假如大家都不怕死，争先恐后地力求先说，偏偏又说得不伦不类，似是而非，究竟是好还是坏?

先说当然也可能不死，因为中国人的道理一向是相对的，“先说先死”固然是事实，“有话直说”也可能得到许多好处。既然二者各有利弊，强调“先说先死”，就是站在“先说先死”的立场来有话直说，才是中国人真正的功夫。“先说先死”是“根本”，“有话直说”不过是“作用”。有些人本末倒置，强调“有话直说”，等到自己吃亏上当了，再来怨天尤人，又有何用？

有话直说而不致害死自己的人，才是真正有本事的人。他们以“先说先死”的基本精神作为直说而不死的最佳保障。


\section{不说也死}

“先说先死”固然是事实，“不说也死”亦是不可否认的真理。只是“不说也死”是说给懂得“先说先死”的人听的，唯有深谙“先说先死”的道理，才有资格讲求“不说也死”。仅仅保持“我有话要说”的心态，根本没有必要深究“不说也死”的意义。

长久以来，“先说先死”把中国人害得很惨，浪费了很多宝贵的时间，也养成不善沟通的坏习惯。所以明白“先说先死”的道理之后，还要反过来告诉自己：不说也死。

小丽是老板的秘书，一向勤勤恳恳、规规矩矩，从不出大错。星期四她得到通知，说星期五公司有个舞会，小丽很想参加。虽然按照公司的规定，星期五可以不穿正装，但是身为老板的秘书，小丽每天都要穿职业套装，她不敢穿得太随便。可是既然有舞会，总不能穿正装参加吧？因此，小丽破例换上连衣裙，把自己打扮得漂漂亮亮的。她在老板办公室进进出出，老板看着很不舒服，但没说什么。下午，老板通知她：“3点钟有个紧急会议，你准备一下，负责会议记录。唉，你怎么穿成这个样子，赶快换掉。”小丽这才说：“公司有舞会，何况今天是星期五，公司规定……”老板火了：“到底是舞会重要还是工作重要？”

小丽认为自己并没有违反公司的规定，回答得理直气壮。殊不知小丽如果回答得没有道理，老板还可以批评她。她回答得有道理，老板更是下不了台，于是恼羞成怒，逼迫小丽换掉连衣裙，否则“炒鱿鱼”。结果小丽强忍泪水，赶快打车回家，换衣服。

如果小丽一开始就向老板暗示今天是星期五，可以穿便装，也许老板就会不以为意了。

不要以为多说多错，不说不错。有话不说往往会使你陷入被动的局面，举几个例子：

第一，如果你的上司交给你一项很复杂的任务，你完成不了，又一直不敢开口，最后任务完不成，那所有的过错都是你的。如果你早说了，你的上司就会想其他的办法解决：要么找其他的人做，要么他自己做，实在不行就从外面聘请专家。在一开始，凡事都有转圜的余地，而你明明完不成任务，还一声不响、硬着头皮继续做，往往贻误了时机。

第二，你很少说话，别人就很难了解你，不知你整天想什么，所以有晋升的机会也轮不到你，因为你的上司根本不了解你，又怎么敢提升你？孔子欣赏木讷的人，却也主张言词必须通达。少说话很好，但是少说话绝对不是不说话。

第三，如果你本来是个有说有笑的人，结果哪天了解到“先说先死”的法则，突然变得沉默寡言，别人会觉得有些蹊跷。平日少说话，忽然话多起来，或者一向多话，现在居然不说了，都会令人起疑。这种比较明显的变化，多半被认为是心理不平衡的反应。

第四，如果你和老板一起去拜访客户，老板不小心说错了话，你却不提醒，老板很可能把过错都推到你的身上，指责你隔岸观火，居心不良。其实老板选你一起去拜访客户，必然是经过考虑的，认为你会对他有帮助。但是你却漠不关心，难怪老板会火冒三丈。

“先说先死”和“不说也死”相反相成，凡事在说与不说之间，看情势、论关系、套交情，衡量此时、此地、此事对此人应该说到什么地步，才算合理。

换句话说，不能够由于害怕“先说先死”而不说，也应该顾虑“不说也死”的不良后果，慎重思考怎样说才不致一开口就闯祸。

“不说也死”是告诫我们，不沟通就难以协调与其他人的关系。大家都不说，根本无法沟通。不能沟通，当然无法协调。由于很多人受到“先说先死”的影响，不敢沟通，所以特别提醒大家，不说也死，希望大家早日摆脱“先说先死”的阴影，走出沟通的良好道路。


\section{说到不死}

“先说先死”和“不说也死”，构成沟通的两难。我们既然明白说也不好、不说也不好的困境，便应该设法加以突破。换句话说，最好能够“说到不死”。“说到不死”是一门高深的学问，需要在合适的时候、合适的地点，对合适的人，以合适的方式说出合适的话。如何判断合不合适，就要看你的功夫了。

要想“说到不死”，必须选对合适的人，有的话你对甲说没问题，对乙说就要触霉头。我们说“事无不可对人言”，又说“逢人只说三分话”，就是因为说话的对象不同。对知心朋友，当然“事无不可对人言”；而对一般人，则“逢人只说三分话”。比如一般人问你：“听说你要买辆跑车？”你的反应可能是：“没这么回事，我哪有那么多钱啊？”而熟悉的朋友若问你相同的话，你再否认的话，你的朋友就会认为你信不过他，所以你可能说：“我最近炒股票赚了点钱，是打算换辆车，但还没选好，你帮我参谋参谋。”说与不说，需要你衡量轻重，对一般人选择保密到家的策略，以免“先说先死”，而对朋友则采取私下透露，以“不说也死”的方式，以求“说到不死”。

下面几个案例就说明如何才能“说到不死”：

在工作年会上，总经理正在讲话，大家都在聚精会神地听，行政主管发现总经理遗漏了一项重要的行政决定，他不慌不忙地在便条纸上写下“关于……的决定”等，然后偷偷地递给总经理，希望提醒他，把此决定在会上公布一下。

行政主管的做法就很明智，如果等总经理讲完话，行政主管急忙站起来，补充说明一番，相信总经理必定很生气，不但不感激他的补充，而且事后必定气冲冲地责备行政主管：“你以为我把那项决定忘在脑后了？我记得比谁都清楚，只不过我认为暂时不宜在会上宣布，没想到你自作聪明，招呼都不打一声，就宣布了。”而行政主管必定会因“先说”而“先死”。

如果总经理真的忘了，而行政主管不说，那行政总管就会落到“不说也死”的境地：总经理会认为他根本心不在焉，这么重要的事都不提醒一下，以后根本不能信任他。

在说与不说之间，行政主管选择了一种合适的方式，即不明言，该提醒的也提醒了，至于总经理说不说出来，由总经理决定。无论以后有什么结果，总经理都不会怪到他的头上。

业务经理陪老板到客户那里谈判，客户提出让利3\%，业务经理当场拿出计算器，熟练地计算一番，然后把结果显示给老板看，嘴上说：“不行，这样我们就无利可图了！”老板看看结果，心里明白，接着说：“虽然如此，但是看在老客户的分上，再想想办法吧。”

明明可以接受，业务经理嘴上却说不行，实则将决定权交给老板。老板若同意，等于给对方一个人情；老板若不同意，则有充分的理由拒绝。所以，业务经理真正做到了“说了不死”。如果他计算完，不和老板商量一下，马上说“接受”或“不接受”，等于没把老板放在眼里，势必“先说先死”；如果他计算完，一句也不说，就等着老板做决定，老板就比较为难，因为他的做法摆明了告诉对方可以接受，老板再拒绝，岂不是让对方嘲笑？

“说到不死”其实就是说到合理的意思。只要合理，大家都能够接受，当然可以不死。


\chapter{沟通的艺术}
要想“说到不死”，就要掌握沟通的艺术。

\section{使对方听进去}

对方如果听不进去，就算你有千言万语，他全当耳旁风。对方听得进去，是良好沟通的第一步。所以开口之前，必须谨慎，以免徒劳无功。当对方听不进去的时候，我们宁可暂时不说，也不要逼死自己。能拖即拖，并非完全没有道理。运用得合理，也是一种有效的沟通方式。

中国人往往情绪反应激烈，一语不合，就可能翻脸。在沟通的时候，我们不能确保每一句话都说得很妥当，但至少从第一句话开始就特别小心，以诚恳的语气来使对方放心，使对方了解我们不会采取敌对或者让对方没有面子的方式来进行沟通。这样，对方才会逐渐放松。

第一句话就引起对方的戒心，使他觉得自己可能会吃亏，或者可能会没有面子，他就会采取躲避的策略；躲不开的时候，也会且战且走。一旦对方想“溜”想“躲”，就不可能获得圆满的结果。

中国人说话很少开门见山，而是先寒暄一番，看看对方的反应如何。如果对方心情不错，才可以进一步沟通。如果没说两句话，对方就很不耐烦，甚至要端茶送客，那你就算有再重要的事也要忍一忍，因为此时多说无益，“话不投机半句多”便是此理。

有人可能认为中国人的寒暄是在浪费时间，有正事不说，非得在无关紧要的事上大费唇舌，是不分轻重的表现。其实，他们根本不懂寒暄的妙处。东拉西扯，说一些没有用的寒暄话，目的在于了解对方的情绪状态，并且产生稳定对方情绪的作用。不急着讲，先摸清楚情况再说，乃是上策。

寒暄可以试探对方适宜不适宜沟通。比如，人们一见面，通常会说些无关紧要的话：

\section{“你最近气色不错。”}

对方如果说：“我最近吃不好、睡不好，气色怎么会好？”那你就知道对方心情不佳，不管什么事都需要延后，贸然说出来而对方一口回绝的话，连个商量的余地都没有了。

如果对方回答：“还好，最近没什么烦心事。”说明他心情不错，有什么事都可以说了。

寒暄可以缓解人们的紧张甚至是排斥的情绪，如果对方摆明不想听你说话，你通过寒暄可以渐渐使对方放松对你的戒备。春秋时期，触龙说服赵太后时就巧妙利用寒暄打动了严阵以待的赵太后。

赵太后刚执政的时候，秦国加紧攻赵。赵国向齐国求救，齐国的条件是让长安君作为人质。赵太后不肯答应，大臣们极力劝说，结果赵太后明言：“有哪个再来说要长安君为人质的，我就要把唾沫吐在他的脸上。”在这种情况下，触龙出马了。他并没有一上来就痛陈利害，而是一开始说自己脚上有毛病，竟不能快步走，但是因为好久都没见太后，怕太后身体不舒服，所以来看望一下。接着又问太后饮食起居的情况，然后把话题引到自己如何疼爱儿子上，引起赵太后的共鸣。最后说起赵太后与其女儿燕后的事，点明父母疼爱自己的孩子，就必须为他的长远利益考虑。触龙自始至终都没有提到让长安君做人质的事，但是他以聊家常的方式，让赵太后明白“父母之爱子，则为之计深远”的道理。

人们为什么不直接问“你此刻的心情如何?可以和你沟通吗”？那是因为，中国人十分机警，遇到这样的询问，大多不敢明说，以避免吃亏上当。在没有摸清楚对方的用意之前，谁也不愿意冒险回答这样的问题。这样问，对方大多沉默不语，问了等于没问。还不如稍微拐个弯儿，反而容易获得比较真实的答案。

当然，也可以不用寒暄就切入正题，但是要满足一定的条件：

其一，双方比较熟识，且要谈的事多半比较重要，前因后果双方也比较清楚，这时可以不用铺垫，直接切入正题。

其二，自己有把握吸引对方的注意力，让对方不得不按照自己的思路走。

《三国演义》中有几段关于说客的描写：

军士引阚泽至，只见帐上灯烛辉煌，曹操凭几危坐，问曰：“汝既是东吴参谋，来此何干？”泽曰：“人言曹丞相求贤若渴，今观此问，甚不相合。黄公覆，汝又错寻思了也！”

阚泽诈降曹操，必须要让曹操相信，如果直接表述忠心，以曹操的疑心病重，绝不会相信。所以阚泽反其道而行之，直言曹操徒有虚名，反而使曹操想听听他到底想说些什么。

马超端坐帐中不动，叱李恢曰：“汝来为何？”恢曰：“特来作说客。”超曰：“吾匣中宝剑新磨。汝试言之，其言不通，便请试剑！”恢笑曰：“将军之祸不远矣！但恐新磨之剑，不能试吾之头，将欲自试也！”超曰：“吾有何祸？”……

李恢去说马超，第一句话就说明来意，然后以一句“将军之祸不远矣”引起马超的关注，痛陈利害，顺利说服马超归顺刘备。

这些说客都是能言善辩之士，自然有能力牵着对方的鼻子走。但是，在现实生活中，这样的人确实不多见，所以中国人常常以迂回的方式，先让对方表明不会生气，然后说出可能会冒犯对方的话。

比如试探地问：“我有一句话不知当讲不当讲？”不要以为这是废话，这话是告诉对方：“我要说几句话，而这几句话可能会引起你的不快，你要有个心理准备。”

在《红楼梦》中，袭人对王夫人说贾宝玉该打这种绝对会冒犯王夫人的话时，也采取一波三折的方式：

袭人道:“别的缘故实在不知道了。我今儿在太太跟前大胆说句不知好歹的话。论理……”说了半截忙又咽住。王夫人道：“你只管说。”袭人笑道：“太太别生气，我就说了。”王夫人道：“我有什么生气的，你只管说来。”袭人道：“论理，我们二爷也须得老爷教训两顿。若老爷再不管，将来不知做出什么事来呢。”

为了让对方听得进去，我们很容易采取讨好的方式，尽量说一些好听的话，让对方听起来很高兴，而易于接受。其实，想讨好中国人，并不简单。

既不能单纯地讨好对方，又要让对方听得进去，唯一的办法就是把话说妥当。只是说得对并没有用，有时候，你说得越对，对方会觉得越没面子，以致恼羞成怒，更加听不进去。说错了，后果更加严重。

说得对，还不如说得妥当来得有效。但每一句话都要说得很妥当，实在不容易。任何话一出口，对方大多不会“就听到的话来判断”，反而多半“在听到的话之外去猜测用意”。弦外之音，往往比说出来的话更重要。

\section{要学会察言观色}

我们在开口说话之前，要谨慎小心，努力让对方听得进去。另一方面，我们也要用心听取对方所说的道理，不要过于在意对方怎么说。

任何一句话，认真去听，都可能听出某些道理，不可能毫无价值。但是，我们常常不在乎这些道理，却斤斤计较于对方表达时的态度和语气。换句话说，我们不认真听对方在讲什么，却十分介意对方是怎么讲的。事实上，愈有道理，愈容易引起听者的反感，所谓忠言逆耳。只要双方都认真地听，那么中国人沟通起来就会顺畅得多。

但中国人很奇怪，心中有话不一定说出来，而要等着对方来猜；就算我们勉强说出来，也必定说得含含糊糊、不清不楚；而当我们说得很肯定的时候，对方就更小心了，因为说得斩钉截铁的未必是真话。

中国人这样做，不是没有道理的，只是不了解的人，很难明白我们的真实意图，造成沟通障碍。既听他的话，又看他说话的样子，综合判断，才可以决定信或不信。

这时就需要我们发挥察言观色的本领，关注对方说话时脸部的表情。表情比言语本身更能表达内心的动态。人类五官之中，眼睛是最敏锐也最诚实的。《孟子\&\#8226;离娄篇》说：“存乎人者，莫良于眸子。眸子不能掩其恶。胸中正，则眸子了焉；胸中不正，则眸子眊焉。听其言也，观其眸子：人焉廋哉！”意思是，观察人的邪正，没有比观察他的眼睛更准确的了。眼睛不能遮掩人的恶念。心正，眼睛就明亮，心不正，眼睛就昏昧。听了他的话，再看他的眼睛，人的邪正，哪里隐藏得过去呢？

说话的速度、说话的音调、说话的节奏等，能帮助我们揣摩对方的心理。比如，说话的速度常常能反映一个人的心情，说话快的人突然慢下来，那他可能有些不满，而说话慢的人忽然加快语速，他可能在说谎，或者心中怀有愧疚。

再比如说话的音调。一般人说谎时，由于害怕事情被揭穿，音调会不自主地提高。同时，为了反对他人的意见，也可能提高自己的音调。

说话的节奏也很重要。节奏比较顺畅时，说明他很有信心；若张口结舌、吞吞吐吐，说明他缺乏自信。

喜欢复述说话者的言辞，表示自己一直在注意听；一边听话一边点头，表示全神贯注，心无旁骛；自问自答的人，多半相当顽固；既不肯定又不否定的人，往往具有神经质。

要准确做出这一类的判断，最好提醒自己：每一个人的观念，都不太一样，必须平日多做沟通，促进了解，把对方的价值观和人生观摸清楚，然后再来评断，通常比较准确。否则把坏人当成好人，将好人看成坏人，不但自己吃亏，也会引起他人的不满。

特别是对一些老于世故的人，喜怒不形于色，很难从其表情上看出其内心活动。所以若非经过多次观察，最好不要轻率地加以判断。

古语说：“衣不如新，人不如故。”大抵因为相处时间长了，比较容易了解对方的心理。新人相处不久，彼此互不了解，看他怎么说，事实上和听他怎么说同样困难。

\section{站在对方的立场}

在沟通方面，中国人最重视圆满，也就是设法站在每个人的立场上，让大家都有面子。如果是很多人在一起的时候，不能只照顾几个人而冷落其他人。被冷落的人觉得很没面子，就会引起情绪上的反弹，故意制造很多问题，不但增加沟通的困难，还会产生难以预料的不良后果。《红楼梦》中的王熙凤在初见林黛玉时，说她“况且这通身的气派，竟不像老祖宗的外孙女儿，竟是个嫡亲的孙女”，林黛玉远来是客，夸奖她是应该的，但是当时迎春姐妹都在场，如何只夸奖黛玉的话，恐怕她们会觉得不快，所以王熙凤一句“竟是个嫡亲的孙女”，在夸奖黛玉的同时，又肯定了迎春姐妹，使大家都很有面子。

当然，在人数众多的情况下，让每个人都有面子，确实很难。但在沟通的过程中，尽量站在对方的立场上，则有助于沟通的顺利进行。《三国演义》中，李恢劝降马超，始终是替马超考虑，先指出“将军之祸不远矣”，接着说明马超的困境——“目下四海难容，一身无主”，最后说“刘皇叔礼贤下士，吾知其必成，故舍刘璋而归之。公之尊人，昔年曾与皇叔约共讨贼，公何不背暗投明，以图上报父仇，下立功名乎”。一番话，既帮马超解决眼下的难题，又替马超筹划了未来，马超又怎会不投降刘备呢？


\chapter{沟通的层次}
沟通一般可以划分为四个层次：第一是不沟不通；第二是沟而不通，不管你怎么沟就是不通；第三是沟而能通，比较顺利；第四是不沟就通，这是比较高的层次。


\section{不沟不通}

从本质上讲，不沟不通算不上沟通，甚至可以说是沟通的反面。但是，本书将不沟不通作为沟通的起点。任何沟通进行前，都是不沟不通的状态。不沟不通，是指人们没有沟通的欲望或沟通的必要，处于不相往来的状态。比如，两人虽然彼此认识，但是工作、生活基本没有交集，不需要“通”，所以也没有“沟”的必要。


\section{沟而不通}

中西方人在沟通方面存在很大差异，西方人有什么事都会讲出来，其看法是“我有事情，想跟你讲”；中国人有什么事决不会讲出来，其看法是“我有事情，根本不跟你讲”——我有事，为什么要跟你讲？由此可见，西方人的出发点是“我要说”，中国人的出发点则是“我不跟你说”。

就因为出发点不一样，结果也不一样。很多中国人学习了西方的沟通理论后，想在中国社会广泛应用，却往往造成沟而不通的现象。以西方人的沟通原则为基础，无论你开口说话，或者闭嘴不说，都可能沟而不通。

有时候，人们虽然在滔滔不绝地说话，但是一方说的话，另一方根本没听进去。甚至，听的人很生气，当面加以指责，使说的人简直下不了台，或者表面上装成无所谓的样子，内心却恨死说话的人。这就是沟而不通的现象。其实，在现实生活中，很多人都停留在沟而不通的层次上。说了很多话，却收不到效果，无法达成预期的沟通目标。

古人常说：“武死战，文死谏。”“文死谏”无非是沟而不通导致的悲惨结局。中国历朝都有专门的官员负责向皇帝进谏，但常常因沟而不通，触怒皇帝，最后丢了性命。在现代社会中，沟而不通虽然不至于让我们没了性命，但会产生巨大的障碍，使我们寸步难行。

中国人相互之间为什么不容易沟通？首先，和面子有关系。沟而不通，和所说内容的对错没有必然的关系，并不是说得对就“通”，说得不对就“不通”。有时候，说得对也很危险，对方常常会因为你说得对而觉得没有面子：不是觉得你在指责他不对，就是认为你把他当做傻瓜，嘲笑他连浅显的道理都不懂。人没面子就会恼羞成怒，以致蛮不讲理，当然沟而不通。修养好一些的人，表面上很高兴，其实心里颇不以为然，同样是沟而不通。

在这种情况下，闭嘴不说也不行。对方如果看到你不开口，也来个相应不理，造成尴尬的气氛，想说的人更不知从何说起，也会沟而不通。

因此，说与不说、说对说错都可能沟而不通。中国人之所以先说先死，是因为任何人所说的，不过是片面的道理，充其量只能说是自圆其说。别人要支持你，固然可以说出一大堆道理，使我们觉得很有面子；若是不支持，照样可以陈述许多道理，让我们颜面无光，倍感羞惭。实际上，在中国社会，某一个人赞成另一个人的观点，并不意味着是对其观点的完全认同和肯定，而是因为前者能感觉到后者对他的友善而予以的支持，这是一种道义的表现。所以说话的人，为求立于不败之地，必须先摸清楚能不能获得别人的支持。

其实，情绪化也是中国人不易沟通的主要原因。虽然全世界人都有情绪，但只有中国人情绪起伏最大，动不动就发脾气，你还跟他沟通什么？不挨骂就不错了。

中国人的情绪，说起来十分有趣。由于警觉性很高，所以怀疑心很重，因此情绪的起伏很大。往往一句话听得不顺，就会想得很多。而且愈想愈多，也愈想愈火。以“你把我惹火了，我当然不讲理”为借口，干脆蛮不讲理到底，看你还有什么本事，可以和我沟通?

在这种情况下，唯有小心翼翼，步步为营，一句话都不能讲错，不把对方惹火，才能顺利沟通。

如果预料到得不到沟通的效果，那还不如不说，最多引起对方的不愉快，大家的损失并不大。如果对方是有识之士，知道“不说话并不代表无话可说，而是不知道该不该说，要怎么说才有效”，于是制造沟通的渠道，来增强沟通的信心，添加沟通的气氛，反而容易造成沟而能通的美景。

不说毕竟是消极的做法，要想改变沟而不通，最好是站在对方的立场来沟通，凡事以对方的利益为出发点，就算你说的话他不接受，至少他会思量一下。久而久之，就能获得对方的信任，这样沟通起来就方便得多。如果得不到对方的信任，你说的每一句话他都要考虑一下是何用意，沟通起来就相当困难。

当事情有多种选择时，还可以瞄准对方的需求，模拟出若干可能的方案，再依对方的立场和行事风格，评估、分析、选择其中最合适的。刘备欲取西川，庞统献了三条计策，让刘备选择。最后刘备认为“上计太促，下计太缓；中计不迟不疾，可以行之”。当然了，如果几种方案均不可行，还可以再行调整，以求获得合理回应。


\section{沟而能通}

沟而能通当然是人们喜闻乐见的情况。误会也好，分歧也好，只要沟而能通，都不是问题。

沟通也要讲缘分，有缘的人在一起，肯定会谈得很投机，好像没有什么禁忌，什么都可以谈，怎样说都可以。无缘的人在一起，多数是话不投机，或者是各说各的，说的时候很热闹，但不会产生任何效果。

其实，只要关系够、交情深，而且场合适宜，即掌握了天时、地利、人和，中国人照样能够有话直说、有话实说，不但彼此畅快，而且效果良好。

前面说过，中国人之所以不易沟通，首先是因为面子问题。也就是说，当对方觉得很有面子的时候，大多比较容易沟通。

其次，情绪问题也是影响沟通的因素，但是当中国人情绪好的时候，不管说话的人如何唐突、冒犯、无礼，都能够心平气和地合理响应，一副大人不记小人过的模样。如果双方都能够有诚意、能包容、不计较，当然沟而能通，一点障碍都没有。

要想让别人感觉有面子，并且有好的情绪听你讲话，关键是你怎么说。这样就存在一个问题：如何处理妥当话与真话的关系。有时候，真话反而是最刺耳的。不要以为反正是实话实说，怎么说都一样。其实，同样一句话，说得妥当一些，对方比较容易接受；说得十分真实，对方往往受不了，反而听不进去。

比如说，“你这件衣服太难看了，我送你件新的吧”，这么说的话，虽然送别人东西是好意，但绝对不妥当，对方很难接受。我们常说有人好心办坏事，就是因为办得不够妥当，才事与愿违。


\section{不沟而通}

不沟而通是一种高度的艺术。我们先求说到不死，再求沟而能通，逐渐走入不沟而通的境界，当然十分美妙!

中国人十分讲究人与人之间的默契，高度的默契便是不沟而通，是一种难得的沟通美景。有时候人们不需要说话，光靠眼神、动作就能传达意思。赤壁之战时，“一阵风过，刮起旗角于周瑜脸上拂过。瑜猛然想起一事在心，大叫一声，往后便倒，口吐鲜血”，只有诸葛亮知道周瑜的心病，开出“欲破曹公，宜用火攻；万事俱备，只欠东风”的药方，这也是不沟而通的例子。

当然，我们不可能都像诸葛亮那样聪明，但是要想做到不沟而通，却也不是毫无章法可循。

不沟而通的关键在于双方的默契，而要建立默契，就要关注对方，随时随地注意对方的举动，不依赖对方的言语表达，却主动地捕捉对方的肢体语言。毫不关心对方，不注意观察对方的举动，当然无法不沟而通。我们只有将心比心，通过心与心的感应，使对方的心意能够畅通地传过来。心意相通，自然不沟而通。


\chapter{沟通的原则}
中国人普遍重视诚信，对于不诚无信的人，非常厌恶。但是，沟通的时候，却常常要求对方：

“我告诉你，你不要告诉别人。”

稍为放宽一些：“你如果要告诉别人，就不要说是我说的。”

“如果你告诉别人是我说的，我一定说‘我没说’。”

中国人说这些话时，根本没有欺骗的感觉。可见它和诚信并没有关系，也就是不属于不诚无信的范围。听起来怪怪的，却具有相当的道理。

第一句话表示你我关系不同，所以我才告诉你，我说的话，你相不相信、相信到什么程度、要不要转述、转述到什么地步，都必须由你自己决定，不要赖在我的身上。我告诉你不要告诉别人，事实上并没有什么约束力，只是好意提醒你，一切由你自己做主。

第二句话的含义是，你如果决定要告诉别人，表示你已经充分明了、相信我所说的话，是经过考虑后告诉别人的，这时候你所说的话，已经是你自己研判之后的资讯；而你所要告诉的对象，也是你自己审慎选择、决定的。一切都与我无关，所以不必再把我拉扯进去，说什么是我说的。

也就是说，我既然可以告诉你，你在必要的时候当然可以告诉别人。但是你必须审慎地选择对象，就好像我选择你一样。只要关系够亲密，可以放心地告诉他，不必担心产生什么不良的后果。但是，你告诉别人的时候，千万不要提及我，以免使我难做。要不要告诉别人，我充分尊重你的决定。但在保密信息来源方面，请你充分尊重我。

我告诉你的话，是经过我的分析辨别之后才告诉你的，所以你告诉别人的时候，怎么说，说多少，都由你自己决定。你所说的内容，当然也是你自己认为可以说的，因此已经和我说的没有什么关系了。如果此时你再说这些内容是我说的，那就太不尊重我了。

而听话的人却希望说话的人必须负起责任，“都是你说的，不然我不会做这种事”。意思是，我并不是故意犯错，或者能力太差，而是因为你这么说，我又很相信，才弄得如此不堪，罪不在我。发生这种事，是怪说话的人不负责任，随便乱说，还是怪听话的人不懂分辨，太轻信别人？其实，双方都有责任，应各打五十大板。不过，为了杜绝此类事情发生，我们就要严格约束自己。不管我们是说话的人，还是听话的人，都应该对自己的行为负责，既不要以讹传讹，也不要偏听偏信。

从听话的人的角度来说，说出是谁说的，以给别人有求证的可能，不但伤了与说话人的和气，也会使自己处于被动的局面。因为说话的人若言之凿凿地否认自己说过这种话，会使别人对听话的人的信用产生怀疑，认为他是在搬弄是非。

如果说话的人纯心以假话骗人，等听话的人相信了又转告他人的时候，再拆穿假话，就会使听话的人百口莫辩。因此，为了安全，听话的人若真要转告他人时，最好以模棱两可的语言转述，以免有人较真儿，产生不良后果。

第三句话的含义是，如果你一定要告诉某人，却又指名是我说的，鉴于这个对象根本不是我选定的，要说哪些话，说到什么程度，也不是我所能控制的，因此我只有表态：我并没有说这些话，至少我不是这样说的，语气、用词都不相同。

明明是自己说的，却常常加以否认。因为我告诉你的时候，已经提醒你不要告诉别人。你一定要告诉别人，实际上我也禁止不了。所以我只要求你，不要说是我告诉你的。现在你不但告诉别人，而且公开信息的来源，明白指出是我说的，那我就加以否认。这样有什么后果你自己负责，怨不得我。因为你不理会我的要求，指明是我说的，不但不尊重我，还把责任推给我，我只好郑重否认，把责任再推回给你。

再者，你告诉别人是我说的，我又弄不清楚你同谁说的，当时是怎么说的，现在居然要我承担责任，我怎么能够承认?如果我的否认使你没面子的话，我也没办法。你既然告诉别人，就应该自己承担所有的后果。


\chapter{沟通的方向}
沟通，按信息流动的方向来分，可以分为向上沟通、向下沟通和平行沟通。我们说，中国人的人际关系实际上是一种人伦关系，而人伦是有大有小、有上有下的，人们不能以大欺小，也不能以下犯上。

向上沟通是指居下者向居上者陈述实情、表达意见，即我们通常所说的下情上达，如臣对君、子对父、下属对上司等。

向下沟通与向上沟通正好相反，是居上者向居下者传达意见、发号施令等，即我们通常所说的上情下达。

平行沟通是指同阶层人员的横向联系，如公司内部同级部门之间都需要平行沟通，以促进彼此的了解、加强合作，免得产生隔阂，影响团结。平行沟通的目的是交换意见，以求心意相通。

这三种方向的沟通，对任何人而言，都是常用的。而且流动的方向不是一成不变的，随着具体情况的不同而随时改变。


\section{向上沟通}

在向上沟通中，“下”应是主体。以下属与上司进行沟通为例，下属经常要汇报、请示、建议，甚至犯错误以后要辩解，这都属于向上沟通。在向上沟通时，下属应该谨记“上下”观念，不可以下犯上，当然也不必奴颜婢膝。

下属应该安守本分，小心自己的言行，让上司感受到尊重，否则上司只好摆出上司的架子，下属反而觉得没有面子。有些下属，上司稍微对其客气些，便得意忘形，没大没小起来，弄得上司不得不摆出高姿态。

最好的办法是，表现得不卑不亢：太“亢”，会使上司的面子受损，上司真要还以颜色，吃亏的还是下属；太“卑”，上司会觉得这个下属太无能，靠赔笑脸混日子。下属不能让上司感到功高震主的威胁，也不能让上司觉得自己无能。

有时候，下属与上司的意见发生分歧时，下属应该合理地坚持，才能树立信用。否则，上司一说“不”，下属就见风转舵，只会给上司留下不可靠的印象。即使坚持自己的想法，也要先肯定上司的意见，再表达自己的意见；或者提出问题，反过来请教上司。当然，也不能过分坚持，否则就成了顶撞，上司以为下属在挑战他的权威，难免会以势压人，容易伤了上下级的和气。当出现严重分歧时，下属最好不说话，显出一直在深思的样子，上司自然看得出来，但是下属不要说出来，等上司让说时才说。意见相同时，要给予热烈的回应，不要说“我也这么想”，上司会以为你是随声附和，最好说：“我想了许久都没有想通，原来这样最好。”

打开上下级言路的最好方法是让上司感到下属心中有他。一些企业老板毫不隐晦地说，他们最重视下属心中有没有老板。如果你能让上司感到你心中有他，你说的话他就会听得进去，因为他觉得你的所作所为都是为他着想的。相反，如果你做不到这一点，你提出任何意见，他都会怀疑你的动机，也就不容易接受你的意见。


\section{向下沟通}

向下沟通时，“上”应是主体。要想沟通顺畅，上司要降低自己的姿态，不要一副高高在上的样子，使下属畏惧，产生不愿意沟通的反感。中国人重视身份地位，所谓的“大人不计小人过”，就是“大人”不愿意放下身份去同“小人”斤斤计较。所以，越是位高权重的人，越会表现出平易近人的样子，同下属说话的时候就如同仁慈长者，多数是谆谆教导的口吻。凡是那些动不动就大发雷霆、咄咄逼人的上司，一般是火候没到，还不懂得中国人的“为官之道”。

三国时，吴国大将吕蒙骁勇善战。他自幼家贫，没有读过书，吴王孙权常劝吕蒙多读点书，但是吕蒙以军务繁忙为借口推脱。在这种情况下，孙权并没有指责他不学无术且不听劝告，而是说：“你说事多，难道还多过我吗？我常常抽时间读书，感觉受益颇多……”经过孙权的劝导，吕蒙终于发愤读书，这才有了“士别三日，刮目相看”的成语。

某公司副总第一次到香港出差，回来后满腔热情，在开会的时候不谈正经事，大谈特谈香港的风土人情。与会的人听得不耐烦，又不好打断他，只好都看着总经理。总经理该怎么办？他已经暗示过副总了，可是副总视而不见。一个总经理连会场都控制不住，还当什么总经理？于是，总经理只好用强制的方法，对副总说：“这事以后再谈，我们现在谈正事。”副总当面被驳，面子上有点挂不住，为了不在下属面前颜面扫地，他依然故我，宁可散会以后挨骂，也要在下属面前保住面子。

这下子，轮到总经理面子上挂不住了，其实总经理可以这样说：“李科长，你安排两个小时的时间，我们专门请副总来说说他到香港的事情，今天我们抓紧时间解决问题。副总，你的事情十分钟也讲不完，我给你安排两个小时，做一个专题演讲，好不好？”这样，既给副总留了面子，又限制了他说题外话。当领导的，被别人牵着鼻子走，下面的人谁会服你？可是一旦领导太?了，所有人都会转而同情弱者，让领导进退两难。


\section{平行沟通}

对上沟通、对下沟通，彼此之间都会保留三分的礼让空间，比较容易找到合理的平衡点。平级之间，大家一样大，很容易产生“谁怕谁”的心态，对沟通十分不利。中国人其实谁都不服谁，当大家处在同等地位的时候，这种现象更明显，甚至一句话、一个动作，都会引起他人的不满。比如，你说“听我说”，别人会想：“别自以为是了，凭什么听你的？”

在这种情况下，要想进行顺利的沟通，要先从自己做起，尊重对方，对方才会用同样的态度对待你。中国人是讲交互的，你敬我一尺，我敬你一丈，所以，你尊重对方，对方也自然会尊重你，这样才方便沟通。

沟通的时候要以真诚的心去交流。中国人疑心比较重，稍不留神，就会触动其怀疑的神经，一旦中国人产生疑心，你说什么都没有用，因为他什么都不会相信。中国人说“一次不忠，百次不用”，不要以为中国人太决绝，这只是小心谨慎的表现。千万不要存心占便宜，也不要设陷阱存心让别人出丑，中国人很怕吃亏，吃过一次亏后，就不愿再和你打交道。所以，建立“和我打交道，一定不吃亏”的信用，增强大家对自己的信任感，才能方便沟通。

\end{document}